\documentclass[11pt]{article}
\usepackage[T1]{fontenc}
\usepackage{lmodern}
\usepackage{amsmath,amssymb,amsthm,mathtools}
\usepackage[margin=1in]{geometry}
\usepackage{booktabs,array,longtable,enumitem}
\usepackage{microtype}
\usepackage[colorlinks=true,linkcolor=blue,citecolor=blue,urlcolor=blue]{hyperref}
\hypersetup{pdftitle={Theta manifold geometry and signed heat-flow collision control},pdfauthor={Drafted for Edward Baker},pdfsubject={Scalar stationary geometry, complete prime-dilation heat identities and conditional theta collision control}}
\allowdisplaybreaks
\setlength{\emergencystretch}{2em}
\numberwithin{equation}{section}
\newtheorem{theorem}{Theorem}[section]
\newtheorem{proposition}[theorem]{Proposition}
\newtheorem{lemma}[theorem]{Lemma}
\newtheorem{corollary}[theorem]{Corollary}
\newtheorem{classicalinput}[theorem]{Classical input}
\theoremstyle{definition}
\newtheorem{definition}[theorem]{Definition}
\theoremstyle{remark}
\newtheorem{remark}[theorem]{Remark}
\newcommand{\R}{\mathbb R}
\newcommand{\C}{\mathbb C}
\newcommand{\Lag}{\mathcal L}
\newcommand{\dd}{\,d}
\newcommand{\Ephys}{E}
\newcommand{\researchnote}[2]{\href{../../papers/investigations-2/active/22_theta_manifold_geometry/#1}{#2}}
\title{Theta manifold geometry and\\signed heat-flow collision control}
\author{Drafted for Edward Baker}
\date{Research manuscript --- 15 October 2026}
\begin{document}
\maketitle
\begin{abstract}
We study the collision detector $E=H^2+H_x^2$ for the complete theta
heat flow through signed jet geometry and arithmetic source identities.
Every finite collision forces a negative fixed-endpoint chord of order
$\sqrt E$. Exponential averaging and endpoint-only approximation yield
a conditional exclusion criterion with independently chosen spatial
lengths and cutoffs. Near each finite collision, the complete backward
stationary locus is one analytic curve, its physical first jet has a scalar factor of the full
collision order, and prescribed positive cones reduce the sufficient
trace estimate to a scalar one-sided bound. Joint source and coefficient
currents have one gauge-invariant total boundary trace; signed quotient
constants may remain bounded when individual coefficients diverge.
The normalized coefficient recurrence identifies theta and gives exact
physical defect formulas and sharp finite-hierarchy stability, but its
inversion leaves the unknown observation at a lower boundary.
Multiplicative source dilations implement classical prime excision and
M\"obius inversion. We prove uniform endpoint remainder bounds and retain
the complete transformed lower intervals and primitive edges.
At positive heat time, the separated complete-transform and lower-interval
M\"obius series fail even the term test, while their paired differences
converge rapidly. Whole-line inversion loses heat integrability, and
nonlocal recurrence, Euler-summation and Poisson candidates expose
further specific boundary or collective-cancellation obstructions.
Signed forms, holomorphic jet surfaces and theta connections provide
the geometric setting; positive-source collision controls distinguish
arithmetic rigidity from a favorable physical sign. The remaining
problem is an independent theta-specific signed trace estimate with
full domain coverage. No new global theta sign bound or proof of RH
is claimed.
\end{abstract}
\setcounter{tocdepth}{2}
\begingroup
\small
\tableofcontents
\endgroup
\clearpage

\section{The geometric problem and the program}\label{intro:section}

The theta heat flow presents a tension between positive source data and
a signed spectral observation.  The complete source is positive and has
very rapid decay, but its oscillatory transform can have small first
jets.  The physical collision detector is
\[
 E(t,x)=H(t,x)^2+H_x(t,x)^2.
\]
The research objective is to exclude $E=0$ at finite positive heat time
on the full real spectral line.  The established de Bruijn--Newman
framework supplies the relation to RH, stated with its domain
requirements below.  This paper develops the differential geometry of
that objective; it does not prove the missing theta-specific estimate.

There are two natural manifolds.  On the physical side, a primitive
$P_z=H$ gives a holomorphic surface
$X_t(z,w)=(P(t,z),H(t,z),w)$ whose induced metric contains the complex
first-jet energy as a coefficient.  Its regular locus is a K\"ahler
manifold; a common zero of $H$ and $H_z$ is a branch event for the
parametrization.  The real projective tangent direction retains the
signed Laguerre quantity through
\[
 \partial_x\arg(H+iH_x)=-\frac{H_x^2-HH_{xx}}{E}.
\]
The associated two-point observation is
\[
 \mathcal C_j(t;x,y)=h_j(t,x)\{h_j(t,y)-h_{j+1}(t,y)\},
 \qquad h_j=\partial_x^jH,\quad j=0,1.
\]
At a full stationary start $h_{j+1}(t,x)=0$, on an interval with
$h_j\ne0$, its nonnegativity is precisely a $\pi/4$ upper bound on
the rightward projective turning. More usefully, every collision forces
$\mathcal C_j/\sqrt{E(t,x)}$ to approach a strictly negative constant
at a fixed nearby endpoint. A lower bound with error
$o(\sqrt{E(t,x)})$ therefore suffices across all finite multiplicities.
The energy is evaluated at the starting point on both parity branches.

On the source side, the actual theta kernel is a covariant derivative
expression in a degree-one line bundle over a curved elliptic family.
The complete two-source transform retains its spectral phases, positive
heat multiplier, transverse jet corrections and primitive endpoints.
Its directly restricted connection is flat, so curvature alone adds no
new tangential sign. The arithmetic question is sharper: which coupling
of the theta coefficients can control this complete signed observation?
A positive weighted-square source can collide while retaining the first
source endpoint; full modular reflection excludes those free weights
by forcing the theta coefficients. Turning this rigidity into an
inequality remains open.

\subsection{Results and their logical status}
The principal results and the remaining candidate are as follows.
\begin{enumerate}[leftmargin=*,itemsep=3pt]
 \item The fixed-endpoint chord theorem gives conditional collision
 exclusion from an $o(\sqrt E)$ starting-energy error at every finite
 multiplicity. The exponent $1/2$ is sharp for arbitrary power-law
 allowances; the inner heat profiles resolve the turning concentration.
 \item Exponential averaging gives the endpoint observation
 $W^2-e^{-\delta}W(x)W(x+\delta)$. Endpoint-only approximation preserves
 the full starting factor. A mixed bound $-K_NE-\varepsilon_N\sqrt E$
 suffices with independently chosen cutoffs, even when the locally
 bounded constants $K_N$ grow arbitrarily with $N$. The earlier
 additive-tail transfer remains available under its stated rate bounds.
 \item Exact normalized recurrence defects control physical mode
 observations with Gaussian dual weights. Positive recurrence
 relaxations through index $N$ have sharp $O(1/N)$ stability, improved
 to a Gaussian beta tail by a slope anchor. These are stability results,
 not favorable-sign estimates at the exact theta coefficients.
 \item Finite neighboring-mode and joint source currents give
 $O(L^{-1/2})\sqrt E$ transfer errors under their stated local source
 bounds. Their combined boundary trace retains the unknown observation
 and is invariant under admissible changes of current. Polynomial source
 divergences do not remove the specified finite-range no-inverse result.
 \item Multiplicative source shifts give exact prime-excision and
 M\"obius identities, with explicit moment remainders and every lower
 interval retained after the heat transform. Their separated
 positive-time series diverge, although complete paired terms converge.
 A massive recurrence current, a floor-counting representation and
 finite-source Poisson duality identify additional scoped obstructions.
 These results use classical arithmetic identities and supply no sign
 for the remaining stationary cross observation.
 \item Complete signed source formulas and coupled addition retain
 independent source variables, parity, exchange terms and primitive
 endpoints. Positive-source controls show why weakened arithmetic
 hypotheses are insufficient; complete reflection fixes the input.
 \item Signed forms, holomorphic jet surfaces, calibrated inclusions
 and theta connections supply the geometric setting. The recurrence
 trace has a one-dimensional cokernel, and bounded endpoint descent to
 the physical-jet quotient requires control through degeneration. A
 signed trace estimate is weaker than absolute endpoint reconstruction;
 its optimal quotient constant permits some singular coefficients.
 \item Near a finite collision, the entire backward stationary locus is
 one analytic curve and the opposite parity stationary set is empty.
 The physical first jet has the principal factor $s^r$, without losing
 its radial order. Restricting starts to any prescribed positive cone
 retains collision coverage and gives a scalar signed-distance target
 with explicit mixed errors and reconstruction costs.
\end{enumerate}
The remaining task is an independently proved signed boundary estimate
using the complete theta coupling. The analytic collision criteria are
conditional; the control sources are not theta; finite checks validate
specified calculations. These roles must not be conflated with a new
theta sign domain or an RH conclusion. We make no novelty claim for the
standard geometric constructions.

\subsection{Relation to the preceding investigation}
This investigation separates the manifold and connection direction
from the broader signed-theta program in \cite{BakerTheta}.
The preceding manuscript develops complete-source correlations,
cutoff accounting, outer-region estimates and finite certificates.
Its self-dual construction and its later manifold notes are the starting
point of the present narrative.  The outer estimates and certificates
retain their original statements and verification records; they are not
reproved or used as a replacement for the missing global geometric
estimate here.

The author's exploratory quaternion-form notes of January and February
2014 \cite{BakerQuaternion} motivate preserving a signed, oriented
pairing.  We build the required identities directly and state their
boundary hypotheses.  The new line-bundle construction is abelian and
has an explicit connection and curvature; it invokes no unproved
general nonabelian boundary formula.

Throughout, a regular jet interval means an interval on which the
particular pair defining the angle is nowhere zero.  Source-prescribed
constants or moduli must be derived from the complete theta data and
remain uniform through neighborhoods of hypothetical collisions.  A
constant chosen retrospectively from an observed positive energy floor
does not provide such an estimate.

\subsection{Conventions}
All metrics are real Riemannian metrics.  We use
$|dz|^2=dx^2+dy^2$ and the K\"ahler form
$(i/2)\,dz\wedge d\overline z$.  The physical time $t$ is an external
parameter for the four-dimensional exterior derivatives.  It is
distinct from the elliptic modulus $\tau$, whose imaginary part is
positive.  Hodge signs refer to the orientation stated in each
construction.  For a Hermitian holomorphic line bundle in a holomorphic
frame with squared norm $h$, we use the Chern connection
$\nabla= d+\partial\log h$ and curvature
$F_\nabla=\overline\partial\partial\log h$; a frame change transforms
the local connection in the usual way.
% Prepared for Edward Baker, 15 October 2026.
% Substantial LLM assistance: coordinating gpt-6-astra, effort ultra,
% verified and supplied by the parent session. Independently review before use.

\section{The physical heat field and its collision obstruction}
\label{phy:physical-field}

The geometric constructions in this paper must retain a specified physical
observation.  We first fix that observation and prove the local analytic
facts used later.  The conclusions in this section about arbitrary analytic
heat fields are local theorems.  Their implication for the Riemann hypothesis
uses an explicitly identified earlier theorem for the complete theta source.

For real $u\ge0$, define
\begin{align}
 \Phi(u)&=\sum_{n\ge1}
  (2\pi^2n^4e^{9u}-3\pi n^2e^{5u})e^{-\pi n^2e^{4u}},
 \label{phy:source}\\
 H(t,z)&=\int_0^\infty e^{tu^2}\Phi(u)\cos(zu)\,du.
 \label{phy:heat-field}
\end{align}
The variable $t$ is an external heat parameter.  In the geometry below it
is fixed when spatial derivatives, metrics, and curvature are computed.
On real $(t,x)$ put
\begin{equation}
 h_j=\partial_x^jH,\qquad E=h_0^2+h_1^2,\qquad
 \mathcal L_j=h_{j+1}^2-h_jh_{j+2},\qquad
 S_j=\{h_{j+1}=0\},\quad j=0,1.
 \label{phy:observations}
\end{equation}
The same notation $h_j(t,z)$ denotes the holomorphic spectral derivatives.
The four consecutive components $h_0,h_1,h_2,h_3$ have one scalar source;
they are not independent fields to which arbitrary geometric constraints
may be assigned.

\begin{lemma}[Analyticity and the sign of heat time]
\label{phy:entire}
The integral \eqref{phy:heat-field} defines a jointly entire function of
$(t,z)\in\mathbb C^2$.  It is even in $z$, obeys real reflection for real
$t$, and satisfies
\begin{equation}
 H_t=-H_{zz}.
 \label{phy:backward-heat}
\end{equation}
For every real $t$, $H(t,0)>0$.  Consequently no real-time slice is
identically zero and every spatial zero has finite multiplicity.
\end{lemma}
\begin{proof}
Every term of \eqref{phy:source} is positive for $u\ge0$, since
$2\pi n^2e^{4u}-3>0$.  Dropping the negative term and comparing successive
terms of $n^4e^{-\pi n^2e^{4u}}$ gives
\[
 0<\Phi(u)<4\pi^2e^{9u-\pi e^{4u}}.
\]
Indeed the ratio of successive majorant terms is at most
$16e^{-3\pi}<1/2$.  On each compact set of complex $(t,z)$, every fixed
parameter derivative of the integrand is bounded by a constant times
\[
 (1+u)^M\exp\{C u^2+C u+9u-\pi e^{4u}\},
\]
which is integrable.  Uniform convergence and differentiation under the
integral prove joint entirety and \eqref{phy:backward-heat}.  The symmetries
follow from the cosine representation.  At real $t$ and $z=0$ the
integrand is strictly positive.  Finally, zeros of a nonzero holomorphic
function have finite order.
\end{proof}

Direct differentiation, with $x$ held fixed, gives
\begin{equation}
 \frac12E_t=-h_0h_2-h_1h_3
 =\begin{cases}\mathcal L_0,&h_1=0,\\
                 \mathcal L_1,&h_2=0.
   \end{cases}
 \label{phy:radial-energy}
\end{equation}
These are the stationary sets of the original field and its first
derivative.  A normalization of the field depending on $x$ would change
its stationary set and cannot be substituted in this identity.

\begin{remark}[The imported theta-to-RH bridge]
\label{phy:imported-bridge}
For the normalization \eqref{phy:source}--\eqref{phy:heat-field}, the earlier
investigation establishes
\[
 H(0,z)=\frac18\xi\!\left(\frac12+\frac{iz}{2}\right)
\]
and the finite-collision criterion
\begin{equation}
 \mathrm{RH}\quad\Longleftrightarrow\quad
 E(t,x)>0\quad(0<t\le1/2,\ x\in\mathbb R).
 \label{phy:imported-rh-criterion}
\end{equation}
We use this result as an imported theorem, with its normalization and
range unchanged; see \cite[section on positive-time collisions]{BakerTheta}.
Its global inputs are the de Bruijn--Newman real-zero threshold, including
its endpoint and upper bound $1/2$, the strip bound for nonreal zeros, and
a confinement estimate uniform for $t$ in compact positive-time intervals.
The last input prevents a limiting first collision from escaping to
infinite real height.  Joint analyticity and conjugation symmetry then
force a finite multiple real zero if RH fails.  Conversely, local
backward-heat splitting at a positive-time multiple zero produces nonreal
zeros immediately before it.  Thus a geometric construction alone does
not supply \eqref{phy:imported-rh-criterion}; its global theta hypotheses
remain part of the imported theorem.  The local exclusion arguments below
are proved independently of that import.
\end{remark}

\subsection{Full stationary predecessors at every finite multiplicity}
\label{phy:predecessors}

\begin{lemma}[Even-order stationary predecessor]
\label{phy:even-predecessor}
Let $W(t,x)$ be real analytic near $(T,a)$, satisfy $W_t=-W_{xx}$, and
have an exact spatial zero of order $2r$ there, where $r\ge1$.
Write $A=\partial_x^{2r}W(T,a)\ne0$ and
$B=\partial_x^{2r+1}W(T,a)$.  For sufficiently small $s\ge0$ there is
a real analytic curve $\chi(s)$ such that
\begin{align}
 W_x(T-s,\chi(s))&=0,&
 \chi(s)&=a-\frac{B}{rA}s+O(s^2),\label{phy:predecessor-location}\\
 W(T-s,\chi(s))&=\frac{A}{r!}s^r+O(s^{r+1}),&
 W_{xx}(T-s,\chi(s))&=\frac{A}{(r-1)!}s^{r-1}+O(s^r).
 \label{phy:predecessor-values}
\end{align}
In particular the central value of $W$ is nonzero for sufficiently small
$s>0$.
\end{lemma}
\begin{proof}
Analyticity and the heat equation give the convergent Taylor expansion
\begin{equation}
 W(T-s,a+y)=\sum_{p,q\ge0}
   \frac{s^p y^q}{p!q!}\,\partial_x^{2p+q}W(T,a).
 \label{phy:heat-taylor}
\end{equation}
Set $y=s\eta$ and differentiate once in $y$.  Every surviving term has
$2p+q+1\ge2r$, hence $p+q\ge r$.  It follows that
$s^{-r}W_x(T-s,a+s\eta)$ extends analytically across $s=0$.
At $s=0$ its value is
\[
 \frac{A}{(r-1)!}\eta+\frac{B}{r!}.
\]
Its derivative in $\eta$ is nonzero.  The analytic implicit-function
theorem supplies $\eta(s)$ with initial value $-B/(rA)$, proving
\eqref{phy:predecessor-location}.  Substitution in
\eqref{phy:heat-taylor} proves \eqref{phy:predecessor-values}:
the leading terms are respectively $p=r,q=0$ and, after two spatial
derivatives, $p=r-1,q=0$.
\end{proof}

\begin{proposition}[The physical stationary pole]
\label{phy:stationary-pole}
Suppose a nonzero analytic backward-heat field $H$ has exact spatial
multiplicity $m=2r+j$ at $(T,a)$, with $r\ge1$ and $j\in\{0,1\}$.
Apply Lemma~\ref{phy:even-predecessor} to $W=h_j$.  On its full stationary
predecessor, as $s\downarrow0$,
\begin{equation}
 h_{j+1}=0,\qquad E>0,\qquad
 \frac{\mathcal L_j}{E}=-\frac{r}{s}+O(1).
 \label{phy:stationary-pole-formula}
\end{equation}
For $j=0$ one has $E=W^2$ exactly.  For $j=1$,
\begin{equation}
 H(T-s,\chi(s))=O(s^{r+1}),\qquad
 E=W^2(1+O(s^2)).
 \label{phy:odd-energy-comparison}
\end{equation}
\end{proposition}
\begin{proof}
On the constructed branch $\mathcal L_j=-W W_{xx}$.
Equation~\eqref{phy:predecessor-values} gives
$W_{xx}/W=r/s+O(1)$.  If $j=0$, stationarity is $H_x=0$, so
$E=W^2$ and the claim follows.  If $j=1$, $H$ has order $2r+1$ at the
event.  In its heat Taylor expansion evaluated at $x-a=O(s)$, each
surviving term has $2p+q\ge2r+1$, which implies $p+q\ge r+1$.
Thus $H=O(s^{r+1})$, whereas $W$ is a nonzero constant times $s^r$.
This proves \eqref{phy:odd-energy-comparison} and then
\eqref{phy:stationary-pole-formula}.  All divisions are taken on the
proved nonzero predecessor, not at the event itself.
\end{proof}


\begin{theorem}[The complete local backward stationary locus]
\label{phy:complete-stationary-locus}
Let $H$ be a real analytic backward-heat field with exact spatial
multiplicity $2r+j$ at $(T,a)$, $r\ge1$, $j\in\{0,1\}$, and put $W=h_j$.
For some $s_0,d>0$, throughout
$0<s<s_0$, $|x-a|<d$, both $W(T-s,x)$ and $W_{xx}(T-s,x)$
are nonzero with the same fixed sign. The equation $W_x=0$ has
exactly one root in this rectangle at each $s$, namely the analytic
curve $x=\chi(s)$ of Lemma~\ref{phy:even-predecessor}.
The opposite parity stationary set is empty there: $H_{xx}\ne0$
when $j=0$, and $H_x\ne0$ when $j=1$.
\end{theorem}
\begin{proof}
If an analytic backward-heat field $G$ has exact even spatial order
$2k$ at the event and leading coefficient $c\ne0$, its heat Taylor
expansion gives, for $s\ge0$ and $R=\sqrt{s}+|y|$,
\[
 G(T-s,a+y)=cP_{2k}(s,y)+O(R^{2k+1}),\qquad
 P_{2k}(s,y)=\sum_{p=0}^k\frac{(2k)!s^py^{2k-2p}}{p!(2k-2p)!}.
\]
The ordinary Taylor theorem applied to $G(T-\eta^2,a+y)$ proves
the displayed remainder. Comparing coefficients with $(s+y^2)^k$
gives $P_{2k}\ge(s+y^2)^k\ge2^{-k}R^{2k}$.
Thus $G$ has sign $c$ on a sufficiently small punctured backward
half-neighborhood. For $k=0$ this follows by continuity, including
the center. Apply this argument to $W$ and $W_{xx}$, of orders
$2r$ and $2r-2$ with leading coefficients of the same sign.
The fixed sign of $W_{xx}$ makes $W_x$ strictly monotone in $x$;
the already constructed root $\chi(s)=a+O(s)$ is therefore unique.
The opposite parity equation is $W_{xx}=0$ or $W=0$, respectively.
\end{proof}

\begin{corollary}[One analytic factor for the physical first jet]
\label{phy:scalar-stationary-factor}
On the entire stationary locus of
Theorem~\ref{phy:complete-stationary-locus}, write
$W=s^ru(s)$, where $u$ is analytic and
$u(0)=\partial_x^{2r}W(T,a)/r!\ne0$.
On branch zero, $(H,H_x)=(W,0)$. On branch one,
\begin{equation}
 H_x=s^ru(s),\qquad H=s^{r+1}v(s),\qquad
 \frac{H}{H_x}=s\lambda(s),\quad \lambda=v/u,
 \label{phy:scalar-odd-factor}
\end{equation}
with analytic $v,\lambda$. Consequently the analytic ideal generated
by $(H,H_x)$ on the curve is $(W)=(s^r)$, and
\begin{equation}
 E=W^2Q(s),\qquad
 Q(s)=\begin{cases}1,&j=0,\\1+s^2\lambda(s)^2,&j=1,\end{cases}
 \qquad Q(0)=1.
 \label{phy:scalar-energy-factor}
\end{equation}
The radial size $\sqrt E=|W|\sqrt Q$ retains exact order $s^r$.
\end{corollary}
\begin{proof}
The analytic predecessor expansion gives $W=s^ru$.
For the odd branch put $h(s)=H(T-s,\chi(s))$ and
$w(s)=H_x(T-s,\chi(s))$. The heat equation and $H_{xx}=0$ give
\begin{equation}
 h'(s)=\chi'(s)w(s),\qquad
 h(s)=\int_0^s\chi'(v)w(v)\,dv.
 \label{phy:odd-stationary-transport}
\end{equation}
Since $w=s^ru$ and $h(0)=0$, the integral is $s^{r+1}$ times
an analytic function. Division by the analytic unit $u$ proves
\eqref{phy:scalar-odd-factor}; the remaining conclusions follow.
\end{proof}

In fact $w'=W_{xx}$ on the stationary curve has the sign of $w$,
so $|w|$ increases with $s$ and
\[
 \left|\frac H{H_x}\right|
 \le s\sup_{0\le v\le s}|\chi'(v)|.
\]
For an odd event, with
$\alpha=\partial_x^{2r+1}H(T,a)$ and
$\beta=\partial_x^{2r+2}H(T,a)$, the leading ratio is
\[
 \lambda(0)=-\frac{\beta}{r(r+1)\alpha}.
\]
Thus a generally asymmetric odd curve is not the exact $H_x$ axis.
These results classify the small backward stationary neighborhood,
not the forward side of the event or an entire time slice.
They are developed further in \cite{ThetaScalarGeometry}.

\begin{theorem}[Signed stationary exclusion]
\label{phy:stationary-exclusion}
Let $H$ be a real analytic backward-heat field whose spatial slices are
not identically zero.  Suppose nonnegative functions $K_0,K_1$ are locally
bounded through every candidate event in a stated time-space domain and
\begin{equation}
 \mathcal L_j\ge-K_jE\quad\hbox{on the full set }S_j,\qquad j=0,1.
 \label{phy:stationary-target}
\end{equation}
Assume the domain contains the sufficiently close predecessors of each
candidate event.  Then it contains no multiple real zero covered in this
way.  For the complete theta field, if the hypothesis covers
$(0,1/2]\times\mathbb R$, it implies RH by
Remark~\ref{phy:imported-bridge}.
\end{theorem}
\begin{proof}
Every multiple zero has a finite order $m\ge2$, uniquely expressible as
$2r+j$ with the stated ranges.  Proposition~\ref{phy:stationary-pole}
then contradicts the locally bounded coefficient $K_j$ on its corresponding
predecessor.  At the endpoint $t=1/2$, a left-hand time neighborhood
suffices.  No boundedness as $t\downarrow0$ or as $|x|\to\infty$ is required.
\end{proof}

Both branches are essential for this argument.  The backward-heat polynomial
$H(t,x)=x^3-6(t-T)x$ has
$\mathcal L_0=3x^4+36(t-T)^2\ge0$ everywhere and nevertheless has a
triple zero at $(T,0)$.  On the second stationary predecessor $x=0$,
$\mathcal L_1/E=-1/(T-t)$.

\section{Signed forms and the physical observation}\label{forms:section}

The two Laguerre expressions are polarized quadratic forms in three
consecutive derivatives.  A geometric representation must preserve that
polarization.  Replacing it by the norm of the whole jet changes the
question, because the higher derivatives need not vanish when $E$ does.
The quaternion-form working notes of Edward B.~Baker~III
\cite{BakerQuaternion} motivate the use of oriented two-form pairings.
The construction below is derived directly; it uses none of the
unfinished singular-kernel or nonabelian boundary formulas in those notes.

\subsection{The Hodge pairing in four dimensions}
Fix the Euclidean metric on $\mathbb R^4$ with coordinates $(x,y,a,b)$,
orientation $dx\wedge dy\wedge da\wedge db$, and complex coordinates
$z=x+iy$, $w=a+ib$.  This use of $a,b$ avoids confusing the auxiliary
manifold coordinates with the source variable $u$.  The Hodge star gives
\[
 \Lambda^2=\Lambda^2_+\oplus\Lambda^2_-,\qquad
 \alpha\wedge\alpha=(|\alpha_+|^2-|\alpha_-|^2)\,dV.
\]
The wedge pairing has signature $(3,3)$; the inner product
$\alpha\wedge *\alpha$ is positive definite.  Both structures are useful,
but they answer different questions.

For orthonormal forms $\sigma^+_1,\sigma^+_2\in\Lambda^2_+$ and
$\sigma^-_1\in\Lambda^2_-$, the form
\[
 \Omega=B\sigma^+_1+\tfrac12(A-C)\sigma^+_2
                     +\tfrac12(A+C)\sigma^-_1
\]
has wedge square $(B^2-AC)dV$ and squared norm
$B^2+(A^2+C^2)/2$.  This realizes a Laguerre expression algebraically,
but its norm pays for the higher jet $C$, and a variable-coefficient
version need not be closed.  We instead realize each input by a closed
form and retain the signed polarization explicitly.

\begin{proposition}[Closed self-dual inputs]\label{forms:closed}
For a holomorphic function $f(z)$, define
\[
 \alpha_f=\frac1{\sqrt2}\operatorname{Re}(f(z)\,dz\wedge dw).
\]
Then
\begin{equation}\label{forms:pairing}
 d\alpha_f=0,\qquad *\alpha_f=\alpha_f,\qquad
 \alpha_f\wedge\alpha_g=\operatorname{Re}(f\overline g)\,dV.
\end{equation}
The form is Hodge-harmonic, and in the fixed Euclidean metric
\begin{equation}\label{forms:parallel}
 |\alpha_f|^2=|f|^2,\qquad |\nabla\alpha_f|^2=2|f'|^2.
\end{equation}
In particular, it is parallel on a connected domain exactly when $f$
is constant.
\end{proposition}
\begin{proof}
Put $A=dx\wedge da-dy\wedge db$ and
$B=dx\wedge db+dy\wedge da$.  Both are self-dual,
$A^2=B^2=2dV$, and $A\wedge B=0$.  If $f=p+iq$, then
$\alpha_f=(pA-qB)/\sqrt2$, which proves the pairing and norm.
Holomorphicity gives $d(f\,dz\wedge dw)=0$.  Self-duality then implies
$d^*\alpha_f=-*d*\alpha_f=0$.  The Cauchy--Riemann equations give
$|\nabla p|^2+|\nabla q|^2=2|f'|^2$, proving the last identity.
\end{proof}

For $f=h_j(t,\cdot)$, heat time is held fixed in this exterior
derivative.  Writing $\alpha_j=\alpha_{h_j}$, we obtain
\begin{equation}\label{forms:laguerre}
 \alpha_{j+1}\wedge\alpha_{j+1}-\alpha_j\wedge\alpha_{j+2}
 =\bigl(|h_{j+1}|^2-\operatorname{Re}(h_j\overline{h_{j+2}})\bigr)dV.
\end{equation}
On $y=0$ the coefficient is $\mathcal L_j$.  This is evaluation of a
four-form coefficient on the real slice, not pullback of a four-form
onto a one-dimensional line.  The coefficient-space matrix has spectrum
$\{1,1/2,-1/2\}$.  Its negative eigenvalue persists even though all
three input forms are closed and self-dual.

\subsection{Primitives and a signed boundary observation}
If $P_f'=f$, then
\[
 \beta_f=\frac1{\sqrt2}\operatorname{Re}(P_f\,dw),\qquad
 d\beta_f=\alpha_f.
\]
Stokes' theorem gives, for a smooth bounded four-domain with the forms
defined on a neighborhood of its closure,
\begin{equation}\label{forms:stokes}
 \int_\Omega\alpha_f\wedge\alpha_g
 =\int_{\partial\Omega}\beta_f\wedge\alpha_g.
\end{equation}
This retains a mixed signed pairing.  It does not replace the boundary
term by a positive norm or remove it.

The following elementary identity makes the same distinction on a
spectral circle.  Write $\langle q\rangle=(2\pi)^{-1}
\int_0^{2\pi}q(\theta)\,d\theta$.
\begin{proposition}[Complete circle current]\label{forms:circle}
Let $f$ be holomorphic near the closed disk $|z-x_0|\le R$, real on
its real diameter, and $f'(x_0)=0$.  Set
\[
 \varphi(\theta)=f(x_0)\operatorname{Re} f(x_0+Re^{i\theta}).
\]
Then
\begin{equation}\label{forms:circle-identity}
 -f(x_0)f''(x_0)
 =-\frac{4f(x_0)^2}{R^2}
   +\frac4{R^2}\langle(1-\cos2\theta)\varphi(\theta)\rangle.
\end{equation}
For $f=h_j$ this is exactly $\mathcal L_j$ on its full stationary set.
\end{proposition}
\begin{proof}
The Taylor coefficients are real.  The mean of $\operatorname{Re} f$
is $f(x_0)$ and its second cosine coefficient has mean pairing
$\langle\cos2\theta\operatorname{Re}f\rangle=R^2 f''(x_0)/4$.
Substitution proves the formula.
\end{proof}
The current $\varphi$ is also the coefficient of
$\alpha_{f(x_0)}\wedge\alpha_f$ evaluated on the circle.  A useful bound
must control the complete weighted current relative to $E$, with all
radius costs accounted for.  Separately estimating its negative mass can
discard essential favorable cancellation.

\subsection{An integrated positive identity and its scale}
The following result, inherited from the signed-form development in
\cite{BakerTheta}, explains what positivity of a completed integral can
provide without supplying the central pointwise estimate.
\begin{theorem}[Disk completion]\label{forms:disk}
For $f$ holomorphic near $\overline{D_R(x_0)}$, define
\[
 \ell_f=|f'|^2-\operatorname{Re}(f\overline{f''}),\qquad
 g_R(z)=\left(1-\frac{(z-x_0)^2}{R^2}\right)f(z).
\]
Then
\begin{equation}\label{forms:disk-id}
 \int_{D_R}\ell_f\,dA+\frac1{2R}\int_{\partial D_R}|f|^2\,ds
 =\frac12\int_{D_R}|g_R'|^2\,dA\ge0.
\end{equation}
The boundary coefficient cannot be decreased, even for real-reflection
polynomials stationary at the center.
\end{theorem}
\begin{proof}
Write $f(x_0+R\zeta)=\sum_{n\ge0}b_n\zeta^n$ and set $b_n=0$ for
$n<0$.  Angular orthogonality gives
\begin{align*}
 \pi^{-1}\int_{D_R}\ell_f\,dA
 &=\sum_{n\ge0}n|b_n|^2-\operatorname{Re}
       \sum_{n\ge0}(n+2)b_n\overline{b_{n+2}},\\
 (2\pi R)^{-1}\int_{\partial D_R}|f|^2\,ds&=\sum_{n\ge0}|b_n|^2,\\
 \pi^{-1}\int_{D_R}|g_R'|^2\,dA&=\sum_{n\ge1}n|b_n-b_{n-2}|^2.
\end{align*}
Expansion proves \eqref{forms:disk-id}.  For
$f_N=\sum_{m=0}^N(1-m/(N+1))((z-x_0)/R)^{2m}$, direct summation gives
\[
 B_N:=\frac1{2R}\int_{\partial D_R}|f_N|^2\,ds
 =\pi\frac{(N+2)(2N+3)}{6(N+1)},\qquad
 \frac{\int_{D_R}\ell_{f_N}\,dA}{B_N}=-\frac{2N}{2N+3}\longrightarrow-1.
\]
These polynomials have $f_N(x_0)=1$ and $f_N'(x_0)=0$.
\end{proof}
Their boundary energy grows while their central first-jet energy stays
equal to one.  Thus the integrated positive completion cannot be
converted to a uniform central-energy payment for arbitrary holomorphic
fields.  The actual theta source must supply the additional restriction.

\section{Hyperk\"ahler lifts and their degenerations}\label{hk:section}

A hyperk\"ahler metric in real dimension four has a parallel compatible
triple of complex structures satisfying the quaternion relations; its
holonomy is contained in $\operatorname{Sp}(1)=\operatorname{SU}(2)$.
In particular it is Ricci-flat.  These are statements about the metric
and its parallel tensors, not a consequence of closed self-duality of
one form; see \cite{Joyce2001}.

\subsection{A flat lift with an exact physical coefficient}
\begin{proposition}[Pullback hyperk\"ahler metric]\label{hk:flat}
Where a holomorphic $f(z)$ is nonzero, the metric
\begin{equation}\label{hk:gf}
 g_f=|f(z)|^2|dz|^2+|dw|^2
\end{equation}
is locally flat and hyperk\"ahler.  Its parallel triple is
\begin{align*}
 \omega_1&=\frac i2\bigl(|f|^2dz\wedge d\overline z+
                                      dw\wedge d\overline w\bigr),\\
 \omega_2&=\operatorname{Re}(f\,dz\wedge dw),\qquad
 \omega_3=\operatorname{Im}(f\,dz\wedge dw).
\end{align*}
For $f=H+iH_z$, its real-slice coefficient is exactly $E$.
\end{proposition}
\begin{proof}
A local primitive $F'=f$ supplies the nonsingular holomorphic coordinate
$\zeta=F(z)$.  Then $g_f=|d\zeta|^2+|dw|^2$ and the displayed forms
are the standard constant hyperk\"ahler triple in those coordinates.
Equivalently they are closed and
$\omega_i\wedge\omega_j=2\delta_{ij}|f|^2dV$.
For real $x$, both $H$ and $H_x$ are real, so
$|H+iH_x|^2=H^2+H_x^2$.
\end{proof}

This construction has trivial restricted holonomy, contained in
$\operatorname{Sp}(1)$; it does not have holonomy equal to the full
group.  It is a precise geometric realization of the physical scale.
At a zero of $f$ of order $m$, however, $F-F(z_0)$ has order $m+1$ and
\eqref{hk:gf} degenerates.  After a nonsingular holomorphic change of
coordinate its curve factor is a constant times
$|\zeta|^{2m}|d\zeta|^2$.  Its metric completion has cone angle
$2\pi(m+1)$ and finite distance to the vertex.  Passing to the branched
coordinate $F$ is not a nonsingular coordinate repair on the original
domain.

At a nondegenerate physical collision $H=H_z=0$, $H_{zz}\ne0$, the
coefficient $f=H+iH_z$ has a simple zero since $f'=iH_{zz}$.  Thus the
real physical collision produces an angle-$4\pi$ degeneration.
Away from the real slice, $H+iH_z$ can vanish without both jets
vanishing; the exact equivalence with the physical detector is a
real-slice statement.  The fixed background metric remains smooth at
every such event, but its theta-dependent forms are generally harmonic
rather than parallel.

\subsection{Conformal and compact obstructions}
In dimension four the Hodge star on two-forms is unchanged by
$g\mapsto e^{2\sigma}g$.  Thus closed self-dual forms remain harmonic
under arbitrary conformal changes.  For example, $e^{2x}g_0$ has
\[
 \operatorname{Ric}=-2(dy^2+da^2+db^2),\qquad
 \operatorname{Scal}=-6e^{-2x}.
\]
It is not hyperk\"ahler, despite preserving all the closed self-dual
inputs.  The top-form pairing itself is unchanged; only its coefficient
relative to the volume changes, by a positive factor $e^{-4\sigma}$.
No sign improvement follows from this conformal replacement.

Pointwise normalization also needs care.  Where $f\ne0$,
\[
 d(\alpha_f/|f|)=-d\log|f|\wedge(\alpha_f/|f|).
\]
A nonzero self-dual two-form is symplectic, so wedge multiplication by
it is an isomorphism from one-forms to three-forms.  Consequently this
normalized form is closed only if $|f|$ is constant, which for
holomorphic $f$ forces $f$ constant.  Normalization can therefore lose
the differential equation needed for a holonomy construction.

\begin{proposition}[Compact exactness obstruction]\label{hk:compact}
On a compact oriented four-manifold without boundary, a globally exact
self-dual two-form vanishes.
\end{proposition}
\begin{proof}
If $\alpha=d\beta$ and $*\alpha=\alpha$, then
$\int |\alpha|^2dV=\int\alpha\wedge\alpha
=\int d(\beta\wedge\alpha)=0$.
\end{proof}
Our forms are exact on $\mathbb C^2$.  A nontrivial compact realization
must therefore change their global period, boundary or singularity
data; a descended form may have a primitive that fails to descend.
On the original noncompact space, the forms are independent of $w$
and are not square integrable unless zero.  Compact integration by
parts cannot be imported without controlling the missing boundaries.

There is a further restriction on extending the entire input family to
a compact hyperk\"ahler manifold.  With the orientation in which the
parallel triple is self-dual, the Weitzenb\"ock formula for a harmonic
self-dual form is
\[
 0=\nabla^*\nabla\alpha+
       (\operatorname{Scal}/3-2W^+)\alpha=\nabla^*\nabla\alpha.
\]
Here scalar curvature and $W^+$ vanish.  Integration implies
parallelness.  This makes metric-relative pairings constant, although
coordinate coefficients of a parallel form need not be constant.
For the actual family a stronger local argument applies: at a real
point with $H\ne0$, the forms $\alpha_H$ and $\alpha_{H_z}$ are
proportional with a real ratio $c$.  If both were parallel for the same
connection, their difference $\alpha_{H_z}-c\alpha_H$ would vanish
throughout the connected patch.  This forces $H_z=cH$, incompatible
with a nonzero even entire field decaying along the real axis.  This
obstruction concerns preservation of the whole particular family, not
every possible compact geometric representation.

\section{A holomorphic surface with the physical metric coefficient}
\label{phy:jet-surface}

Fix real $t$, abbreviate $f(z)=H(t,z)$, and choose a holomorphic primitive
$P'=f$.  For the entire theta field a global primitive exists.  Define
\[
 \mathcal E(z)=|f(z)|^2+|f'(z)|^2,\qquad
 \Omega_t=\{z:\mathcal E(z)>0\}.
\]
On $\Omega_t\times\mathbb C$ consider
\begin{equation}
 X_t(z,w)=(P(z),f(z),w)\in\mathbb C^3.
 \label{phy:immersion}
\end{equation}
The complex tangent vectors $(f,f',0)$ and $(0,0,1)$ are independent
exactly where $\mathcal E>0$.  Thus \eqref{phy:immersion} is a holomorphic
immersion with induced real metric and Kähler form
\begin{align}
 g_t&=\mathcal E(z)|dz|^2+|dw|^2,
 \label{phy:metric}\\
 \omega_t&=\frac{i}{2}
    \{\mathcal E\,dz\wedge d\bar z+dw\wedge d\bar w\}.
 \label{phy:kahler-form}
\end{align}
The parameter domain with this metric is a smooth Riemannian four-manifold.
Its image is an immersed complex surface; global injectivity and absence
of self-intersections are not asserted.  On the physical real slice,
$\mathcal E(x)=E(t,x)$ exactly.  At a common zero of $H$ and $H_z$, the
map loses one complex rank and \eqref{phy:metric} is not positive definite.

\subsection{Curvature and induced holonomy}

\begin{proposition}[Curvature of the jet surface]
\label{phy:curvature}
Set $D=ff''-(f')^2$.  With the real metric convention
$|dz|^2=dx^2+dy^2$, the Gaussian curvature of the first factor of
\eqref{phy:metric} is
\begin{equation}
 K_t=-\frac{2|D|^2}{\mathcal E^3}.
 \label{phy:gaussian-curvature}
\end{equation}
The scalar curvature of the four-dimensional product is $2K_t$.  Its Ricci
tensor equals $K_t$ times the curve metric in the first factor and zero in
the flat factor; its Ricci form is
\begin{equation}
 \rho_t=-i\frac{|D|^2}{\mathcal E^2}\,dz\wedge d\bar z.
 \label{phy:ricci-form}
\end{equation}
In particular, on the physical line,
\begin{equation}
 K_t(x)=-\frac{2\mathcal L_0(t,x)^2}{E(t,x)^3}.
 \label{phy:curvature-squares-sign}
\end{equation}
\end{proposition}
\begin{proof}
The two-by-two Hermitian Gram determinant gives
\begin{align*}
 \partial_z\partial_{\bar z}\log\mathcal E
 &=\frac{\mathcal E(|f'|^2+|f''|^2)
        -|f'\bar f+f''\overline{f'}|^2}{\mathcal E^2}\\
 &=\frac{|ff''-(f')^2|^2}{\mathcal E^2}.
\end{align*}
For $ds^2=\mathcal E(dx^2+dy^2)$ the Gaussian curvature is
$-(2\mathcal E)^{-1}\Delta\log\mathcal E$, and
$\Delta=4\partial_z\partial_{\bar z}$.  This proves
\eqref{phy:gaussian-curvature}.  Product curvature proves the Ricci tensor
and scalar statements.  Equivalently the Kähler Ricci form is
$-i\partial\bar\partial\log\det(g_{a\bar b})$; the irrelevant constant
in the Hermitian determinant disappears upon differentiation, giving
\eqref{phy:ricci-form}.  On the real line $D=-\mathcal L_0$, proving the
last identity.
\end{proof}

The metric has a parallel splitting into a curve and a flat complex line,
so its holonomy is contained in $U(1)\times\{1\}\subset U(2)$.
It is Ricci-flat on a connected patch precisely when $D\equiv0$ there.
That condition forces
\begin{equation}
 f(z)=A e^{bz},\qquad A\ne0,
 \label{phy:flat-exponential}
\end{equation}
on that patch.  Indeed on a subpatch where $f\ne0$ one has
$(f'/f)'=D/f^2=0$, and holomorphic identity continuation proves
\eqref{phy:flat-exponential}.  If this equality held on an open patch of a theta slice, the identity
theorem would extend it to the whole spectral plane.  A nonzero even theta
field cannot be such an exponential: evenness would force $b=0$, whereas
$H_{zz}(t,0)=-\int_0^\infty u^2e^{tu^2}\Phi(u)\,du<0$.
Thus this induced metric is not a hyperkähler metric for the theta field.
This is a statement about \eqref{phy:metric}, not a prohibition on other
auxiliary metrics or other geometric realizations.

The squaring in \eqref{phy:curvature-squares-sign} is decisive.  A lower or
upper curvature norm does not select the favorable sign of
$\mathcal L_0$.  In addition, the factor $E^{-3}$ makes intrinsic curvature
very sensitive to degeneration of the physical metric.

\subsection{Genuine calibrated inclusions and their limits}

Every locally embedded image patch of \eqref{phy:immersion} is calibrated
by the ambient Kähler calibration $\omega_0^2/2$ and hence is minimal.
The relevant algebraic and variational statements are the complex
submanifold calibration and the fundamental theorem of calibrations;
see \cite[Theorems 3.1 and 2.3]{Lotay2018}.  This conclusion uses a proved
holomorphic immersion, rather than an analogy between quaternionic
coefficients and tangent vectors.

There are literal exceptional-calibration realizations as well.  On
$\mathbb R_r\times\mathbb C^3$ use the standard $G_2$ form
\[
 \varphi_0=dr\wedge\omega_0+\operatorname{Re}\Omega_3,
 \qquad *_0\varphi_0=\tfrac12\omega_0^2-dr\wedge\operatorname{Im}\Omega_3.
\]
The image of \eqref{phy:immersion} in $\{0\}\times\mathbb C^3$ is
coassociative: the second term vanishes and the first restricts to its
induced volume.  Alternatively, regard $(P,f,w,0)$ as a surface in
$\mathbb C^4\cong\mathbb R^8$.  For the standard Cayley calibration
\begin{equation}
 \Psi_0=\tfrac12\omega_0^2+\operatorname{Re}\Omega_4,
 \label{phy:cayley-calibration}
\end{equation}
the holomorphic four-form pulls back to zero, while the first term is the
volume.  Hence this image is Cayley.  These standard inclusions are also
recorded in \cite[Propositions 6.5 and 7.3]{Lotay2018}.
The ambient metrics used here are flat, with trivial holonomy contained
in the indicated groups.  We have not produced an ambient metric whose
holonomy is exactly $G_2$ or $\operatorname{Spin}(7)$.

Calibration does not remove a branch degeneration.  If at a fixed time
\[
 f(z)=c(z-a)^m+O((z-a)^{m+1}),\qquad c\ne0,\quad m\ge2,
\]
then
\begin{equation}
 \mathcal E(z)\sim m^2|c|^2|z-a|^{2m-2}.
 \label{phy:cone-metric}
\end{equation}
The radial distance to $a$ is asymptotic to $|c||z-a|^m$, and the
circumference of a small coordinate circle is asymptotic to
$2\pi m|c||z-a|^m$.  Thus the completion has asymptotic cone angle
$2\pi m$ in the curve factor.  It is at finite distance but is not a
nondegenerate smooth point of the original metric.
One can instead make the everywhere immersive graph
$(z,P,H,w)\in\mathbb C^4$.  Its induced metric is
$(1+\mathcal E)|dz|^2+|dw|^2$, and it remains a calibrated complex surface
through a physical joint zero.  This repairs immersion by introducing a
background metric term; it does not prove nonvanishing of $E$.

\section{The marked real Gauss map retains the sign}
\label{phy:real-gauss-map}

The nonconstant tangent direction of \eqref{phy:immersion} is the
projective Gauss map
\begin{equation}
 G_0(t,z)=[H(t,z):H_z(t,z)]\in\mathbb{CP}^1.
 \label{phy:gauss-map-zero}
\end{equation}
The second tangent direction is fixed, so this projective line captures
the varying part of the full tangent plane.  Normalize the Fubini--Study
metric by
$ds_{\rm FS}^2=|d\zeta|^2/(1+|\zeta|^2)^2$.  Differentiating in either
projective chart gives the coordinate-independent identities
\begin{equation}
 G_0^*ds_{\rm FS}^2=\frac{|D|^2}{\mathcal E^2}|dz|^2,
 \qquad \rho_t=-2G_0^*\omega_{\rm FS}.
 \label{phy:fubini-study-pullback}
\end{equation}
Real reflection puts the physical image in $\mathbb{RP}^1$.  Its
orientation along the marked real spectral line is extra information
which the nonnegative metric in \eqref{phy:fubini-study-pullback} omits.

For both branches put
\begin{equation}
 E_j^{\rm jet}=h_j^2+h_{j+1}^2,\qquad
 G_j=[h_j:h_{j+1}],\qquad j=0,1.
 \label{phy:branch-gauss-map}
\end{equation}
On a real interval on which this pair never vanishes, choose a continuous
real lift $\alpha_j=\arg(h_j+i h_{j+1})$.  The lift is defined up to a
constant multiple of $2\pi$; its projective value is defined modulo $\pi$.
Its derivative and every continuously lifted increment are unambiguous.
Differentiating the argument proves
\begin{equation}
 \boxed{\partial_x\alpha_j
  =\frac{h_jh_{j+2}-h_{j+1}^2}{h_j^2+h_{j+1}^2}
  =-\frac{\mathcal L_j}{E_j^{\rm jet}}.}
 \label{phy:angle-derivative}
\end{equation}
For $j=0$, $E_0^{\rm jet}=E$.  For $j=1$ the denominator is different
and can vanish even when $E>0$.  No global regularity of $G_1$ as a
physical unit vector is being presumed.

\begin{proposition}[An oriented angular certificate]
\label{phy:angular-certificate}
Suppose prescribed functions $\kappa_j\ge0$ are locally bounded through
the covered candidate events and
\[
 \partial_x\alpha_j\le\kappa_j
 \quad\hbox{on }S_j\cap\{h_j\ne0\},\qquad j=0,1.
\]
Then $\mathcal L_j\ge-\kappa_jE$ on the full stationary sets.
Consequently the hypotheses of Theorem~\ref{phy:stationary-exclusion}
are satisfied on the same covered domain.
\end{proposition}
\begin{proof}
On $S_j$ with $h_j\ne0$, \eqref{phy:angle-derivative} gives
$\mathcal L_j\ge-\kappa_jh_j^2\ge-\kappa_jE$.
The second inequality uses both $h_j^2\le E$ for $j=0,1$ and
$\kappa_j\ge0$.  At a remaining point of $S_j$ one has $h_j=0$,
so $\mathcal L_j=0$ and the same inequality holds without defining an
angle there.
\end{proof}

For $j=0$ this certificate is the original signed stationary bound.
For $j=1$ it is in general a stronger sufficient bound, since it pays by
$h_1^2$ on $S_1$ instead of $E$.  On an odd collision predecessor their
ratio tends to one by \eqref{phy:odd-energy-comparison}.  These distinctions
are necessary when comparing an auxiliary projective metric with the
physical two-component energy.

\subsection{A sign-reversing geometric control}

Fix $\varepsilon>0$ and let
\[
 f_+(z)=\varepsilon+z^2,\quad f_-(z)=\varepsilon-z^2,
 \qquad P_\pm(z)=\varepsilon z\pm z^3/3.
\]
At the real center,
\begin{equation}
 E_\pm=\varepsilon^2,\quad
 \mathcal L_{0,+}=-2\varepsilon,\quad
 \mathcal L_{0,-}=2\varepsilon,\quad
 K_+=K_-=-8/\varepsilon^4.
 \label{phy:opposite-sign-control}
\end{equation}
The corresponding complex surfaces are unitarily congruent after rotating
the parameter:
\begin{equation}
 \operatorname{diag}(-i,1,1)\,X_+(iz,w)=X_-(z,w).
 \label{phy:unitary-congruence}
\end{equation}
This follows by substituting $iz$ into each of the two polynomial
components.  Their intrinsic and calibrated geometry therefore cannot
distinguish the two signs without retaining the marked real slice and
its orientation.  Each polynomial has an exact backward-heat extension
$f_\pm(z)\mp2(t-t_0)$ from its displayed time slice.
In particular, $H(t,z)=z^2-2(t-T)$ is calibrated at every regular point
of its associated surface but develops a double real zero at $(T,0)$.
Its predecessor at $x=0$ has $\alpha_0'=1/(T-t)$.

% Prepared for Edward Baker with substantial LLM assistance.
% Recorded model gpt-6-astra; reasoning effort ultra.

\section{Conditional exclusion by signed chords and spatial turning}
\label{phy:turning-criterion}

The marked angle in \eqref{phy:angle-derivative} retains a sign that
curvature magnitude discards.  Its most useful finite-interval observation
is a signed chord paid by the physical energy at the stationary starting
point.  We first prove a multiplicity-independent collision obstruction
for that observation, then give the finer turning profiles and the
supporting continuity criterion.  These are local heat-flow results;
the complete-theta estimate required to apply them remains open.
The detailed derivations and finite controls are recorded in
\cite{ThetaTurning,ThetaChord}.

\subsection{A fixed nearby endpoint and starting-point energy}

Put $W=h_j$, $j\in\{0,1\}$.  When $W_x(t,x)=0$ and $W$ is nonzero
on the whole interval $[x,y]$, $y>x$, a continuous angle lift satisfies
\begin{equation}
 \alpha_j(t,y)-\alpha_j(t,x)
 =\arctan\frac{W_x(t,y)}{W(t,y)}.
 \label{phy:chord-angle}
\end{equation}
Indeed, $W$ has one sign on the interval, so its lifted argument stays
in one open half-plane, with no winding ambiguity.  In particular,
\begin{equation}
 \begin{aligned}
 \mathcal C_j(t;x,y)&:=h_j(t,x)\{h_j(t,y)-h_{j+1}(t,y)\},\\
 \mathcal C_j(t;x,y)\ge0
 &\quad\Longleftrightarrow\quad
 \alpha_j(t,y)-\alpha_j(t,x)\le\frac\pi4.
 \end{aligned}
 \label{phy:chord-observation}
\end{equation}
All chord losses below use $e=E(t,x)$ at the starting point.

\begin{theorem}[The fixed-endpoint collision law]
\label{phy:fixed-endpoint-chord}
Suppose $H$ has exact spatial multiplicity $2r+j$ at $(T,a)$,
$r\ge1$, and let $W=h_j$, $\chi(s)$ be its full stationary predecessor,
and $k=\partial_x^{2r}W(T,a)/r!\ne0$.  For every sufficiently small
fixed $\delta>0$, the entire interval $[\chi(s),a+\delta]$ is
$W$-zero-free for sufficiently small $s>0$, and
\begin{equation}
 \begin{gathered}
 \boxed{\frac{\mathcal C_j(T-s;\chi(s),a+\delta)}
 {\sqrt{E(T-s,\chi(s))}}\longrightarrow-d_\delta<0,}\\
 d_\delta:=-\operatorname{sgn}(k)
       \{W(T,a+\delta)-W_x(T,a+\delta)\}>0.
 \end{gathered}
 \label{phy:fixed-endpoint-limit}
\end{equation}
\end{theorem}
\begin{proof}
Write $a_0=\partial_x^{2r}W(T,a)/(2r)!$.  The heat Taylor expansion
\eqref{phy:heat-taylor}, applied to the analytic function
$W(T-\eta^2,a+z)$, gives the ordinary two-variable remainder
\begin{align}
 W(T-\eta^2,a+z)
 &=a_0P_{2r}(\eta^2,z)
   +O\bigl((|\eta|+|z|)^{2r+1}\bigr),\label{phy:fixed-neighborhood-profile}\\
 P_{2r}(\eta^2,z)
 &=\sum_{p=0}^r\frac{(2r)!}{p!(2r-2p)!}
       \eta^{2p}z^{2r-2p},\nonumber\\
 P_{2r}(\eta^2,z)&\ge(\eta^2+z^2)^r
 \ge2^{-r}(|\eta|+|z|)^{2r}.\label{phy:positive-heat-polynomial}
\end{align}
For the first inequality, the ratio of the coefficient indexed by $p$
to the corresponding binomial coefficient is
$\prod_{i=0}^{p-1}2(2r-2i-1)\ge1$, with empty product one.
Thus the remainder divided by $P_{2r}$ is $O(|\eta|+|z|)$.
In a sufficiently small fixed neighborhood, $W$ has the sign of $a_0$
whenever $(\eta,z)\ne(0,0)$.  Since $\chi(s)-a=O(s)$, this proves
zero-freeness on the whole asserted interval, not just at its endpoints.

At time $T$,
\[
 W(T,a+\delta)-W_x(T,a+\delta)
 =-2r a_0\delta^{2r-1}+O(\delta^{2r}),
\]
which has the opposite sign from $k$ for every sufficiently small
positive $\delta$.  Finally,
\eqref{phy:predecessor-values} and \eqref{phy:odd-energy-comparison} give
\begin{equation}
 W(T-s,\chi(s))=ks^r(1+O(s)),\qquad
 E(T-s,\chi(s))=k^2s^{2r}(1+O(s)).
 \label{phy:chord-start-energy}
\end{equation}
Consequently $W(T-s,\chi(s))/\sqrt{E(T-s,\chi(s))}\to\operatorname{sgn}(k)$;
analyticity at the fixed endpoint proves \eqref{phy:fixed-endpoint-limit}.
\end{proof}

\begin{theorem}[Starting-energy chord exclusion]
\label{phy:chord-exclusion}
Let $H$ be a real analytic backward-heat field with no identically zero
spatial slice.  Near each covered candidate event suppose there are a
fixed length $\ell_0>0$, $e_0>0$, and prescribed nonnegative functions
$\omega_j(e)=o(\sqrt e)$ as $e\downarrow0$, such that
\begin{equation}
 \mathcal C_j(t;x,y)\ge-\omega_j(E(t,x)),\qquad j=0,1,
 \label{phy:chord-target}
\end{equation}
uniformly whenever $h_{j+1}(t,x)=0$, $0<y-x\le\ell_0$,
$0<E(t,x)\le e_0$, and $h_j$ is nonzero on the whole segment, which
must lie in the covered neighborhood.  Assume that neighborhood contains
the sufficiently close predecessors and these short intervals.
Then no covered finite multiple real zero occurs.  In particular, a
bound $\mathcal C_j\ge-K_jE(t,x)^\gamma$ suffices for any fixed
$\gamma>1/2$, with $K_j\ge0$ locally bounded uniformly in both endpoints.
For the complete theta field, coverage of $0<T\le1/2$ implies RH by
Remark~\ref{phy:imported-bridge}.
\end{theorem}
\begin{proof}
At a putative event use its parity branch and choose $\delta>0$ in
Theorem~\ref{phy:fixed-endpoint-chord}, smaller than $\ell_0/2$ and
inside the covered neighborhood.  The lengths $a+\delta-\chi(s)$ are
positive and smaller than $\ell_0$ for small $s>0$.  Divide
\eqref{phy:chord-target} by the square root of the starting energy and
use \eqref{phy:chord-start-energy}.  Its lower bound tends to zero,
contradicting the strictly negative limit
\eqref{phy:fixed-endpoint-limit}.  At a terminal time a left-hand time
neighborhood suffices.
\end{proof}

The fixed positive allowable length and uniformity through the event
are essential hypotheses.  A bound available only on one prescribed
shrinking family need not cover the fixed nearby endpoint used here.
Both branches are required because the event multiplicity can have
either parity.

\begin{remark}[Comparison with stationary and angular estimates]
\label{phy:criterion-comparison}
At multiplicity $2r+j$, the stationary law and the fixed-endpoint law are
\[
 \mathcal L_j\sim-\mathrm{const}\,E^{1-1/(2r)},\qquad
 \mathcal C_j\sim-\mathrm{const}\,E^{1/2}.
\]
Thus any fixed exponent $\gamma>1/2$ works in the chord target at every
finite multiplicity.  For a stationary loss $KE^\gamma$, any fixed
$\gamma<1$ eventually becomes too weak as $r$ grows.  This comparison
does not order the full hypotheses: chord estimates concern whole
intervals, whereas Theorem~\ref{phy:stationary-exclusion} imposes a
condition only on the stationary sets.
\end{remark}

\begin{proposition}[Sharpness of the energy scale and location]
\label{phy:chord-sharpness}
The exponent $1/2$ with an arbitrary bounded coefficient, or a bounded
loss paid by endpoint energy, does not exclude collisions in general.
\end{proposition}
\begin{proof}
The backward-heat field $H(T-s,x)=x^2+2s$ has a double collision,
stationary predecessor $x=0$, and zero-free rightward intervals for
every $s>0$.  With $y=\ell\ge0$,
\[
 e=E(T-s,0)=4s^2,\qquad
 \mathcal C_0=2s(2s+\ell^2-2\ell)\ge-2s=-\sqrt e.
\]
For endpoint energy put $q=H_x(T-s,\ell)/H(T-s,\ell)\ge0$ and
$b=H(T-s,0)/H(T-s,\ell)\in(0,1]$.  Then
\[
 \frac{\mathcal C_0}{E(T-s,\ell)}
 =b\frac{1-q}{1+q^2}\ge-\frac14,
 \qquad
 1-q+\frac{1+q^2}{4}=\frac{(q-2)^2+1}{4}>0.
\]
For $q\le1$ the quotient is nonnegative; for $q>1$, multiplication
by $b\le1$ improves the lower bound.  The still weaker payment by the
sum of starting and endpoint energies also permits this collision.
\end{proof}

\subsection{Transfer from a complete family of cutoff estimates}

The chord criterion permits quantitative losses that need not be
linear in $E$. For the transfer theorem in this subsection, choose the
cutoff only after proving a single prescribed family of estimates for
every integer cutoff. The endpoint-only criterion below has weaker
cutoff requirements.

\begin{proposition}[All-cutoff transfer]
\label{phy:cutoff-transfer}
Let $Q$ be a stationary observation with $e=E(t,x)$, or a chord
observation with $e=E(t,x)$ at its starting point.  Suppose
\begin{equation}
 Q\ge-Ae\exp(aN^\beta)-B\exp(-cN^2)
 \quad\hbox{for every integer }N\ge N_0,
 \label{phy:all-cutoff-family}
\end{equation}
where $A,a,B\ge0$, $c>0$, and $N_0\ge0$ is an integer.
Assume fixed local upper bounds $A_*,a_*,B_*,N_*$, with $N_*\ge1$,
and a positive lower bound $c_0$ for $c$.  For chords all these bounds
are uniform over the admissible intervals of
Theorem~\ref{phy:chord-exclusion}.  The coefficient rules are prescribed
before selecting the cutoff and are not fitted to an observed energy
floor.

If $0<\beta<2$, set
$M=\max(1,B_*)$, $L(e)=1+\max(0,\log(M/e))$, and
$D=N_*+1+c_0^{-1/2}$.  For $e>0$,
\begin{equation}
 Q\ge-(A_*+1)e\exp\{a_*D^\beta L(e)^{\beta/2}\}.
 \label{phy:subgaussian-transfer}
\end{equation}
This excludes finite collisions for the stationary target with its
predecessor coverage, and for the chord target with its full interval
coverage.

If $\beta=2$ and $a_*<c_0$, put
\[
 \kappa=a_*+c_0,\qquad p=\frac{c_0}{\kappa}>\frac12,
 \qquad D_2=N_*+1.
\]
For $0<e\le M$, with $L_2=\log(M/e)$, one has
\begin{equation}
 Q\ge-M^{1-p}e^p
 \left[A_*\exp\left\{a_*D_2^2+
        2a_*D_2\sqrt{L_2/\kappa}\right\}+1\right].
 \label{phy:gaussian-gap-transfer}
\end{equation}
This loss is $o(\sqrt e)$ and therefore excludes collisions for the
chord target.  Equivalently, every $1/2<\gamma<p$ yields a bound
$Q\ge-K_\gamma e^\gamma$ for sufficiently small $e$.
\end{proposition}
\begin{proof}
For the first claim choose
\[
 N=\left\lceil\max\left(N_*,
       \sqrt{\max(0,\log(M/e))/c_0}\right)\right\rceil.
\]
Then $B\exp(-cN^2)\le e$ and $N\le D\sqrt{L(e)}$, proving
\eqref{phy:subgaussian-transfer}.  Its relative loss is
$\exp(o(\log(1/e)))$.  Along \eqref{phy:chord-start-energy} it is
$s^{-o(1)}$, which cannot absorb the stationary pole $r/s$.
Its absolute loss is also $o(\sqrt e)$, so
Theorem~\ref{phy:chord-exclusion} applies.

For the Gaussian claim choose
$N=\lceil\max(N_*,\sqrt{L_2/\kappa})\rceil$.
The inequalities
$\sqrt{L_2/\kappa}\le N\le D_2+\sqrt{L_2/\kappa}$ give
\eqref{phy:gaussian-gap-transfer} directly.  The factor
$\exp(O(\sqrt{\log(1/e)}))$ pays integer rounding and the cutoff
threshold.  It is dominated by $e^{-(p-\gamma)}$ for any fixed
$\gamma<p$ as $e\downarrow0$.  At $e=0$, letting $N\to\infty$ in
\eqref{phy:all-cutoff-family} gives $Q\ge0$.
\end{proof}

The strict local rate gap $a_*<c_0$ cannot be replaced here by a
nonuniform pointwise gap that closes at the event.  Nor may the
rounding factor in \eqref{phy:gaussian-gap-transfer} be dropped when
claiming an exact $O(e^p)$ loss.  At equality of the rates the preceding
double-collision control satisfies, for every real $N$,
\[
 e\exp(N^2)+\exp(-N^2)\ge2\sqrt e=4s,
 \qquad \mathcal L_0=-4s,\qquad \mathcal C_0\ge-2s.
\]
Thus both observations obey \eqref{phy:all-cutoff-family} with
$\beta=2$ and $A=B=a=c=1$ despite the collision.  These transfer
statements are conditional analytic implications, not consequences of
checking finitely many cutoffs.  A separately truncated observation
requires its own complete error estimate; a proved error
$Q_0\sqrt e\exp(-c_1N^2)$ can be paid by
$e/2+(Q_0^2/2)\exp(-2c_1N^2)$.

\subsection{Exponential averages and complete-start cutoff transfer}
\label{obs:exponential-average}

The endpoint derivative can be removed by a positive average without
losing the collision obstruction.  For $W=h_j$, $j=0,1$, define
\begin{equation}
 \mathcal A_{j,\delta}(t,x)
 =W(t,x)^2-e^{-\delta}W(t,x)W(t,x+\delta),\qquad \delta>0.
 \label{obs:average-definition}
\end{equation}
The fundamental theorem of calculus gives, without stationarity,
\begin{equation}
 \mathcal A_{j,\delta}(t,x)
 =\int_0^\delta e^{-z}\mathcal C_j(t;x,x+z)\,dz.
 \label{obs:average-identity}
\end{equation}
Indeed, $\partial_z(e^{-z}W(t,x+z))=e^{-z}(W_x-W)(t,x+z)$.
The normalization by $1-e^{-\delta}$ is a positive probability
average, but its density does not vanish at the endpoints; no rapid
Fourier-decay claim for a compactly supported smooth density is used.

\begin{proposition}[Fixed-length averaged collision law]
\label{obs:average-collision}
At a multiplicity $2r+j$ event $(T,a)$, $r\ge1$, let $\chi(s)$ be the full
stationary predecessor and set
$a_0=\partial_x^{2r}W(T,a)/(2r)!\ne0$.
For every sufficiently small fixed $\delta>0$, the whole interval
$[\chi(s),\chi(s)+\delta]$ is $W$-zero-free for sufficiently small
$s>0$, and
\begin{align}
 \frac{\mathcal A_{j,\delta}(T-s,\chi(s))}
 {\sqrt{E(T-s,\chi(s))}}&\longrightarrow-d_{j,\delta}<0,
 \label{obs:average-collision-limit}\\
 d_{j,\delta}
 &=e^{-\delta}\operatorname{sgn}(a_0)W(T,a+\delta)
 =e^{-\delta}|a_0|\delta^{2r}(1+O(\delta)).\nonumber
\end{align}
The energy is $E=H^2+H_x^2$ on both branches.
\end{proposition}
\begin{proof}
The positive heat polynomial in the proof of
Theorem~\ref{phy:fixed-endpoint-chord} proves the whole-interval
assertion, since $\chi(s)-a=O(s)$.
Equation~\eqref{phy:chord-start-energy} gives
$W/\sqrt E\to\operatorname{sgn}(a_0)$ on the predecessor.
Moreover $W^2/\sqrt E\le\sqrt E\to0$.
Continuity at the moving endpoint, whose limit is $a+\delta$, proves
the limit; the Taylor expansion
$W(T,a+\delta)=a_0\delta^{2r}+O(\delta^{2r+1})$ proves its sign.
\end{proof}

For theta, let $\Phi_N$ be the first $N$ summands of
\eqref{phy:source}, and put
\[
 W_N(t,y)=\partial_y^j\int_0^\infty e^{tu^2}\Phi_N(u)\cos(yu)\,du.
\]
Truncate only the endpoint:
\begin{equation}
 \mathcal A_{j,\delta,N}
 =W(t,x)^2-e^{-\delta}W(t,x)W_N(t,x+\delta).
 \label{obs:average-endpoint-cutoff}
\end{equation}
Both the starting $W$ and stationarity $h_{j+1}(t,x)=0$ remain
complete.  For $t\le T_*$ and real $y$, the source estimate gives
\begin{align}
 |W(t,y)-W_N(t,y)|
 &\le B_j(T_*)e^{-\pi(N+1)^2/4},\label{obs:average-value-tail}\\
 B_j(T_*)&=2\pi^2\left(\sum_{n\ge1}n^4e^{-\pi n^2/4}\right)
 \int_0^\infty u^j e^{T_*u^2+9u-(\pi/2)e^{4u}}\,du<\infty.
 \nonumber
\end{align}
To check this bound, dominate the $n$th source summand by
$2\pi^2n^4e^{9u-\pi n^2e^{4u}}$ and use
$n^2e^{4u}\ge(n^2+e^{4u})/2$.
Splitting the remaining Gaussian sum into two equal exponent shares
extracts $e^{-\pi(N+1)^2/4}$.
Consequently
\begin{equation}
 |\mathcal A_{j,\delta}-\mathcal A_{j,\delta,N}|
 \le e^{-\delta}B_j(T_*)e^{-\pi(N+1)^2/4}\sqrt{E(t,x)}.
 \label{obs:average-mixed-tail}
\end{equation}
Preserving the complete starting factor is what produces $\sqrt E$;
a separately truncated starting factor does not inherit this estimate.

\begin{theorem}[Fixed-length mixed-error exclusion]
\label{obs:average-mixed-exclusion}
For the complete theta field, fix prescribed positive lengths $\delta_k\to0$.
Near each covered candidate event, for each branch $j=0,1$ and every
sufficiently large fixed $k$, suppose there is a prescribed unbounded
cutoff sequence $N_{k,\ell}\to\infty$ with estimates
\begin{equation}
 \mathcal A_{j,\delta_k,N_{k,\ell}}(t,x)
 \ge-K_{j,k,\ell}(t,x)E(t,x)
       -\varepsilon_{j,k,\ell}\sqrt{E(t,x)}.
 \label{obs:average-mixed-hypothesis}
\end{equation}
Here $K_{j,k,\ell}\ge0$ has a finite local upper bound through the
event for each fixed pair $(k,\ell)$, and the nonnegative local error
bounds satisfy $\varepsilon_{j,k,\ell}\to0$ as $\ell\to\infty$ for
each fixed $j,k$.
The estimates hold on declared full-stationary domains containing all
sufficiently close predecessors and their entire $W$-zero-free intervals
of length $\delta_k$.
Each fixed-index time threshold may differ; the error bounds must be
uniform over these approaching starts, and the coefficient rules must
be prescribed independently of an observed positive energy floor.
Then no covered finite multiple real zero occurs.
\end{theorem}
\begin{proof}
At a putative event choose its parity branch, and first fix $k$ so
that Proposition~\ref{obs:average-collision} gives
$d_{j,\delta_k}>0$ and the interval lies in the covered neighborhood.
Next choose one fixed $\ell$ with
\[
 \varepsilon_{j,k,\ell}
 +e^{-\delta_k}B_j(T_*)e^{-\pi(N_{k,\ell}+1)^2/4}
 <d_{j,\delta_k}/2.
\]
Combine \eqref{obs:average-mixed-hypothesis} and
\eqref{obs:average-mixed-tail}, divide by $\sqrt E$, and let
$s\downarrow0$ along the full predecessor with both indices fixed.
Local boundedness makes $K_{j,k,\ell}\sqrt E\to0$.
The lower bound contradicts \eqref{obs:average-collision-limit}.
\end{proof}

No growth restriction on $K_{j,k,\ell}$ as the indices increase is
needed, and an unbounded certified cutoff subsequence suffices.
The proof uses neither a simultaneous cutoff/time limit nor one
constant uniform in spectral height or as the event time tends to zero.
Every positive-time finite-height event must nevertheless be covered.
For the complete theta field, coverage of $0<T\le1/2$ has the
consequence stated in Remark~\ref{phy:imported-bridge}.

If only one shrinking length $\delta_N$ is available per cutoff,
the same argument works when the total mixed-error coefficient obeys
$\varepsilon_N=o(\delta_N^m)$ for every fixed positive integer $m$.
Indeed, choose a single $N$ with error smaller than the collision
signal $d_{j,\delta_N}\asymp\delta_N^{2r}$ before taking the time
limit.  A Gaussian total mixed error permits $\delta_N=N^{-q}$ for
every fixed $q>0$; an algebraic total error suffices for
$\delta_N=1/\log(N+2)$.
The independent-cutoff version requires only convergence at each
fixed length, but asks for a broader certificate domain.
These refinements are developed in
\cite{ThetaFactorChord,ThetaExponential,ThetaResearchPaths}.
The theta core estimate \eqref{obs:average-mixed-hypothesis} remains
unproved; the tail bound alone supplies no favorable sign.

\subsection{The inner profile and quarter-turn concentration}

\begin{lemma}[The real scaled heat profiles]
\label{phy:scaled-profile}
Use the notation of Theorem~\ref{phy:fixed-endpoint-chord}.  Uniformly
on each fixed compact interval of $\xi$,
\begin{align}
 W(T-s,\chi(s)+s\xi)&=ks^r\{1+O(s)\},\nonumber\\
 W_x(T-s,\chi(s)+s\xi)&=ks^r\{r\xi+O(s)\},\nonumber\\
 W_{xx}(T-s,\chi(s)+s\xi)&=ks^{r-1}\{r+O(s)\}.
 \label{phy:inner-jet-profile}
\end{align}
On each fixed interval $0\le y\le c$, $c>0$, one also has
\begin{align}
 W(T-s,\chi(s)+\sqrt{s}\,y)
 &=ks^r\{F_r(y)+O(\sqrt{s})\},\label{phy:scaled-W}\\
 W_x(T-s,\chi(s)+\sqrt{s}\,y)
 &=ks^{r-1/2}\{F_r'(y)+O(\sqrt{s})\},\label{phy:scaled-Wx}\\
 F_r(y)&=\sum_{\ell=0}^r
       \frac{r!}{(r-\ell)!(2\ell)!}y^{2\ell}.
 \label{phy:profile-polynomial}
\end{align}
In particular, $W$ has the sign of $k$ on each such whole interval
for small $s>0$; $F_r\ge1$ and $F_r'(c)>0$.
\end{lemma}
\begin{proof}
In \eqref{phy:heat-taylor} put
$x-a=s(\eta(s)+\xi)$, where $\eta(s)=(\chi(s)-a)/s$ is analytic.
The leading coefficients of $W,W_x,W_{xx}$ after removing respectively
$s^r,s^r,s^{r-1}$ are $k,kr\xi,kr$; for $W_x$, the constant term
cancels because
$\eta(0)=-\partial_x^{2r+1}W(T,a)/(r\partial_x^{2r}W(T,a))$.
The resulting analytic expansions and their fixed-order derivatives
are uniform on compact $\xi$ intervals.  For the parabolic scale put
$s=\eta^2$ instead.  The coefficient of $\eta^{2r}$ in
$W(T-\eta^2,\chi(\eta^2)+\eta y)$ is $kF_r(y)$; the displacement
$\chi(\eta^2)-a=O(\eta^2)$ does not affect it.  Uniform analytic
remainders and their first $y$ derivatives prove the two expansions.
Positivity of the displayed polynomial proves the sign assertions.
\end{proof}

Equations \eqref{phy:angle-derivative} and \eqref{phy:inner-jet-profile}
give the finer angular core, uniformly on compact $\xi$ intervals:
\begin{align}
 \alpha_j(T-s,\chi(s)+s\xi)-\alpha_j(T-s,\chi(s))
 &=\arctan(r\xi)+O(s),\nonumber\\
 s\,\partial_x\alpha_j(T-s,\chi(s)+s\xi)
 &=\frac{r}{1+r^2\xi^2}+O(s).
 \label{phy:inner-cauchy-profile}
\end{align}
These compact-window limits assert no weak limit on a fixed spatial
neighborhood and require no unproved uniform tail estimate.
For an odd collision the same Taylor counting gives
$H(T-s,\chi(s)+s\xi)=O(s^{r+1})$, so
\begin{equation}
 \frac{E_1^{\rm jet}}{E}(T-s,\chi(s)+s\xi)
 \longrightarrow1+r^2\xi^2.
 \label{phy:inner-odd-energy}
\end{equation}
Thus agreement of the two energies at the stationary center cannot be
extended across a nonzero inner window without retaining this factor.

\begin{theorem}[Positive quarter-turn concentration]
\label{phy:quarter-turn}
For each fixed $c>0$, the collision and predecessor above satisfy
\begin{equation}
 \alpha_j(T-s,\chi(s)+c\sqrt{s})
 -\alpha_j(T-s,\chi(s))\longrightarrow\frac\pi2.
 \label{phy:quarter-turn-limit}
\end{equation}
The defining pair is regular on the whole interval, including on branch
$j=1$.  The chord limits at the two shrinking scales are
\begin{align}
 \frac{\mathcal C_j(T-s;\chi(s),\chi(s)+Cs/r)}
 {E(T-s,\chi(s))}&\longrightarrow1-C\quad(C>0),\nonumber\\
 \frac{\mathcal C_j(T-s;\chi(s),\chi(s)+c\sqrt{s})}
 {E(T-s,\chi(s))}&=-\frac{F_r'(c)}{\sqrt{s}}+O(1).
 \label{phy:shrinking-chord-limits}
\end{align}
\end{theorem}
\begin{proof}
Zero-freeness follows from Lemma~\ref{phy:scaled-profile}.
At the parabolic endpoint,
\[
 W_x/W=s^{-1/2}(F_r'(c)/F_r(c))(1+O(\sqrt{s}))\to+\infty.
\]
Equation~\eqref{phy:chord-angle} proves the angle limit with either
sign of $k$.  Substitution of the two jet expansions into
\eqref{phy:chord-observation}, using \eqref{phy:chord-start-energy},
proves both chord limits.
\end{proof}

The finite inner limit alone can be absorbed by a growing relative
energy allowance.  A bound uniform on all short intervals must also
hold on the parabolic windows and at the fixed nearby endpoint, where
the stronger obstructions apply.  This interval quantifier explains
why the cutoff conclusions above are stronger than a test on one
fixed inner window; see \cite{ThetaChord}.

\subsection{A supporting one-sided continuity criterion}

\begin{definition}[One-sided modulus from stationary starting points]
\label{phy:turning-modulus}
A field has a local one-sided turning modulus on a declared domain if,
near each covered candidate event, there are a fixed $\ell_0>0$ and a
prescribed nonnegative function $\omega(\ell)\to0$ as $\ell\downarrow0$,
uniform in time and starting point, such that for $j=0,1$
\begin{equation}
 \alpha_j(t,x+\ell)-\alpha_j(t,x)\le\omega(\ell)
 \label{phy:modulus-condition}
\end{equation}
whenever $0<\ell\le\ell_0$, $h_{j+1}(t,x)=0$, $h_j(t,x)\ne0$,
the segment lies in the covered neighborhood, and
$(h_j,h_{j+1})\ne(0,0)$ throughout it.  The increment uses its
continuous lift along the segment.
\end{definition}

\begin{theorem}[Oriented-continuity exclusion]
\label{phy:turning-exclusion}
Let $H$ be a real analytic backward-heat field with nonzero spatial
slices.  If Definition~\ref{phy:turning-modulus} holds with coverage of
the sufficiently close full predecessors and their short rightward
intervals, then no covered finite multiple real zero occurs.
For the complete theta field, coverage of $0<T\le1/2$ implies RH by
Remark~\ref{phy:imported-bridge}.
\end{theorem}
\begin{proof}
At a candidate event take $\ell=c\sqrt{s}$ on its parity branch.
Lemma~\ref{phy:scaled-profile} supplies the required regularity on the
whole interval.  The upper bound tends to zero, contradicting
Theorem~\ref{phy:quarter-turn}.
\end{proof}

On every regular test interval the observed increment is exactly
\begin{equation}
 \int_x^{x+\ell}
 \frac{h_jh_{j+2}-h_{j+1}^2}{h_j^2+h_{j+1}^2}(t,y)\,dy.
 \label{phy:signed-turning-integral}
\end{equation}
A uniform upper bound on the angular derivative throughout these
intervals implies a linear modulus.  More generally a uniform strict
gap below $\pi/2$ on all sufficiently short regular intervals already
contradicts \eqref{phy:quarter-turn-limit}.  Both formulations retain
the orientation and allow negative turning; neither follows from a
curvature-magnitude bound or from continuity on individual regular
slices.  No global regularity of the second branch is assumed.

The resulting hierarchy identifies concrete conditional source targets.
The averaged criterion in Theorem~\ref{obs:average-mixed-exclusion}
retains the sharp starting-energy scale while relaxing the endpoint
observations and cutoff growth required of a certificate.  The exact polynomial controls in \cite{ThetaTurning,ThetaChord}
test the local identities and their scope; they establish no sign bound
for the theta source.  The remaining task is a complete-theta estimate
with the required signs, starting-point energy, and uniform interval
coverage.


\section{The complete theta source on a curved elliptic family}
\label{theta:section}

The preceding geometric realizations retain the physical jets, but a local
realization alone does not distinguish the theta source from a collision
control.  We now place the complete source in a family carrying its actual
integer lattice.  There are three separate structures: the holomorphic theta
line bundle, the physical heat-weighted transform, and a nonflat
hyperk\"ahler metric on the elliptic family.  Their relations below are exact.
They do not yet give the signed estimate at the physical energy scale.
Our source normalization is that of~\cite{BakerTheta}.

\subsection{The complete primitive and its boundary contribution}
\label{theta:primitive}

For $u\geq0$ set
\begin{equation}
 \Phi(u)=\sum_{n\geq1}
  \bigl(2\pi^2n^4e^{9u}-3\pi n^2e^{5u}\bigr)
                  e^{-\pi n^2e^{4u}},
 \qquad
 H(t,x)=\int_0^\infty e^{tu^2}\Phi(u)\cos(xu)\,du.
 \label{theta:physical-source}
\end{equation}
Every fixed source derivative decays faster than any Gaussian as
$u\to+\infty$.  The sums and integrals below therefore converge locally
uniformly with all stated derivatives.  In particular, $H$ is entire in
$(t,x)\in\mathbb C^2$ and satisfies $H_t=-H_{xx}$.

Write
\begin{equation}
 \vartheta(z,\tau)=\sum_{n\in\mathbb Z}
                   e^{\pi i n^2\tau+2\pi i nz},\qquad \Im\tau>0,
 \quad
 P(u)=e^u\vartheta(0,ie^{4u}),\quad
 \psi(u)=e^u\sum_{n\geq1}e^{-\pi n^2e^{4u}}.
 \label{theta:primitive-definitions}
\end{equation}
Thus $P=e^u+2\psi$ on the half-line.  The Gaussian Poisson identity
$\vartheta(0,i/y)=\sqrt y\,\vartheta(0,iy)$ implies $P(-u)=P(u)$.
The function $P$ grows like $e^u$ at positive infinity; it is $\psi$, not
$P$, that supplies a convergent positive-time primitive transform.

\begin{proposition}[Complete differential reconstruction]
\label{theta:reconstruction}
The complete source and primitive satisfy
\begin{equation}
 \Phi=\frac{P''-P}{16}=\frac{\psi''-\psi}{8},
 \qquad \psi^{(2k+1)}(0)=-\frac12\quad(k\geq0).
 \label{theta:primitive-ode}
\end{equation}
For
\begin{equation}
 J(t,x)=\int_0^\infty e^{tu^2}\psi(u)\cos(xu)\,du
 \label{theta:J-definition}
\end{equation}
one has the exact identity
\begin{equation}
 \boxed{\displaystyle
 8H=\frac12+(2t-x^2-1)J+4txJ_x-4t^2J_{xx}.}
 \label{theta:reconstruction-identity}
\end{equation}
All the theta indices occur in these identities.
\end{proposition}

\begin{proof}
Termwise differentiation of $P$ gives the first identity in
\eqref{theta:primitive-ode}; $(e^u)''-e^u=0$ gives the second expression.
Evenness of $P=e^u+2\psi$ gives every displayed endpoint derivative.
For $F(u)=e^{tu^2}\cos(xu)$, two integrations by parts yield
\[
 \int_0^\infty\psi''F\,du
 =-\psi'(0)F(0)+\psi(0)F'(0)+\int_0^\infty\psi F''\,du.
\]
The upper boundary terms vanish, $F(0)=1$, and $F'(0)=0$.  Directly,
\[
 F''-F=e^{tu^2}
 \bigl[(2t+4t^2u^2-x^2-1)\cos(xu)-4txu\sin(xu)\bigr].
\]
Insert $\psi'(0)=-1/2$ and differentiate~\eqref{theta:J-definition}.
This gives~\eqref{theta:reconstruction-identity}, including its boundary
constant.
\end{proof}

\begin{proposition}[Positive primitive; signed reconstruction]
\label{theta:primitive-positive}
For $0\leq t\leq1/2$ and every real $x$, $J(t,x)>0$.  Uniformly on this
heat interval,
\begin{equation}
 J(t,x)=\frac1{2x^2}-\frac{1+6t}{2x^4}+O(x^{-6})
 \qquad(|x|\to\infty).
 \label{theta:J-asymptotic}
\end{equation}
At time zero the exact reconstruction is
\begin{equation}
 J(0,x)=\frac1{2(1+x^2)}-\frac{8H(0,x)}{1+x^2}.
 \label{theta:J-time-zero}
\end{equation}
\end{proposition}

\begin{proof}
Let $q_n=\pi n^2e^{4u}$ and
$f_n(u)=\exp(tu^2+u-q_n)$, so that $f=e^{tu^2}\psi=\sum_{n\geq1}f_n$.
For $u\geq0$ and $0\leq t\leq1/2$,
$2tu+1\leq u+1\leq q_n$.  Hence
\[
 \frac{f_n''}{f_n}
 =(2tu+1-4q_n)^2+2t-16q_n
 \geq9q_n^2-16q_n>0.
\]
Thus $f''>0$.  Moreover $f'(0)=-1/2$ and $f'(\infty)=0$.
For $x\ne0$, integration by parts gives
\begin{equation}
 J(t,x)=\frac1{x^2}\int_0^\infty f''(u)(1-\cos(xu))\,du>0.
 \label{theta:convex-cosine}
\end{equation}
At $x=0$ the defining integral is positive.  Repeated integration by parts
uses $f'(0)=-1/2$ and $f'''(0)=-(1+6t)/2$ to give
\eqref{theta:J-asymptotic}.  The sixth derivative is uniformly integrable
on the stated compact heat interval, which bounds its remainder.
Finally~\eqref{theta:J-time-zero} is~\eqref{theta:reconstruction-identity}
with $t=0$.
\end{proof}

\begin{remark}[What this positivity does not prove]
\label{theta:primitive-limit}
The physical field is reconstructed from a signed differential expression
in a strictly positive channel.  At $t=0$ it is the difference between that
channel and an explicit rational profile.  The primitive-compatible
collision controls of~\cite{BakerControl}, have the same boundary
constant, strict primitive positivity, and every algebraic asymptotic
coefficient of $J$, while possessing a positive-time double zero.
Their source is not the integer theta series.  Accordingly, the next
construction retains the theta line bundle and its section, rather than
using primitive positivity alone.
\end{remark}

\subsection{The theta line bundle, its connection, and the fourth jet}
\label{theta:line-bundle}

Use the marked elliptic family
\[
 \mathcal X=(\mathbb H\times\mathbb C)/\mathbb Z^2,
 \qquad (m,n):(\tau,z)\longmapsto(\tau,z+m+n\tau),
\]
where $\mathbb H$ is the upper half-plane.  Write $\tau=p+iq$,
$z=r+\tau s$, so that $q>0$ and $r,s$ are real fiber coordinates modulo
one.  The factors
\begin{equation}
 j_{m,n}(\tau,z)=e^{-\pi i n^2\tau-2\pi i nz}
 \label{theta:multipliers}
\end{equation}
define a holomorphic line bundle $L\to\mathcal X$, and reindexing the
series gives
\begin{equation}
 \vartheta(z+m+n\tau,\tau)=j_{m,n}(\tau,z)\vartheta(z,\tau),
 \qquad 4\pi i\,\partial_\tau\vartheta=\partial_z^2\vartheta.
 \label{theta:automorphy-heat}
\end{equation}
The second equality follows term by term.  These standard theta-section
and heat-connection structures are described in
\cite[Section~6]{AndersenBlaavand2011}; the explicit normalizations needed
here are derived below.

Ordinary derivatives are jets, not separate copies of the same
holomorphic section.  For every nonnegative integer $k$,
\begin{equation}
 (\partial_z^k\vartheta)(z+m+n\tau,\tau)
 =j_{m,n}(\tau,z)
   \sum_{a=0}^k\binom ka(-2\pi in)^{k-a}
                          (\partial_z^a\vartheta)(z,\tau).
 \label{theta:jet-automorphy}
\end{equation}
Thus a global description must keep the lower-jet terms or choose a
connection.  The following choice also identifies the actual curvature
of the theta line bundle.

\begin{proposition}[Hermitian theta connection]
\label{theta:chern}
In the holomorphic frame on $\mathbb H\times\mathbb C$, the metric
\begin{equation}
 h(\tau,z)=\exp\!\left(-\frac{2\pi(\Im z)^2}{q}\right)
 \label{theta:Hermitian-metric}
\end{equation}
descends to $L$.  Its Chern connection and curvature are
\begin{align}
 \nabla&=d+A, & A&=\partial\log h
                  =2\pi i s\,dz-\pi i s^2\,d\tau,
                  \label{theta:chern-connection}\\
 F_\nabla&=\frac\pi q\eta\wedge\overline\eta
                 =-2\pi i\,dr\wedge ds,
 &\eta&=dz-s\,d\tau=dr+\tau ds.
 \label{theta:chern-curvature}
\end{align}
In particular, $\int_{E_\tau}iF_\nabla/(2\pi)=1$.
The vertical covariant operator in the chosen relative frame is
\begin{equation}
 \mathcal D=\nabla_{\partial_z}=\partial_z+2\pi i s.
 \label{theta:vertical-operator}
\end{equation}
Each $\mathcal D^k\vartheta$ transforms by the multiplier
\eqref{theta:multipliers}.
\end{proposition}

\begin{proof}
Put $v=\Im z$.  Under a lattice translation $v$ becomes $v+nq$, while
$|j_{m,n}|^2=e^{2\pi n^2q+4\pi nv}$.  Consequently
$h(\tau,z+m+n\tau)|j_{m,n}(\tau,z)|^2=h(\tau,z)$.
Using $\partial_z v=1/(2i)$ and $\partial_\tau q=1/(2i)$ gives
\eqref{theta:chern-connection}.  Differentiation yields
\[
 dA=2\pi i\,ds\wedge(dz-s\,d\tau).
\]
Since
$ds=(\eta-\overline\eta)/(2iq)$, this is
$(\pi/q)\eta\wedge\overline\eta$.  Finally
$\eta\wedge\overline\eta=-2iq\,dr\wedge ds$.
The fiber integral is one, with orientation $dr\wedge ds$.
Under translation $s$ becomes $s+n$ and
$\partial_z j_{m,n}=-2\pi in j_{m,n}$.  The two extra terms cancel in
$\mathcal D(j_{m,n}a)$ after using the translated $s$.  Iterating proves
the covariant transformation statement.
\end{proof}

\begin{remark}
The covariant derivatives here are smooth sections.  They need not be
holomorphic: $[\partial_{\bar z},\mathcal D]=-\pi/q$.
Also, $\mathcal D^k$ means repeated vertical differentiation in the
translation-invariant frame $\partial_z$ of the relative tangent bundle.
It is not an unspecified full covariant Hessian on the total space.
The Chern connection is distinct from the theta heat connection in
\eqref{theta:automorphy-heat}.
\end{remark}

\begin{theorem}[Exact covariant theta-source formula]
\label{theta:covariant-source}
Let $y=e^{4u}$ and evaluate the theta section at $(\tau,z)=(iy,0)$.
Then
\begin{align}
 \Phi(u)
 &=y^{1/4}\left(y^2\partial_y^2+\frac32y\partial_y\right)
                      \vartheta(0,iy) \label{theta:y-source}\\
 &=y^{1/4}\left(\frac{y^2}{16\pi^2}\partial_z^4
                  +\frac{3y}{8\pi}\partial_z^2\right)
                      \vartheta(z,iy)\big|_{z=0} \label{theta:ordinary-source}\\
 &=\boxed{\displaystyle
   \frac{y^{1/4}}{16\pi^2}
      \bigl(y^2\mathcal D^4\vartheta-3\pi^2\vartheta\bigr)(0,iy).}
                      \label{theta:covariant-source-formula}
\end{align}
The final expression is evaluated in the stated holomorphic frame along
the zero section of the marked family.  It is a source identity, not a
norm identity or a positivity statement.
\end{theorem}

\begin{proof}
For $D_y=y\partial_y$, one has $\partial_u=4D_y$ and
$y^{-1/4}D_y y^{1/4}=D_y+1/4$.  Thus
\[
 \frac1{16}y^{-1/4}(\partial_u^2-1)y^{1/4}
       =D_y^2+\tfrac12D_y
       =y^2\partial_y^2+\tfrac32y\partial_y,
\]
which proves~\eqref{theta:y-source}.  The theta heat equation gives
$\partial_y\vartheta(z,iy)=(4\pi)^{-1}\partial_z^2\vartheta(z,iy)$,
proving~\eqref{theta:ordinary-source}.
At fixed $\tau$, let $a=2\pi is$, so $\mathcal D=\partial_z+a$.
Then $\partial_z a=\pi/q$ and $\partial_z^2a=0$.  At $z=0$, $a=0$,
and the operator product gives
\begin{align*}
 (\mathcal D^2\vartheta)|_0
   &=(\partial_z^2\vartheta)|_0+\frac\pi q\vartheta|_0,\\
 (\mathcal D^4\vartheta)|_0
   &=(\partial_z^4\vartheta)|_0
     +\frac{6\pi}{q}(\partial_z^2\vartheta)|_0
     +\frac{3\pi^2}{q^2}\vartheta|_0.
\end{align*}
Set $q=y$ and substitute the last equality into
\eqref{theta:ordinary-source}.  The second-derivative terms cancel,
leaving~\eqref{theta:covariant-source-formula}.
\end{proof}

The physical reconstruction is now an exact integral along the real
path $\tau=iy$, $y\geq1$, in this marked family:
\begin{align}
 H(t,x)
 &=\frac14\int_1^\infty
    e^{t(\log y)^2/16}y^{1/4}
    \left(y^2\partial_y^2+\frac32y\partial_y\right)\vartheta(0,iy)
       \cos\!\left(\frac x4\log y\right)\frac{dy}{y}
       \label{theta:physical-Mellin}\\
 &=\frac1{64\pi^2}\int_1^\infty
    e^{t(\log y)^2/16}y^{1/4}
    \bigl(y^2\mathcal D^4\vartheta-3\pi^2\vartheta\bigr)(0,iy)
       \cos\!\left(\frac x4\log y\right)\frac{dy}{y}.
       \label{theta:covariant-Mellin}
\end{align}
Both integrals retain all theta indices.  The physical heat multiplier
$e^{t(\log y)^2/16}$ is not the parameter evolution in the Jacobi heat
equation.  Inversion $y\mapsto1/y$ preserves this squared-log multiplier;
that reflection symmetry must not be discarded.  Nevertheless, the
oscillatory integral~\eqref{theta:covariant-Mellin} does not become a
positive Hermitian norm merely because $L$ has positive fiber degree.

\subsection{A nonflat hyperk\"ahler metric with the correct elliptic fibers}
\label{theta:semiflat}

The same marked family has an explicit metric which incorporates its
varying elliptic modulus.  This is a local semi-flat construction in the
sense of~\cite[Section~2]{GrossWilson2000}; we give a direct verification
so that neither its curvature nor its global scope is implicit.

\begin{proposition}[Semi-flat geometry on the marked cover]
\label{theta:semiflat-proposition}
On $\mathbb H\times(\mathbb R/\mathbb Z)^2$ the metric
\begin{equation}
 g=q(dp^2+dq^2)+q^{-1}(dr+p\,ds)^2+q\,ds^2
   =q|d\tau|^2+q^{-1}|\eta|^2
 \label{theta:semiflat-metric}
\end{equation}
is hyperk\"ahler and nonflat.  Its elliptic fibers have modulus $\tau$
and area one.  A parallel hyperk\"ahler triple is
\begin{equation}
 \omega_I=q\,dp\wedge dq+dr\wedge ds,
 \qquad \omega_J+i\omega_K=d\tau\wedge dz.
 \label{theta:hyperkahler-triple}
\end{equation}
For the orientation $\omega_I^2/2>0$, the three forms are self-dual.
The horizontal sectional curvature is $1/(2q^3)$.
\end{proposition}

\begin{proof}
The fiber metric is $q^{-1}|dr+\tau ds|^2$, with area form
$dr\wedge ds$.  Introduce the orthonormal coframe
\[
 e^1=\sqrt q\,dp,\quad e^2=\sqrt q\,dq,\quad
 e^3=q^{-1/2}(dr+p\,ds),\quad e^4=\sqrt q\,ds.
\]
Writing $e^{ab}=e^a\wedge e^b$, direct expansion gives
\begin{equation}
 \omega_I=e^{12}+e^{34},\qquad
 \omega_J=e^{13}-e^{24},\qquad
 \omega_K=e^{14}+e^{23},\qquad
 \omega_a\wedge\omega_b=2\delta_{ab}\operatorname{vol}_g.
 \label{theta:triple-algebra}
\end{equation}
The first form is closed.  Moreover
$d\tau\wedge dz=d\tau\wedge(dr+\tau ds)$, which is closed,
so the other two forms are closed as well.

For a direct K\"ahler verification, $(\tau,z)$ are holomorphic coordinates
on a local lift, and
\[
 \omega_I=\frac i2\bigl(q\,d\tau\wedge d\bar\tau
                              +q^{-1}\eta\wedge\bar\eta\bigr).
\]
The Hermitian coefficient matrix of the metric is
\begin{equation}
 G=\begin{pmatrix}q+s^2/q&-s/q\\-s/q&1/q\end{pmatrix},
 \qquad \det G=1.
 \label{theta:Hermitian-determinant}
\end{equation}
Thus the K\"ahler Ricci form $-i\partial\bar\partial\log\det G$
vanishes.  The induced Chern connection on the canonical bundle has
zero connection form in the frame $d\tau\wedge dz$; since the metric is
K\"ahler, this frame is Levi-Civita parallel.  Its real and imaginary
parts, together with the parallel K\"ahler form, give the parallel
triple~\eqref{theta:hyperkahler-triple}.  The algebra
\eqref{theta:triple-algebra} fixes its normalization and orientation.

Finally, the surfaces with $r,s$ constant are totally geodesic: in these
real coordinates the metric is block diagonal, and its base block is
independent of $r,s$.  Their induced metric is $q(dp^2+dq^2)$.
Its Gaussian curvature is
\begin{equation}
 K_{\mathrm{base}}=-\frac1{2q}(\partial_p^2+\partial_q^2)\log q
                  =\frac1{2q^3}>0.
 \label{theta:nonflat-curvature}
\end{equation}
Total geodesy identifies this with the ambient sectional curvature of
the horizontal plane, proving nonflatness.
\end{proof}

\begin{remark}[The precise global scope]
\label{theta:global-scope}
The formulas descend under the fiber lattice and are defined over the
whole marked upper half-plane.  They do not give a metric on the full
modular quotient.  Under
\[
 \tau'=\frac{a\tau+b}{c\tau+d},\qquad z'=\frac z{c\tau+d},
 \qquad ad-bc=1,
\]
one has
\begin{equation}
 d\tau'\wedge dz'=(c\tau+d)^{-3}d\tau\wedge dz,
 \qquad q'|d\tau'|^2=|c\tau+d|^{-6}q|d\tau|^2.
 \label{theta:modular-obstruction}
\end{equation}
The base metric therefore fails modular invariance as written.
There is also incompleteness: a path with $p,r,s$ constant and
$q\downarrow0$ has finite length $\int_0^{q_0}\sqrt q\,dq$.
No complete K3 metric or extension through singular fibers is asserted.
The global construction in~\cite{GrossWilson2000} introduces local
models and corrections around singular elliptic fibers; that separate
existence theorem does not identify the physical transform with a K3
invariant.  Any modular descent here must supply the missing base-period
and line-bundle data and then control the resulting boundary terms.
\end{remark}

\subsection{An actual abelian anti-self-dual theta connection}
\label{theta:ASD}

The connection in Proposition~\ref{theta:chern} carries the arithmetic
section, but its curvature is not anti-self-dual for
\eqref{theta:semiflat-metric}.  This distinction can be resolved explicitly
on the marked noncompact family, without changing the vertical source
formula.

\begin{proposition}[A Chern connection with anti-self-dual curvature]
\label{theta:ASD-connection}
Let $\omega_b=q\,dp\wedge dq$ and $\omega_f=dr\wedge ds$.
Equip the same holomorphic theta line bundle with the Hermitian metric
\begin{equation}
 \widehat h=\exp\!\left(\frac{2\pi}{3}q^3\right)h.
 \label{theta:ASD-metric}
\end{equation}
Its Chern connection has connection form and curvature
\begin{align}
 \widehat A&=A-\pi i q^2\,d\tau,
                       \label{theta:ASD-form}\\
 F_{\widehat\nabla}
    &=2\pi i(\omega_b-\omega_f),
 & *F_{\widehat\nabla}&=-F_{\widehat\nabla}.
                       \label{theta:ASD-curvature}
\end{align}
It is an abelian Yang--Mills connection, and its vertical covariant
operator remains $\mathcal D$.  Thus
\eqref{theta:covariant-source-formula} is unchanged.
\end{proposition}

\begin{proof}
The extra factor in~\eqref{theta:ASD-metric} depends only on the base,
so it preserves the lattice descent and the holomorphic bundle.
Since $\partial_\tau q=1/(2i)$,
\[
 \partial\!\left(\frac{2\pi}{3}q^3\right)
       =-\pi i q^2\,d\tau,
 \qquad d(-\pi i q^2\,d\tau)=2\pi i q\,dp\wedge dq.
\]
Combine this with~\eqref{theta:chern-curvature}.
In the orthonormal coframe, $\omega_b=e^{12}$ and
$\omega_f=e^{34}$, so $*\omega_b=\omega_f$ and
$*\omega_f=\omega_b$.  This proves anti-self-duality.
The curvature is closed, and hence $d*F_{\widehat\nabla}=0$;
for a line bundle this is the Yang--Mills equation.
The extra connection form vanishes on $\partial_z$, proving the last
assertion.
\end{proof}

\begin{remark}[Scope of the ASD construction]
\label{theta:ASD-limits}
This is a nonflat geometric setting with the actual theta section and
an anti-self-dual line-bundle connection. The unmodified Chern curvature
has self-dual part $F_\nabla^+=-\pi i\omega_I$. The base rescaling
removes that part while preserving the vertical source expression.
It provides neither compactness nor a finite-energy instanton: the
real curvature $F_{\widehat\nabla}/i$ has squared norm $8\pi^2$, while
$\operatorname{vol}_g=q\,dp\,dq\,dr\,ds$ has infinite integral on this
cover. The displayed metric's rapid cusp growth is not an invariant
obstruction; the classification below gives a representative normalized
on the physical path. Neither positive degree nor the Yang--Mills
equation supplies the required signed physical estimate.
\end{remark}

% Prepared for Edward Baker, 15 October 2026.
% Substantial LLM assistance: gpt-6-astra, reasoning effort ultra.
\subsection{Physical restriction and complete signed source identities}
\label{theta:complete-source-bridge}

The ASD construction does not determine the norm of the distinguished
arithmetic section, and its physical restriction is flat.  The useful
data that survive are the transverse jet identities and the complete
signed transform, including its endpoints \cite{ThetaMetric,ThetaBridge}.

\begin{proposition}[Metric freedom and transverse jets]
\label{theta:metric-freedom}
Among base-only changes $h_b=e^{b(p,q)}h$, the ASD metrics are exactly
\begin{equation}
 b=\frac{2\pi}{3}q^3+2\Re f(\tau),\qquad
 f\text{ holomorphic on }\mathbb H.
 \label{theta:ASD-classification}
\end{equation}
They have the same curvature and vertical operator $\mathcal D$.
In particular, $f=-\pi i\tau^3/3$ gives
\begin{equation}
 \widetilde h=e^{2\pi p^2q}h,\qquad
 \widetilde A=A-\pi i(q^2+\tau^2)d\tau,
 \qquad\gamma^*\widetilde h=1,\quad\gamma^*\widetilde A=0,
 \label{theta:path-normalized-ASD}
\end{equation}
where $\gamma(u)=(ie^{4u},0)$.  The natural external product connection
on two physical paths is also flat.

In a marked holomorphic chart set
$X=4q\partial_q|_{p,z,\bar z}$ and
$J_k(u)=(\mathcal D^k\vartheta)(0,iq)$, $q=e^{4u}$.  Compute the
commutators before restriction to the zero section.  Then
\begin{align}
 [\nabla_X^{(b)},\mathcal D^k]
 &=-4k\pi i s\mathcal D^{k-1}
   -\frac{k(k-1)\pi}{q}\mathcal D^{k-2},
 \label{theta:transverse-commutator}\\
 J_k'&=\frac q\pi J_{k+2}-(2k+1)J_k
       -\frac{k(k-1)\pi}{q}J_{k-2}.
 \label{theta:covariant-jet-recurrence}
\end{align}
Terms with zero prefactor are omitted.  The extension $X$ is chart-local;
it is not asserted to be a global horizontal field under the
$\tau$-dependent lattice translations.
\end{proposition}
\begin{proof}
Direct differentiation gives
\[
 F_b=-2\pi i\,dr\wedge ds
       +\frac i2(b_{pp}+b_{qq})\,dp\wedge dq.
\]
Thus ASD is equivalent to $\Delta b=4\pi q$.  The difference from
$2\pi q^3/3$ is harmonic, and the simply connected upper half-plane
supplies its harmonic conjugate.  The stated choice of $f$ gives
\eqref{theta:path-normalized-ASD}.  For comparison,
$\gamma^*\widehat A=d(\pi e^{12u}/3)$; on the external product the
connection is the sum of these two exact one-forms, with zero curvature.

Write $a=2\pi i s$ and $c=\pi/q$.  One has
$[\nabla_X^{(b)},\mathcal D]=-2a$,
$[\mathcal D,a]=c$, and $[\mathcal D,c]=0$; the base-only connection
term commutes with $\mathcal D$.  Iteration gives
\eqref{theta:transverse-commutator}.  The recurrence follows by
differentiating the generating identity
\[
 \sum_{k\ge0}J_k(u)\frac{w^k}{k!}
   =e^{\pi w^2/(2q)}\vartheta(w,iq)
\]
and using the Jacobi heat equation.  The base scalar cancels upon
returning to coefficients in the marked arithmetic frame.
\end{proof}

The rapid cusp growth of $\widehat h$ is therefore a property of that
specified metric, not an obstruction determined by ASD curvature.
Changing the metric while keeping $\vartheta$ fixed changes its norm;
the equivalent holomorphic bundle isometry instead transforms both the
section and its dual observation.  Flat physical restriction does not
remove the second term of \eqref{theta:transverse-commutator}: for
$k\ge2$ it survives at $s=0$.  It also does not make the oscillatory
observation a Hermitian norm.

\begin{proposition}[The complete signed bilocal observations]
\label{theta:bilocal-complete}
Put $c_j(x,u)=u^j\cos(xu+j\pi/2)$ and define the chord
\[
 \mathcal C_j(t;x,y)
   =h_j(t,x)\{h_j(t,y)-h_{j+1}(t,y)\},\qquad j=0,1.
\]
For $x<y$, $h_{j+1}(t,x)=0$, and $h_j$ nonzero on the whole interval,
its nonnegativity is equivalent to an angular increase of at most
$\pi/4$.  Define the symmetric source kernels
\begin{align}
 K_j(x;u,v)
 &=\frac{(uv)^j}{4}\left[(u+v)^2\cos(x(u-v))
       +(-1)^j(u-v)^2\cos(x(u+v))\right],\label{theta:K-complete}\\
 M(x;u,v)
 &=\frac12\left[(1+uv)\cos(x(u-v))
                       +(1-uv)\cos(x(u+v))\right],\label{theta:M-complete}\\
 T_j(x,y;u,v)
 &=\frac12\bigl[c_j(x,u)\{c_j(y,v)-c_{j+1}(y,v)\}
 \notag\\[-2pt]
 &\hspace{35mm}+c_j(x,v)\{c_j(y,u)-c_{j+1}(y,u)\}\bigr].
 \label{theta:T-complete}
\end{align}
With $\mathcal B(u)=e^{8u}J_4(u)-3\pi^2J_0(u)$, all three
observations have the same complete covariant insertion:
\begin{equation}
 \begin{pmatrix}\mathcal L_j(t,x)\\ \Ephys(t,x)\\
                       \mathcal C_j(t;x,y)\end{pmatrix}
 =\frac1{256\pi^4}\int_0^\infty\!\int_0^\infty
 e^{t(u^2+v^2)+u+v}\mathcal B(u)\mathcal B(v)
 \begin{pmatrix}K_j(x;u,v)\\M(x;u,v)\\T_j(x,y;u,v)\end{pmatrix}
 \,du\,dv.
 \label{theta:complete-bilocal-insertion}
\end{equation}
\end{proposition}
\begin{proof}
Write $h_j=\int_0^\infty e^{tu^2}\Phi(u)c_j(x,u)\,du$.
Multiply the defining observations and symmetrize in $u,v$.
The product-to-sum identities give $K_j$ and $M$; the directed chord
gives $T_j$.  Insert \eqref{theta:covariant-source-formula} in both
variables.  Complete-source decay justifies every product and
differentiation locally uniformly in the parameters.  The angular
statement follows from the continuous lift
$\alpha_j(t,y)-\alpha_j(t,x)=\arctan(h_{j+1}(t,y)/h_j(t,y))$,
because $h_j$ has one sign throughout the interval.
\end{proof}

The primitive in \eqref{theta:primitive-definitions} displays every
lower-boundary contribution without splitting the cancelling theta
terms.  Define
\begin{equation}
 I_j(t,x)=\int_0^\infty\psi(u)
       (\partial_u^2-1)\{e^{tu^2}c_j(x,u)\}\,du.
 \label{theta:primitive-I}
\end{equation}
Two integrations by parts, using $\psi'(0)=-1/2$, give
\[
 \int_0^\infty(\psi''-\psi)f\,du
 =\tfrac12 f(0)+\psi(0)f'(0)
      +\int_0^\infty\psi(f''-f)\,du.
\]
Here $f=e^{tu^2}c_j(x,u)$ is even, $f(0)=\delta_{j0}$, and the
upper endpoints vanish by double exponential decay.  Consequently
$8h_j=I_j+\delta_{j0}/2$, and the full signed reconstructions are
\begin{align}
 64\mathcal L_0&=I_1^2-I_0I_2-\tfrac12 I_2,&
 64\mathcal L_1&=I_2^2-I_1I_3,\label{theta:Lag-affine}\\
 64\Ephys&=\tfrac14+I_0+I_0^2+I_1^2,&&\label{theta:E-affine}\\
 64\mathcal C_0(t;x,y)
 &=I_0(t,x)\{I_0(t,y)-I_1(t,y)\}
 \notag\\[-2pt]
 &\quad+\tfrac12 I_0(t,x)
       +\tfrac12\{I_0(t,y)-I_1(t,y)\}+\tfrac14,
 &&\label{theta:C0-affine}\\
 64\mathcal C_1(t;x,y)
 &=I_1(t,x)\{I_1(t,y)-I_2(t,y)\}.&&\label{theta:C1-affine}
\end{align}
Unspecified arguments in \eqref{theta:Lag-affine}--\eqref{theta:E-affine}
are $(t,x)$.  The affine terms and constants record the lower edges and
corner of the bilocal integration by parts; none are disposable.
Every $I_j$ retains the physical heat multiplier inside its differential
operator.  These identities specify complete quantities to estimate,
but assert no sign.  In a chord estimate the comparison energy is
$\Ephys(t,x)$ at the starting point, also on branch $j=1$; neither
endpoint energy nor the auxiliary norm $h_1^2+h_2^2$ may replace it.

% Prepared for Edward Baker, 15 October 2026.
% Substantial LLM assistance; recorded model gpt-6-astra, effort ultra.
\section{Arithmetic discrimination: support, endpoints and reflection}
\label{arith:discrimination}

The signed bridge requires information that distinguishes the actual theta
section from positive weighted square sums.  The following control rejects
support-only moment certificates, even with the first source endpoint and
primitive normalization imposed.  Complete reflection, by contrast, forces
the theta coefficients themselves.  Neither statement proves the remaining
signed estimate.

\subsection{A validated collision within the square-support relaxation}

Put $\lambda_n=\pi n^2$, $q=e^{4u}$, and
\[
 \phi_n(u)=e^u(2\lambda_n^2q^2-3\lambda_nq)e^{-\lambda_nq},
 \qquad H_n(t,x)=\int_0^\infty e^{tu^2}\phi_n(u)\cos(xu)\,du.
\]
For fixed nonnegative weights $a_n$, write
$\Phi_a=\sum_{n\ge1}a_n\phi_n$,
$\psi_a=e^u\sum_{n\ge1}a_ne^{-\lambda_nq}$, and $H_a=\sum a_nH_n$.
For finite support these definitions give entire physical transforms,
$H_{a,t}=-H_{a,xx}$, and $8\Phi_a=\psi_a''-\psi_a$.
The source is positive on $u\ge0$ whenever some weight is positive.
Moreover, for every real polynomial $V$ and every polynomial $Q$ nonnegative
on the square nodes,
\begin{equation}
 \sum_{n\ge1}a_ne^{-\lambda_nq}Q(\lambda_n)V(\lambda_n)^2\ge0
 \qquad(q>0).
 \label{arith:localizers}
\end{equation}
This includes every finite moment matrix and adjacent-square localizing
matrix, by taking $Q(\lambda)=(\lambda-\lambda_m)(\lambda-\lambda_{m+1})$.

\begin{proposition}[Validated weighted-square control]
\label{arith:weighted-control}
There are fixed positive weights on $n=1,2,3$, zero otherwise, and a fixed
positive normalization such that
\[
 \Phi_a'(0)=0,\qquad \psi_a'(0)=\psi_a'''(0)=-\tfrac12,
 \qquad H_a(\tfrac14,x_*)=H_{a,x}(\tfrac14,x_*)=0,
 \quad H_{a,xx}(\tfrac14,x_*)\ne0,
\]
where $44.1303<x_*<44.1304$.  The third source derivative is nonzero, so
this source is not even in $u$.  The control satisfies
\eqref{arith:localizers} at every degree and every $q>0$.
\end{proposition}
\begin{proof}
The construction and validated record are given in
\cite{ThetaSquareControl}.  Define
\[
 P_1(a)=-8a^3+30a^2-15a,\quad d_n=P_1(\lambda_n),\quad
 G_n=e^{\lambda_n}H_n,\quad b_n=-d_1/d_n\quad(n=2,3).
\]
Here $d_1>0>d_2,d_3$, and
$G_\theta=G_1+\theta b_2G_2+(1-\theta)b_3G_3$ has positive weights
for $0<\theta<1$.  Its source derivative at zero is exactly
$d_1+\theta b_2d_2+(1-\theta)b_3d_3=0$.
At $T=1/4$, set $A=G_1+b_3G_3$, $B=b_2G_2-b_3G_3$, and
$D=A'B-AB'$.  On $I=[44.1303,44.1304]$, the validated enclosures include
\begin{align*}
 A&<0<B,&
 D(44.1303)&\in(5.0206,5.0208)10^{-14},\\
 D'&\in(-1.178,-1.121)10^{-9},&
 D(44.1304)&\in(-6.4752,-6.4750)10^{-14}.
\end{align*}
Thus $D$ has exactly one zero $x_*$ in $I$.  At that zero,
$\theta_*=-A(x_*)/B(x_*)\in(0.87916,0.87919)$ gives a double zero, with
\[
 \partial_x^2G_{\theta_*}(T,x_*)\in(-5.95,-5.65)10^{-6}.
\]
The raw weights are
$(a_1,a_2,a_3)=(1,\theta_*b_2e^{3\pi},(1-\theta_*)b_3e^{8\pi})$.
They are fixed after this construction.  Multiplication by
$c=-1/(2\psi_a'(0))>0$ imposes the stated primitive endpoint; then
$\psi_a'''(0)=\psi_a'(0)+8\Phi_a'(0)=-1/2$.
The raw third source derivative lies in $(-40.00,-39.98)$, so positive
scaling does not restore source reflection.

The enclosure calculation uses validated integration on $[0,1]$, together
with the complete tail bound, for $k\le3$,
\[
 \int_1^\infty u^ke^{u^2/4}e^{\lambda_n}\phi_n(u)\,du
 \le\lambda_n\exp\{\lambda_n-4-(\lambda_n/2)e^4\}.
\]
Indeed $k\log u+u^2/4+9u\le(\lambda_n/2)e^{4u}$ for $u\ge1$;
its elementary endpoint constant is also enclosed.  Centered jets and
absolute moments extend the bounds across all of $I$.  Fifty assertions
at 160-bit precision and an independent replay support these finite
claims; no sampled near-zero is substituted for root existence.
\end{proof}

The restriction to finite support is inessential.  Add a sufficiently
small positive multiple of $\sum_{n\ge4}\phi_n$, and solve
$\Phi_a'(0)=0$ for the third weight.  After this elimination, the Jacobian
of $(H_a,H_{a,x})$ with respect to the remaining weight parameter and $x$
is a nonzero multiple of $B(x_*)H_{a,xx}(T,x_*)$.  The implicit function
theorem preserves the collision and endpoint, with every weight positive.
Every finite ordinary and adjacent-square localizing matrix is then
strictly positive definite.  This continuation is qualitative: no numerical
perturbation radius is asserted.  All resulting weights and normalizations
remain independent of physical heat time.

On the control's full stationary predecessor,
$\mathcal L_0/E=-1/(T-t)+O(1)$; its short admissible chords also become
negative.  Thus the relaxed assumptions cannot supply the required uniform
signed estimates.  The control retains periodicity in $z$ and period two in
$\tau$, and obeys Jacobi heat, but fails the full elliptic multiplier.
Its evenness in the spectral variable $x$ does not imply evenness of its
source in $u$.

\subsection{Complete reflection fixes the coefficients}

\begin{theorem}[Weighted-comb reflection rigidity]
\label{arith:reflection-rigidity}
Let $(a_n)_{n\in\mathbb Z}$ be an even complex sequence with
$|a_n|\le C(1+|n|)^M$ for some $C>0$ and finite $M\ge0$.  Define
$F_a(q)=\sum_{n\in\mathbb Z}a_ne^{-\pi n^2q}$ for $q>0$.
If
\begin{equation}
 F_a(q)=q^{-1/2}F_a(1/q)\qquad(q>0),
 \label{arith:reflection}
\end{equation}
then $a_n=a_0$ for every integer $n$.  No coefficient positivity is needed.
\end{theorem}
\begin{proof}
Use $\widehat f(\xi)=\int f(x)e^{-2\pi ix\xi}\,dx$ and the even tempered
comb $\mu=\sum a_n\delta_n$.  For $g_q(x)=e^{-\pi qx^2}$,
$\widehat g_q=q^{-1/2}g_{1/q}$, so \eqref{arith:reflection} gives
$\langle\nu,g_q\rangle=0$, where $\nu=\mu-\widehat\mu$.
Differentiation at $q=1$ in the Schwartz topology annihilates every even
Gaussian moment of $\nu$; odd moments vanish by evenness.
The function
\[
 \mathcal G(z)=\langle\nu,e^{-\pi x^2+zx}\rangle
\]
is entire, since the test functions and their parameter derivatives have
uniform Schwartz seminorm bounds on compact $z$-sets.  Every derivative
at zero vanishes, hence $\mathcal G=0$.  On the imaginary axis this
represents the Fourier transform of $e^{-\pi x^2}\nu$; Fourier injectivity
and local division by the nowhere-zero Gaussian give $\nu=0$.
Thus $\mu=\widehat\mu$.  Integer support makes $\widehat\mu$ period one
as a distribution.  Therefore $\mu$ is period one, and tests supported
near successive integers give $a_{n+1}=a_n$.  This proves the theorem;
\cite{ThetaReflection} records the full scope and defect identity.
\end{proof}

\begin{corollary}[Source reflection and the primitive endpoint]
\label{arith:source-rigidity}
In the class of Theorem~\ref{arith:reflection-rigidity}, take $a_0=1$ and
set, on the entire real $u$ axis,
\[
 P_a(u)=e^uF_a(e^{4u})=e^u+2\psi_a(u),\qquad
 \Phi_a=(P_a''-P_a)/16.
\]
If $\Phi_a(-u)=\Phi_a(u)$ and $\psi_a'(0)=-1/2$, then $a_n=1$ for all $n$.
\end{corollary}
\begin{proof}
Local uniform convergence makes these functions real analytic.
The difference $R(u)=P_a(u)-P_a(-u)$ satisfies
$R''-R=0$, $R(0)=0$, and
$R'(0)=2[1+2\psi_a'(0)]=0$.  Hence $P_a$ is even, which is exactly
\eqref{arith:reflection}; apply the theorem.
\end{proof}

Together with the normalized primitive endpoint, the complete odd source
endpoint hierarchy likewise gives reflection by real analyticity.  Proposition~\ref{arith:weighted-control} retains
only its first condition.  Alternatively, imposing the full elliptic law
on $\sum_n a_ne^{\pi in^2\tau+2\pi inz}$ directly forces
$a_{n+1}=a_n$ by reindexing.  These relations supply coefficient coupling
that support-only positivity omits.  They identify the theta input, but
an additional argument must still transfer that coupling to a favorable
sign of the physical chord or stationary observation.

\subsection{A normalized coefficient recurrence}
\label{core:recurrence-subsection}

The projected elliptic law supplies a convenient arithmetic input for the
adjoint calculation below.  Its normalization is essential: the projected
law alone retains an affine ambiguity in the square coefficients.

\begin{proposition}[The normalized recurrence core]
\label{core:normalized-recurrence}
Let $(a_n)_{n\ge0}$ have polynomial growth, extend it evenly to $\mathbb Z$,
and put
\[
 \beta_n=(4\pi n^2-1)e^{-\pi n^2},\qquad
 \Delta a_n=a_{n-1}-2a_n+a_{n+1}.
\]
Then the conditions
\begin{equation}
 a_0=1,\qquad \Delta a_n=0\quad(n\ge1),\qquad
 \sum_{n\ge1}\beta_na_n=\tfrac12
 \label{core:normalized-input}
\end{equation}
force $a_n=1$ for every $n$.  No positivity assumption on $a_n$ is needed.
Moreover, the recurrence is exactly the vanishing of the complete
one-shift defect after the physical source projection.
\end{proposition}
\begin{proof}
For the last assertion define
\[
 \Theta_a(z,\tau)=\sum_{n\in\mathbb Z}
             a_ne^{\pi i n^2\tau+2\pi inz},\qquad
 (\mathsf TF)(z,\tau)=e^{\pi i\tau+2\pi iz}F(z+\tau,\tau).
\]
At $q=e^{4u}\ge1$, use the complete ordinary source operator
from \eqref{theta:ordinary-source},
\[
 \mathcal S_qF
 =\left.e^u\left(\frac{q^2}{16\pi^2}\partial_z^4
             +\frac{3q}{8\pi}\partial_z^2\right)F(z,iq)\right|_{z=0}.
\]
Reindexing the shifted series and pairing opposite modes gives
\begin{equation}
 \mathcal S_q(\mathsf T\Theta_a-\Theta_a)
       =\tfrac12\sum_{n\ge1}(\Delta a_n)\phi_n(u).
 \label{core:projected-recurrence}
\end{equation}
The zero mode is annihilated.  If the right side vanishes for every
$q\ge1$, its least nonzero coefficient, if any, dominates as $q\to\infty$:
its nonzero quadratic polynomial in $q$ times $e^{-\pi n^2q}$
decays more slowly than every later mode.  Polynomial coefficient growth
controls the remaining tail.  Thus every $\Delta a_n$ vanishes;
the converse is immediate.  This proves the stated equivalence.

The recurrence and $a_0=1$ give $a_n=1+bn$.  The complete theta primitive
has $\sum_{n\ge1}\beta_n=1/2$, so the endpoint in
\eqref{core:normalized-input} implies
\[
 b\sum_{n\ge1}n\beta_n=0.
\]
The last sum is strictly positive and finite, hence $b=0$.
The complete projected calculation and its primitive endpoints are
recorded in \cite{ThetaEllipticCore}.
\end{proof}

The endpoint in \eqref{core:normalized-input} is
$-\psi_a'(0)=1/2$.  It restores information lost by source projection;
it does not turn coefficient rigidity into a signed transform estimate.
In particular, the ramp $a_n=1+bn$ passes the projected recurrence
before that normalization is imposed.

\begin{remark}[Scope of finite endpoint restrictions]
\label{core:finite-endpoint-control}
Positive weights on four square modes can satisfy
$\Phi_a'(0)=\Phi_a'''(0)=0$ and the normalized primitive endpoint
while their physical field has a nondegenerate positive-time double zero
\cite{ThetaTwoEndpoint}.  Thus adding that second source endpoint
does not replace the complete recurrence or reflection hypotheses.
This control does not assert failure of every finite endpoint hierarchy.
\end{remark}

\subsection{Coupled addition and the complete starting factor}
\label{core:coupled-subsection}

The arithmetic product must retain independent source points.  Write
$\Theta=\Theta_1$, $q=e^{4u}$, $r=e^{4v}$, with $u,v\ge0$, and set
$p=n+m$, $k=n-m$, $X=z+w$, $Y=z-w$.  The inverse lattice map is
$n=(p+k)/2$, $m=(p-k)/2$, with $p\equiv k\pmod2$.  Hence
\begin{equation}
 \Theta(z,iq)\Theta(w,ir)
 =\sum_{p\equiv k\ (2)}
       e^{-\pi Q_{q,r}(p,k)+\pi i(pX+kY)},
 \quad
 Q_{q,r}=\frac{q+r}{4}(p^2+k^2)+\frac{q-r}{2}pk .
 \label{core:coupled-addition}
\end{equation}
The quadratic form has eigenvalues $q/2,r/2>0$, so these series and
their fixed jets converge normally on compact parameter sets.

For an explicit parity form, put $A=q+r$, $B=q-r$ and
\[
 \vartheta_\epsilon(Z,iA)=\sum_{a\in\mathbb Z}
 e^{-\pi A(a+\epsilon/2)^2+2\pi i(a+\epsilon/2)Z},
 \qquad \epsilon=0,1.
\]
Then the same identity is
\begin{equation}
 \Theta(z,iq)\Theta(w,ir)
 =\exp\!\left(\frac{B}{2\pi}\partial_X\partial_Y\right)
  \sum_{\epsilon=0}^1
       \vartheta_\epsilon(X,iA)\vartheta_\epsilon(Y,iA).
 \label{core:mixed-parity}
\end{equation}
The derivative expansion is normally convergent by the Gaussian margin
$A-|B|=2\min(q,r)>0$.  Its mixed derivative cannot be discarded:
$q=r$ represents only the diagonal $u=v$ of the physical double integral.

Applying the two complete source operators to
\eqref{core:coupled-addition} gives
\begin{equation}
 \Phi(u)\Phi(v)=\sum_{p\equiv k\ (2)}M_{p,k}(u,v),\qquad
 M_{p,k}=\tfrac14
 \phi_{|(p+k)/2|}(u)\phi_{|(p-k)/2|}(v),\quad \phi_0=0.
 \label{core:coupled-source}
\end{equation}
This follows because each nonzero integer mode contributes
$\phi_{|n|}/2$.  The lines $p=k$ and $p=-k$ are annihilated by the
source operators.  Formula \eqref{core:coupled-source} retains positive
source coefficients, but the physical transform still inserts its
heat weights and oscillatory phases \cite{ThetaCoupled}.

\begin{proposition}[Endpoint cutoff in the coupled lattice]
\label{core:complete-start-strip}
Let $h_{j,n}=\partial_x^jH_n$, $W=h_j(t,x)$ and
$V_N=\sum_{m=1}^Nh_{j,m}(t,y)$.  The product $WV_N$, with its complete
starting factor, uses precisely the coupled lattice strip
\begin{equation}
 p\equiv k\pmod2,\qquad |p-k|\le2N,
 \label{core:endpoint-strip}
\end{equation}
with no restriction on $p+k$.  More generally, bounded endpoint weights
$d_m$ give the mask $d_{|(p-k)/2|}$, independent of the starting index.
\end{proposition}
\begin{proof}
The endpoint index is $m=(p-k)/2$, while
$n=(p+k)/2$ is unrestricted.  Pairing the signs of the original integer
indices gives the positive-index product
$\sum_{n\ge1}\sum_{m=1}^Nh_{j,n}(t,x)h_{j,m}(t,y)$.
The same argument gives the weighted mask; its zero-index value is
irrelevant because $\phi_0=0$.  Absolute source bounds justify the
exchanges \cite{ThetaParityStrip}.
\end{proof}

\begin{remark}[Constraints on a signed use of the product]
\label{core:coupled-scope}
The source addition matrix is not a positive Gram kernel: its principal
block on $\{0,2\}$ has zero diagonal and strictly positive off-diagonal
entries.  The factor-preserving estimates used here recombine arithmetic parity
sectors and the physical exchange sectors of a complete-start/cutoff-endpoint
product before charging the error to the starting energy
\cite{ThetaCoupled,ThetaParityStrip,ThetaReflectedControl}.
Full source reflection and generic source/Fourier Gram positivity
also do not supply the missing sign: an explicit positive even source
retains every odd primitive endpoint and nevertheless collides
\cite{ThetaReflectedControl}.  That source has no nonzero
polynomial-growth integer-square expansion, so it does not contradict
Theorem~\ref{arith:reflection-rigidity}.  The coupled identities fix
the arithmetic bookkeeping; the following estimates must still control
its complete signed physical observation.
\end{remark}

% Prepared for Edward Baker with substantial LLM assistance.
% Integration draft; mathematical scope and internal review in \cite{ThetaAdjoint}.
\section{Quantitative recurrence stability and endpoint currents}
\label{adj:section}

The coefficient recurrence also gives quantitative information.
We prove stability for weighted perturbations and a finite endpoint
approximation whose error retains the complete starting amplitude.
Neither result supplies a theta sign; the unresolved lower boundary
is identified explicitly.

\subsection{Physical mode bounds and the imaginary anchor}

Write $\Phi=\sum_{n\ge1}\phi_n$, where
\[
 \phi_n(u)=e^u(2\pi^2n^4e^{8u}-3\pi n^2e^{4u})
                  e^{-\pi n^2e^{4u}},
 \quad
 H_n(t,x)=\int_0^\infty e^{tu^2}\phi_n(u)\cos(xu)\,du,
\]
and put $h_{j,n}=\partial_x^jH_n$.  Define
\begin{equation}
 \beta_n=(4\pi n^2-1)e^{-\pi n^2},\qquad
 \sum_{n\ge1}\beta_n=\tfrac12,\qquad
 m_\beta=\sum_{n\ge1}n\beta_n\ge\tfrac12.
 \label{adj:beta}
\end{equation}
The sum is the complete primitive endpoint $-\psi'(0)$.

\begin{lemma}[Physical bound and endpoint evaluation]
\label{adj:physical-bound}
For every integer $j\ge0$ and finite real $T_{\max}$, the constants
\[
 M_j(T_{\max})=\sup_{u\ge0}u^j
 e^{T_{\max}u^2+9u-(\pi/2)(e^{4u}-1)},\qquad
 C_j(T_{\max})=\frac{\pi M_j(T_{\max})}{4\pi-1}
\]
are finite, and
\begin{equation}
 |h_{j,n}(t,x)|\le C_j(T_{\max})\beta_n
 \qquad(t\le T_{\max},\ x\in\R).
 \label{adj:mode-bound}
\end{equation}
Moreover $\beta_n=8H_n(0,i)$.  Thus, for
$\sum\beta_n|a_n|<\infty$, the mode sum $H_a=\sum a_nH_n$
and its real spatial jets converge absolutely, and
$\sum\beta_na_n=1/2$ is exactly $H_a(0,i)=1/16$.
\end{lemma}
\begin{proof}
Positivity of $\phi_n$ and its explicit formula give
\[
 \frac{\phi_n(u)}{\beta_n}\le
 \frac{2\pi^2}{4\pi-1}n^2
 e^{9u-\pi n^2(e^{4u}-1)}.
\]
Use $e^{4u}-1\ge4u$ and $n^2-1/2\ge n^2/2$,
then integrate $n^2e^{-2\pi n^2u}$, to obtain
\eqref{adj:mode-bound}.  For
$p_n(u)=e^ue^{-\pi n^2e^{4u}}$,
$p_n''-p_n=8\phi_n$ and $p_n'(0)=-\beta_n$.
Integration against $\cosh u$ has zero upper boundary and gives
$8\int_0^\infty\phi_n(u)\cosh u\,du=\beta_n$.
Absolute summability justifies the stated mode sums and evaluation.
\end{proof}

The weighted class in this lemma need not have a smooth source density
at zero.  Its observations are defined by the convergent mode sums;
polynomial growth of the coefficients is not required.

\subsection{An exact normalized defect identity}

For $a_0=1$ and $\sum\beta_n|a_n|<\infty$, set
\[
 d_k=a_{k-1}-2a_k+a_{k+1},\qquad
 e_a=\sum_{n\ge1}\beta_na_n-\tfrac12.
\]
For a real or complex test sequence with $|f_n|\le C\beta_n$, define
\[
 F_a=\sum a_nf_n,\quad F_\theta=\sum f_n,\quad m_f=\sum nf_n,
 \qquad G_f(k)=\sum_{n>k}(n-k)f_n\quad(k\ge0).
\]
Use $G_\beta$ for the same construction with $f=\beta$, so
$G_\beta(0)=m_\beta$.

\begin{theorem}[Normalized recurrence stability]
\label{adj:defect-stability}
Put
\[
 \eta_f=\frac{m_f}{m_\beta},\qquad
 D_k(f)=G_f(k)-\eta_fG_\beta(k).
\]
Then all the following series converge absolutely:
\begin{align}
 F_a-F_\theta&=\eta_fe_a+\sum_{k\ge1}D_k(f)d_k,
 \label{adj:defect-identity}\\
 |F_a-F_\theta|&\le
 C\left(|e_a|+2\sum_{k\ge1}G_\beta(k)|d_k|\right).
 \label{adj:defect-estimate}
\end{align}
The coefficients satisfy
$|\eta_f|\le C$ and
$|D_k(f)|\le2CG_\beta(k)\le8C\beta_{k+1}$.
Their exact lower traces are
\begin{equation}
 \Delta D_k=f_k-\eta_f\beta_k\quad(k\ge1),\qquad
 D_0=0,\qquad D_1=-F_\theta+\eta_f/2.
 \label{adj:dual-traces}
\end{equation}
\end{theorem}
\begin{proof}
Writing $b=a_1-1$, the finite-index recurrence solution is
$a_n=1+nb+\sum_{k<n}(n-k)d_k$.
Since $\beta_{n+1}/\beta_n<1/2$,
$G_\beta(k)\le4\beta_{k+1}$.
Expanding $d_k$ into its three neighboring coefficients shows
$\sum G_\beta(k)|d_k|<\infty$ under the stated weighted assumption.
The recurrence may therefore be summed against both $\beta_n$ and
$f_n$, giving
\[
 e_a=bm_\beta+\sum G_\beta(k)d_k,\qquad
 F_a-F_\theta=bm_f+\sum G_f(k)d_k.
\]
Eliminating $b$ proves \eqref{adj:defect-identity}.
The bounds follow from $|G_f(k)|\le CG_\beta(k)$.
Finally $\Delta G_f(k)=f_k$,
$G_f(0)=m_f$, and $G_f(1)=m_f-F_\theta$, proving
\eqref{adj:dual-traces}.
\end{proof}

For physical jets the estimate is uniform in real $x$ and
$t\le T_{\max}$.  Its coefficients use no division by $H$, $E$,
or stationarity.  Gaussian decay of the dual weights does not,
without a uniform defect-tail bound, imply a uniform Gaussian
remainder after truncating the sum.
At theta, $d=e_a=0$; the identity then retains the unknown
observation in $D_1$, rather than estimating its sign.

\subsection{A finite positive hierarchy and its sharp rate}

\begin{proposition}[Finite recurrence stability]
\label{adj:finite-hierarchy}
Suppose $a_0=1$, $a_n\ge0$, $\sum\beta_na_n=1/2$, and
$d_k=0$ for $1\le k\le N$.  Put $M=N+1$ and
\[
 R^\beta_M=\sum_{n>M}\beta_n,\qquad
 m_M=\sum_{n=1}^M n\beta_n.
\]
Then
\begin{equation}
 \sum_{n\ge1}\beta_n|a_n-1|
 \le2R^\beta_M+\frac{2m_M}{M},\qquad
 |F_a-F_\theta|\le2C\left(R^\beta_M+\frac{m_M}{M}\right).
 \label{adj:finite-bound}
\end{equation}
The first bound is the exact supremum over this class.
The order $M^{-1}$ also occurs for the physical observation $H_a(0,0)$.
If $a_1=1$ is additionally imposed, both bounds improve respectively
to $2R^\beta_M$ and $2CR^\beta_M$, where
$R^\beta_M=O((M+1)^2e^{-\pi(M+1)^2})$.
\end{proposition}
\begin{proof}
The recurrence gives $a_n=1+bn$ for $0\le n\le M$.
Positivity and normalization give
$-1/M\le b\le R^\beta_M/m_M$ and tail mass
$\sum_{n>M}\beta_na_n=R^\beta_M-bm_M$.
Hence the weighted distance is at most
$2R^\beta_M+(|b|-b)m_M$, proving \eqref{adj:finite-bound}.
If $a_1=1$, then $b=0$, giving the improved estimate.

For sharpness, choose $K>M$, set $a_n=(1-n/M)_+$ except at $K$,
and set $a_K=(R^\beta_M+m_M/M)/\beta_K$.
This finite nonnegative sequence has every required constraint.
For sufficiently large $K$, its weighted distance is
$2R^\beta_M+2m_M/M-2\beta_K$.
For the physical assertion put $\omega_n=\int_0^\infty p_n(u)\,du$.
Integration by parts and $e^{4u}\ge1+4u$ give
\[
 H_n(0,0)=\frac{\beta_n-\omega_n}{8},\qquad
 0<\frac{\omega_n}{\beta_n}
 \le(4\pi n^2-1)^{-2}.
\]
The ramp-and-spike deviation consequently tends, as $K\to\infty$, to
\[
 \frac{\sum_{n\le M}n\omega_n+
                     M\sum_{n>M}\omega_n}{8M}
 \sim\frac1{8M}\sum_{n\ge1}n\omega_n>0.
\]
For each $M$, sufficiently large finite $K$ realizes at least half
this positive limit.  Finally the Gaussian tail bound for
$R^\beta_M$ follows directly from the explicit $\beta_n$.
\end{proof}

These controls establish a stability rate, not a collision.
Even the anchored Gaussian absolute error is not automatically
an $o(\sqrt E)$ error near an unknown event.

\subsection{A finite current preserving the starting amplitude}

For a sequence $f$ with $\sum_{n\ge1}n|f_n|<\infty$, adjoin $f_0=0$
and define
\[
 (\mathsf Pf)_0=0,\qquad
 (\mathsf Pf)_n=\tfrac12(f_{n-1}+f_{n+1})\quad(n\ge1).
\]
The zeroth coordinate is reset to zero at every iteration.
For an integer depth $L\ge1$, put
$J^{[L]}=-\frac12\sum_{k=0}^{L-1}\mathsf P^kf$ and $J^{[L]}_0=0$.

\begin{proposition}[Absorption current and remainder]
\label{adj:absorption-current}
Let $c_L=2^{-L}\binom L{\lfloor L/2\rfloor}$ and set
\[
 p_L(k)=2^{-L}\binom L{(L+k)/2},\qquad
 s_L(m)=\sum_{k=1-m}^{m}p_L(k),\qquad q_L(m)=1-s_L(m),
\]
with inadmissible binomial coefficients interpreted as zero.
Then $q_L(m)=0$ for $m>L$, $0\le q_L,s_L\le1$, and
\begin{align}
 \Delta J^{[L]}&=f-\mathsf P^Lf,\qquad
 \sum f_n=B_L(f)+R_L(f), \label{adj:current-identity}\\
 B_L(f)&=-J^{[L]}_1=\sum_{m=1}^Lq_L(m)f_m,\qquad
 R_L(f)=\sum_{m\ge1}s_L(m)f_m, \label{adj:current-weights}\\
 |R_L(f)|&\le c_L\sum_{m\ge1}m|f_m|
 \le\sqrt{\frac2{L+1}}\sum_{m\ge1}m|f_m|.
 \label{adj:current-error}
\end{align}
The first moment $\sum n(\mathsf P^Lf)_n=\sum nf_n$ is conserved.
\end{proposition}
\begin{proof}
On sequences with zero zeroth coordinate, $\Delta=2(\mathsf P-I)$,
so the current identity telescopes.
Reflection at zero gives the nonnegative kernel
$p_L(n-m)-p_L(n+m)$ for $\mathsf P^L$.
Summing over $n\ge1$ gives $s_L(m)$.
Its defining interval contains $m$ admissible parity indices, so
$s_L(m)\le\min(1,mc_L)$.
The elementary central-binomial bound gives
$c_L\le\sqrt{2/(L+1)}$.
Finite iterates preserve weighted absolute summability.
Thus the upper boundaries in
$\sum_{n=1}^Q\Delta J_n=J_0-J_1+J_{Q+1}-J_Q$
and
$\sum_{n=1}^Q n\Delta J_n=J_0-(Q+1)J_Q+QJ_{Q+1}$
vanish as $Q\to\infty$, proving the remaining assertions.
\end{proof}

Apply this construction only at the endpoint, defining
\[
 W^{[L]}(t,y)=\sum_{m=1}^Lq_L(m)h_{j,m}(t,y),\qquad W=h_j.
\]
Lemma~\ref{adj:physical-bound} and
Proposition~\ref{adj:absorption-current} give
$|W-W^{[L]}|\le c_LC_j(T_{\max})m_\beta$ uniformly for real $y$
and $t\le T_{\max}$.
For $j=0,1$ and $\delta>0$, define
\[
 \mathcal A^{[L]}_{j,\delta}
 =W(t,x)^2-e^{-\delta}W(t,x)W^{[L]}(t,x+\delta),
 \quad
 \mathcal A_{j,\delta}=W(t,x)^2-e^{-\delta}W(t,x)W(t,x+\delta).
\]
Because the starting field remains complete,
\begin{equation}
 |\mathcal A_{j,\delta}-\mathcal A^{[L]}_{j,\delta}|
 \le e^{-\delta}c_LC_j(T_{\max})m_\beta\sqrt{E(t,x)}.
 \label{adj:factor-preserving-current}
\end{equation}
This identity requires no stationarity assumption.
Any independent signed-core bound must use the complete condition
$h_{j+1}(t,x)=0$ and its required interval domain.
The coefficient of $\sqrt E$ in this error tends to zero at fixed
$\delta$; it is also smaller than every fixed power of
$\delta_L=1/\log(L+2)$.
The finite primitive endpoint is
$\beta^{[L]}=\sum_{m=1}^Lq_L(m)\beta_m<1/2$, with exact deficit
$R_L(\beta)$, and cannot be replaced by $1/2$.

\begin{remark}[Scope of the negative results]
\label{adj:limitations}
The current gives an approximation, not its sign.
For a fixed residual with $\sum n|r_n|<\infty$, cancelling its
boundary-plus-normalization
term with $\eta_L=-B_L(r)/R_L(\beta)$ leaves exactly
$R_L(r-\eta_L\beta)=\sum r_n$.
The unknown observation has moved to the remainder.
A separate finite-range result rules out the specified exact pointwise
endpoint reconstruction by finitely many neighboring square modes,
polynomial source multipliers and finitely many spectral frequencies.
Its precise hypotheses require a nonzero diagonal target mode beyond
the current's range and exclude an independent source-divergence term.
These scoped obstructions, proved in \cite{ThetaAdjoint}, do not exclude
integrated identities, source-divergence corrections with complete
boundary terms, or one-sided bilocal estimates.
\end{remark}


\subsection{Joint source currents and their remaining boundary}
\label{adj:joint-source}
The source divergence can be retained explicitly. Write
$\lambda_n=\pi n^2$, $q=e^{4u}$,
$p_n=e^{u-\lambda_nq}$, and
$\phi_n=p_n(2\lambda_n^2q^2-3\lambda_nq)$.
For a twice differentiable weight $w_n$, define
\begin{equation}
 r_n=\phi_nw_n,\qquad
 K_n=\frac{p_n'w_n-p_nw_n'}8,\qquad
 g_n=\frac{p_n(w_n''-w_n)}8.
 \label{adj:joint-source-reduction}
\end{equation}
Since $p_n''-p_n=8\phi_n$, direct differentiation gives
$r_n=\partial_uK_n+g_n$.
Apply Proposition~\ref{adj:absorption-current} to $g(u)$ at each
source point, writing
$J^{[L]}=-\frac12\sum_{k=0}^{L-1}\mathsf P^kg$ and
$\varrho^{[L]}=\mathsf P^Lg$. Then
\begin{equation}
 r_n=\Delta J_n^{[L]}+\partial_uK_n+\varrho_n^{[L]}.
 \label{adj:joint-identity}
\end{equation}
Every finite-rectangle boundary must be retained before taking limits;
their exact expression is \eqref{obs:rectangular-trace} below.

For an endpoint cutoff $V_N=\sum_{n\le N}h_{j,n}(t,y)$ with
$y=x+\delta$, and independently prescribed $A,B,C$, take
\begin{align*}
 w_n(u)=e^{tu^2}\bigl[&\mathbf1_{n\le N}u^j\cos(yu+j\pi/2)
       -A\cos(xu)+Bu\sin(xu)\\
       &-Cu^{j+1}\cos(xu+(j+1)\pi/2)\bigr].
\end{align*}
All weights are even, so $w_n'(0)=0$ and
$w_n(0)=\mathbf1_{j=0}\mathbf1_{n\le N}-A$.
Let $G_n=\int_0^\infty g_n\,du$ and
$\mathcal M(g)=\sum_{n\ge1}n\int_0^\infty|g_n|\,du$.
The complete identity is
\begin{align}
 V_N&=AH+BH_x+Ch_{j+1}+B_{\rm src}+B_L(G)+R_L(G),
 \label{adj:joint-physical}\\
 B_{\rm src}
 &=\frac{\mathbf1_{j=0}}8\sum_{n\le N}\beta_n-\frac A{16},
 \qquad |R_L(G)|\le c_L\mathcal M(g).
 \label{adj:joint-source-boundary}
\end{align}
Indeed $p_n(0)=e^{-\lambda_n}$, $p_n'(0)=-\beta_n$, and
$\sum\beta_n=1/2$, which give the source lower boundary shown.
The finite sum in \eqref{adj:joint-source-boundary} cannot be
replaced by $1/2$ without its tail.
For bounded $A,B,C,x,\delta$ and $t$ in a compact real interval
$[t_-,T_*]$, a source majorant has the form
\[
 |g_n(u)|\le C_*(1+u)^d
 e^{T_*u^2+u-(\pi/2)e^{4u}}e^{-\pi n^2/2}.
\]
It is summable with the index weight $n$ and integrable in $u$.
Similar majorants kill the upper boundaries at fixed $L$ and justify
\eqref{adj:joint-physical} locally uniformly. A lower bound on $t$
is included because $w_n''$ contains powers of $t$.
Multiplication by the complete $|W|\le\sqrt E$ turns the remainder
into a mixed error at most $c_L\mathcal M(g)\sqrt E$.

The prescribed choice
$A=2\mathbf1_{j=0}\sum_{n\le N}\beta_n$ cancels $B_{\rm src}$,
but does not cancel the total trace
\begin{equation}
 B_\partial^{[L]}=B_{\rm src}+B_L(G),\qquad
 B_\partial^{[L]}+R_L(G)=V_N-AH-BH_x-Ch_{j+1}.
 \label{adj:joint-total-residual}
\end{equation}
At fixed parameters the trace converges to the same unresolved physical
residual as $L\to\infty$. Boundedness of a signed coefficient norm
alone does not justify the source majorant when $A$ or $B$ diverges;
any such extension needs a separate remainder estimate.

\begin{proposition}[The polynomial source obstruction]
\label{adj:polynomial-source-obstruction}
For real $t,\omega$, $\lambda>0$, and nonzero polynomial $a(u)$,
no polynomial $P(u,q)$ satisfies
\begin{equation}
 \mathcal L_{\lambda,t,\omega}P
 =(2\lambda^2q^2-3\lambda q)a(u),\qquad
 \mathcal L_{\lambda,t,\omega}
 =\partial_u+4q\partial_q+1+2tu+i\omega-4\lambda q.
 \label{adj:polynomial-source-equation}
\end{equation}
\end{proposition}
\begin{proof}
The sole degree-raising term in $q$ is $-4\lambda qP$, so a
solution must have $q$-degree one. Matching its positive powers gives
\[
 P=-\frac\lambda2a q+\frac{(1-2tu-i\omega)a-a'}8.
\]
The constant coefficient then forces $(D^2-1)a=0$, where
$D=\partial_u+2tu+i\omega$.
For $t\ne0$ its highest $u$-degree term is $4t^2u^2a$;
for $t=0$ its leading coefficient is $-(1+\omega^2)$ times that
of $a$. Neither can vanish.
\end{proof}

This also extends the finite-range obstruction to polynomial source
divergences, with the precise pointwise class specified in
\cite{ThetaJointCurrent}. Both currents there are finite sums of
$e^{u+tu^2-\pi m^2q+i\omega u}$ times polynomials in $(u,q)$,
with real frequencies in a fixed finite set. Their mode ranges are
$|m-n|\le L_J$ and $|m-n|\le L_K$, respectively.
Separating square modes and frequencies and taking the discrete first
moment for a mode $m>L_J$ yields
\[
 m(2\lambda_m^2q^2-3\lambda_mq)a_{m,\omega}
 =\mathcal L_{\lambda_m,t,\omega}
       \left(\sum_{n\ge1}nQ_{n,m,\omega}\right),
\]
because its zeroth index-boundary coefficient is absent.
The finite sum is polynomial, so the proposition forces
$a_{m,\omega}=0$. For an endpoint $x+\delta$ with $x\ge0$,
$\delta>0$ and $N>L_J$, its nonzero frequency at mode
$m=L_J+1$ is distinct from all starting frequencies, a contradiction.
Finite square-mode separation follows by dividing by the least square
exponential: a nonzero polynomial-amplitude trigonometric polynomial
cannot decay super-exponentially, as its leading trigonometric part
has positive limiting mean square. The argument does not exclude
nonpolynomial primitives, nonlocal currents, identities after source
integration, or independently proved one-sided estimates.

\section{Arithmetic boundary traces and signed observation}
\label{obs:boundary-geometry}

The recurrence experiment exposes a boundary obstruction that the
physical geometry can identify more precisely.
The conclusions below concern exact sequence spaces, finite linear
relaxations, and conditional physical estimates; they do not establish
a signed theta inequality. Details are recorded in
\cite{ThetaBoundaryGeometry,ThetaSignedQuotient,ThetaScalarGeometry,
ThetaJointCurrent}.

Let $\mathcal S_+$ and $\mathcal S_0$ denote the real rapidly
decreasing sequences on $n\ge1$ and $n\ge0$, respectively, and set
\[
 (\Delta J)_n=J_{n-1}-2J_n+J_{n+1},\qquad
 \mathcal N=\{J\in\mathcal S_0:J_0=J_1\}.
\]
\begin{proposition}[The recurrence boundary class]
\label{obs:trace-cokernel}
The map $\Delta:\mathcal N\to\mathcal S_+$ is injective and
\begin{equation}
 \operatorname{range}(\Delta|_{\mathcal N})
 =\left\{r\in\mathcal S_+:\sum_{n\ge1}r_n=0\right\}.
 \label{obs:trace-range}
\end{equation}
Its algebraic cokernel is one-dimensional, detected by total sum.
\end{proposition}
\begin{proof}
Telescoping gives $\sum_{n\ge1}\Delta J_n=J_0-J_1$.
For any $r\in\mathcal S_+$, the current
\begin{equation}
 J_n=\sum_{m>n}(m-n)r_m
 \label{obs:tail-current}
\end{equation}
is rapidly decreasing, satisfies $\Delta J=r$, and has
$J_0-J_1=\sum r_n$.
Rapid decrease follows from arbitrarily high weighted tail bounds.
A homogeneous solution is affine in $n$, hence is zero if rapidly
decreasing.  This proves the range assertion and injectivity.
\end{proof}

The primitive endpoint sequence
$\beta_n=(4\pi n^2-1)e^{-\pi n^2}$ has $\sum\beta_n=1/2$.
Its class therefore spans this cokernel: uniquely,
\[
 r=\Delta J+\eta\beta,\qquad J\in\mathcal N,
 \qquad \eta=2\sum r_n.
\]
Choosing $\eta$ from the unknown physical residual does not estimate
that residual.  An effective arithmetic construction must control
the lower trace independently, while retaining every finite-cutoff
boundary.  Parameter-dependent estimates additionally require locally
uniform weighted bounds; a pointwise sequence identity does not
supply them.


\subsection{One invariant trace for the two arithmetic boundaries}
\label{obs:invariant-trace-section}
Suppose $r_n=\Delta J_n+\partial_uK_n+e_n$, where $J$ is defined
for $n\ge0$, $K$ for $n\ge1$, and $u\ge0$.
On a finite rectangle its complete trace is
\begin{equation}
 \begin{split}
 \mathcal B_{M,U}(J,K)
 ={}&\int_0^U(J_0-J_1+J_{M+1}-J_M)\,du\\
    &+\sum_{n=1}^M\{K_n(U)-K_n(0)\}
 =\sum_{n=1}^M\int_0^U(r_n-e_n)\,du.
 \end{split}
 \label{obs:rectangular-trace}
\end{equation}
For every sufficiently smooth index-source field $L$, the change
\begin{equation}
 \widetilde J=J+\partial_uL,\qquad
 \widetilde K=K-\Delta L
 \label{obs:current-gauge}
\end{equation}
preserves the residual $e$, the joint divergence, and the whole trace.
Indeed $\Delta\partial_uL=\partial_u\Delta L$, while
$\sum_{n=1}^M\Delta L_n=L_0-L_1+L_{M+1}-L_M$ makes the
changes on the four sides of the rectangle cancel exactly.
Keeping only one lower boundary would lose this invariance.

For the infinite limit, assume every weighted index supremum of
$J(u)$ and $r(u)-e(u)$ has an integrable source majorant,
$k_n=K_n(0)$ is rapidly decreasing, and
$\sum_n|K_n(U)|\to0$ as $U\to\infty$.
These hypotheses, and locally uniform versions for parameter estimates,
justify the following sums, integrals, and vanished upper boundaries.
Write
\[
 a_n=\int_0^\infty(r_n-e_n)\,du,\qquad
 \overline J_n=\int_0^\infty J_n\,du,\qquad
 (\mathcal Gk)_n=\sum_{m>n}(m-n)k_m.
\]
By \eqref{obs:tail-current}, $\Delta\mathcal Gk=k$.
The effective current $\widehat J=\overline J-\mathcal Gk$
satisfies $\Delta\widehat J=a$, whence uniqueness of the rapidly
decreasing inverse gives $\widehat J=\mathcal Ga$. In particular,
\begin{equation}
 \boxed{\widehat J_0-\widehat J_1
 =\int_0^\infty(J_0-J_1)\,du-\sum_{n\ge1}K_n(0)
 =\sum_{n\ge1}a_n.}
 \label{obs:invariant-total-trace}
\end{equation}
Both the effective current and its trace are therefore invariant under
gauges retaining these hypotheses and the same residual.
For example, a smooth compactly supported collar $\chi$ with
$\chi(0)=1$ and $L_n(u)=\chi(u)(\mathcal Gk)_n$ makes
$\widetilde K_n(0)=0$. Since $\int_0^\infty\chi'=-1$, it
transfers $-\sum k_n$ to the integrated lower index flux, leaving
\eqref{obs:invariant-total-trace} unchanged.
This gauge can be nonlocal in the index and need not preserve a finite
polynomial ansatz or improve its norms. It consolidates the obstruction
into one scalar trace; it does not estimate that trace.

\subsection{Absolute and signed physical observation quotients}

\begin{proposition}[Bounded descent to the physical observation quotient]
\label{obs:quotient-norm}
Let $V$ be a finite-dimensional real inner-product space describing a
homogeneous linear relaxation of arithmetic and stationarity
constraints.  Let $\mathsf J:V\to\R^2$ be its physical first-jet map
and $\ell:V\to\R$ an endpoint functional.
There exists a finite $C$ such that
\begin{equation}
 |\ell(c)|\le C\|\mathsf Jc\|\quad(c\in V)
 \label{obs:quotient-bound}
\end{equation}
if and only if $\ker\mathsf J\subseteq\ker\ell$.
In orthonormal coordinates the least constant is
$C_{\min}=\|\ell\mathsf J^\dagger\|$, where
$\mathsf J^\dagger$ is the Moore--Penrose inverse.
\end{proposition}
\begin{proof}
Necessity follows on kernel vectors.  Under kernel inclusion,
$b(\mathsf Jc)=\ell(c)$ is a well-defined covector on
$\operatorname{im}\mathsf J$, with least bound its operator norm.
The row $\ell\mathsf J^\dagger$ represents that covector and
vanishes on the orthogonal complement of the image.
\end{proof}

Affine normalization requires retaining the basepoint contribution
separately.  This proposition does not assert that actual theta
coefficients range over every vector in a chosen relaxation.
Moreover, pointwise descent does not imply local boundedness through
degeneration: for
\[
 \mathsf J(\tau)=\operatorname{diag}(1,\tau^2),\qquad
 \ell(\tau)=(0,\tau),
\]
the kernel condition holds even at $\tau=0$, but
$C_{\min}=|\tau|^{-1}$ for $\tau\ne0$.
This is an analytic matrix example, not a theta heat field.
It makes a vanishing-order estimate necessary after the rank test.


\begin{proposition}[The one-sided observation quotient]
\label{obs:signed-quotient}
In the same homogeneous setting, fix $e\in\R^2$ and put
$W(c)=e\cdot\mathsf Jc$, $E(c)=\|\mathsf Jc\|^2$.
If $W\equiv0$, the estimate $W\ell\le CE$ holds with $C=0$.
Otherwise a finite nonnegative $C$ exists exactly when
$\ker\mathsf J\subseteq\ker\ell$.
Under this inclusion, let $L=\operatorname{im}\mathsf J$,
$a=P_Le$, and $b=\ell\mathsf J^\dagger$, identified with a
vector in $L$. The least nonnegative constant is
\begin{equation}
 C_{\rm sign}=\frac{a\cdot b+\|a\|\|b\|}{2}.
 \label{obs:signed-quotient-constant}
\end{equation}
\end{proposition}
\begin{proof}
Choose $c_*$ with $W(c_*)\ne0$. For $v\in\ker\mathsf J$,
the values $W,E$ are constant on $c_*+\tau v$; if $\ell(v)\ne0$,
$W\ell$ is unbounded above for one sign of $\tau$.
After descent, the inequality is
$(a\cdot z)(b\cdot z)\le C\|z\|^2$ on $L$.
The positive part of the largest eigenvalue of
$(ab^T+ba^T)/2$ is \eqref{obs:signed-quotient-constant}:
in their span the two eigenvalues are
$(a\cdot b\pm\|a\|\|b\|)/2$, with the collinear and zero
cases giving the same nonnegative formula.
\end{proof}

For the two stationary branches, reconstruction by $AH+BH_x$ has
the valid signed costs
\begin{equation}
 q_0=A_+,\qquad q_1=\frac{B+\sqrt{A^2+B^2}}2.
 \label{obs:physical-signed-costs}
\end{equation}
The second is optimal over the full observed plane; it need not be
optimal on the smaller actual theta image. For $M\ge0$,
\begin{equation}
 q_1\le M\quad\Longleftrightarrow\quad
 B\le M,\qquad A^2\le4M(M-B),
 \label{obs:signed-cost-parabola}
\end{equation}
as follows from the principal minors of
$\left(\begin{smallmatrix}M&-A/2\\-A/2&M-B\end{smallmatrix}\right)$.
Thus $B\to-\infty$ can be harmless if $A^2=O(-B)$.
This distinction survives analytic degeneration: for
\[
 \mathsf J(\tau)=\operatorname{diag}(\tau^2,\tau^4),\qquad
 \ell(\tau)=(\tau,-\tau^2),\qquad e=(0,1),
\]
endpoint descent holds at every $\tau$, and for $\tau\ne0$
the descended row is $(\tau^{-1},-\tau^{-2})$.
Its absolute norm diverges, but its signed constant is
$1/[2(\sqrt{1+\tau^2}+1)]\le1/4$; at zero the optimum is zero.
This is an analytic matrix example, not a theta heat field.
Affine normalization still retains a possible endpoint offset
$b\cdot z+d$, whose term $d(e\cdot z)$ is not covered by the
homogeneous quotient formula.

\subsection{Signed reconstruction and radial order}

For the actual field, put $F=(H,H_x)$, $\rho=\|F\|=\sqrt E$,
$W=h_j$, and $S=h_{j+1}$, with $j=0,1$.
If an endpoint observation $V$ has a reconstruction
\begin{equation}
 V=AH+BH_x+CS+R,\qquad |R|\le\varepsilon,
 \label{obs:reconstruction}
\end{equation}
then full stationarity with $E>0$ gives the necessary estimate
\begin{equation}
 \|(A,B)\|\ge\frac{(|V|-\varepsilon)_+}{\sqrt E}.
 \label{obs:reconstruction-cost}
\end{equation}
At a multiplicity $2r+j$ predecessor, $\sqrt E\asymp s^r$.
A sufficiently close fixed endpoint, or a sufficiently accurate fixed
cutoff of that endpoint, has a nonzero limit $v$.
If $\varepsilon<|v|/2$, the coefficients in
\eqref{obs:reconstruction} must therefore grow at least as $s^{-r}$.
An independently bounded reconstruction would exclude the event;
division or pseudoinversion alone does not prove its boundedness.
Bounded nonsingular changes of first-jet frame preserve this
obstruction.  A singular frame moves it into the comparison with
the physical norm.

The actual chord target is less restrictive than absolute endpoint
reconstruction.  For $\rho>0$ put $\theta=F/\rho$ and let $e_j$
select the $H$ or $H_x$ coordinate.  With $\alpha=e^{-\delta}$,
\begin{equation}
 \frac{W^2-\alpha WV}{E}
 =(e_j\cdot\theta)^2
   -\alpha\frac{(e_j\cdot\theta)V}{\rho}.
 \label{obs:radial-signed-observation}
\end{equation}
For fixed $\delta>0$, a lower bound by $-KE$ with locally bounded $K\ge0$ is
equivalent to local boundedness of $[WV]_+/E$.
A large endpoint opposite in sign to $W$ is harmless.
The manuscript's marked real Gauss map retains orientation; here
one must additionally retain the radial vanishing order relative
to $\sqrt E$.

\begin{proposition}[A sufficient signed boundary-trace estimate]
\label{obs:signed-trace}
Suppose a complete arithmetic decomposition gives
\[
 V=AH+BH_x+CS+B_{\partial}+R,\qquad |R|\le\varepsilon.
\]
On full stationary starts with $\rho>0$, suppose
\begin{equation}
 \frac{W}{\rho}B_{\partial}
 \le C_0\rho+\varepsilon_0,
 \label{obs:signed-boundary-target}
\end{equation}
where $C_0,\varepsilon,\varepsilon_0\ge0$.
Then
\begin{equation}
 WV\le(|A|+|B|+C_0)E
       +(\varepsilon+\varepsilon_0)\sqrt E.
 \label{obs:signed-boundary-transfer}
\end{equation}
\end{proposition}
\begin{proof}
Multiply \eqref{obs:signed-boundary-target} by $\rho$, use $S=0$,
and bound $|WH|$, $|WH_x|$ by $E$ and $|WR|$ by
$\varepsilon\sqrt E$.
\end{proof}

Using \eqref{obs:physical-signed-costs} improves
\eqref{obs:signed-boundary-transfer} to
\begin{equation}
 WV\le(q_j+C_0)E+(\varepsilon+\varepsilon_0)\sqrt E.
 \label{obs:sharpened-signed-transfer}
\end{equation}
It therefore suffices that $q_j,C_0$ are locally bounded and the
errors meet the prescribed fixed-length or shrinking-length requirements.
Individual $A,B$ need not be bounded. Uniform source and cutoff remainder
bounds must nevertheless be established independently for the chosen
coefficients; their divergence can invalidate an absolute source majorant.
The complete trace may take the recurrence form
$B_{\partial}=J_0-J_1+\eta/2$, with additional cutoff or source
boundaries retained explicitly.
Neither the cokernel description nor a change of geometry proves
\eqref{obs:signed-boundary-target}.
The new geometric task is an independent oriented estimate for this
arithmetic trace with the required radial order; the nonnegative
metric and its curvature magnitude do not supply that estimate.


\subsection{Scalar signed distance on prescribed cones}
\label{obs:scalar-cones}
The actual stationary geometry admits a further simplification.
On the unique curve of Theorem~\ref{phy:complete-stationary-locus},
let $\sigma=\operatorname{sgn}W$ and retain the positive factor $Q$
of \eqref{phy:scalar-energy-factor}. The exact and mixed trace
conditions become, respectively,
\begin{align}
 \frac W{\sqrt E}B_\partial\le C\sqrt E
 &\quad\Longleftrightarrow\quad \frac{B_\partial}{W}\le CQ(s),
 \label{obs:scalar-exact-trace}\\
 \frac W{\sqrt E}B_\partial\le C\sqrt E+\varepsilon_0
 &\quad\Longleftrightarrow\quad
 \sigma B_\partial\le C|W|Q(s)+\varepsilon_0\sqrt{Q(s)}.
 \label{obs:scalar-mixed-trace}
\end{align}
Because $Q(0)=1$, the zero-error condition is equivalent to boundedness
of $(B_\partial/W)_+$. For an analytic trace with first nonzero term $b_qs^q$, $b_q\ne0$,
it holds precisely when
$\sigma b_q<0$ or $q\ge r$; the identically zero trace also passes.
This finite-order criterion now covers the entire local backward
stationary locus, not merely one selected curve in a larger locus.
It supplies no uniform estimate in additional cutoff or endpoint
parameters. A fixed nonzero $\varepsilon_0$ in
\eqref{obs:scalar-mixed-trace} must not be discarded: after division
by $|W|$, it has the singular size $\varepsilon_0/|W|$.

The actual reconstruction coefficient is also scalar:
\begin{equation}
 AH+BH_x=\kappa_jW,\qquad
 \kappa_0=A,\quad\kappa_1=B+s\lambda(s)A,
 \qquad q_{\rm actual}=\frac{(\kappa_j)_+}{Q(s)}.
 \label{obs:actual-scalar-cost}
\end{equation}
For example $A=1/s$, $B=0$ has bounded actual-curve cost on an odd
predecessor while the full-plane cost is $1/(2s)$.
This is an illustrative coefficient choice conditional on the curve,
not a theta reconstruction rule. Using the unknown quotient $V/W$
to prescribe coefficients would remain circular.

An event-independent domain can be fixed in advance. Choose any
$\kappa>0$, write $Z_0=H_x$, $Z_1=H$, and put
\begin{equation}
 \mathcal D_{j,\kappa}
 =\{h_{j+1}=0,\quad W=h_j\ne0,\quad |Z_j|\le\kappa|W|\}.
 \label{obs:prescribed-cone}
\end{equation}
Only the starting jets are restricted; endpoint intervals and their
complete fields retain their original requirements.
On these domains
\begin{equation}
 W^2\le E\le(1+\kappa^2)W^2.
 \label{obs:cone-energy-comparison}
\end{equation}
Every sufficiently close matching predecessor lies in each such fixed
positive cone: this is exact for $j=0$ and follows from
$H/H_x=s\lambda(s)$ for $j=1$.
Choosing $\kappa=0$ would generally discard asymmetric odd events.

\begin{proposition}[Scalar cone transfer]
\label{obs:scalar-cone-transfer}
On $\mathcal D_{j,\kappa}$ let $\sigma=\operatorname{sgn}W$ and
$\mathcal C_\sigma=\{b\in\R:\sigma b\le0\}$.
For $C_0,\varepsilon_0\ge0$, the scalar distance bound
\begin{equation}
 \operatorname{dist}(B_\partial,\mathcal C_\sigma)
 = (\sigma B_\partial)_+\le C_0|W|+\varepsilon_0
 \label{obs:scalar-distance-target}
\end{equation}
is equivalent to
\begin{equation}
 WB_\partial\le C_0W^2+\varepsilon_0|W|.
 \label{obs:scalar-product-target}
\end{equation}
It implies \eqref{obs:signed-boundary-target} with the same constants;
conversely that physical target implies
\eqref{obs:scalar-distance-target} with constants
$C_0(1+\kappa^2)$ and $\varepsilon_0\sqrt{1+\kappa^2}$.
For the complete reconstruction in Proposition~\ref{obs:signed-trace},
the scalar bound yields
\begin{equation}
 WV\le(q_{j,\kappa}+C_0)W^2
              +(\varepsilon+\varepsilon_0)|W|,
 \qquad q_{0,\kappa}=A_+,\quad
 q_{1,\kappa}=(B+\kappa|A|)_+.
 \label{obs:cone-reconstruction-transfer}
\end{equation}
The displayed $q_{1,\kappa}$ is optimal relative to $W^2$ over the
unrestricted algebraic cone of jet pairs; no sharpness on the actual
theta image is asserted.
\end{proposition}
\begin{proof}
Multiply \eqref{obs:scalar-distance-target} by $|W|$, using
$\sigma|W|=W$. Conversely its right side is nonnegative, so
division by $|W|$ recovers the positive-part bound.
The physical comparisons follow from
\eqref{obs:cone-energy-comparison}.
On branch zero $AH+BH_x=AW$. On branch one, putting $z=H/H_x$
gives $W(AH+BH_x)=(B+Az)W^2$, where $|z|\le\kappa$.
Its maximum coefficient is $B+\kappa|A|$, attainable in the
algebraic cone, and $|WR|\le\varepsilon|W|$ gives the transfer.
\end{proof}

For instance, $A=\tau^{-2}$, $B=-\kappa\tau^{-2}$ has cone cost
zero but full-plane cost
$(\sqrt{1+\kappa^2}-\kappa)/(2\tau^2)\to\infty$.
This enlarges the sufficient coefficient class, without asserting an
admissible theta identity or a uniform remainder estimate for it.
A division-free source certificate for the scalar trace would be
\begin{equation}
 B_\partial=\mu W-\sigma P+R_0,\qquad
 \mu\le C_0,\quad P\ge0,\quad |R_0|\le\varepsilon_0.
 \label{obs:scalar-source-certificate}
\end{equation}
Multiplication by $W$ proves \eqref{obs:scalar-product-target}.
There is no required lower bound on $\mu$ or upper bound on $P$;
the independent work is to construct and bound them from the source.

\begin{corollary}[Fixed-cone mixed-error exclusion]
\label{obs:cone-mixed-exclusion}
Theorem~\ref{obs:average-mixed-exclusion} remains valid when its
starting-point estimates are assumed only on the prescribed domains
$\mathcal D_{j,\kappa}$ for any fixed $\kappa>0$ and both branches.
In particular a scalar endpoint estimate
\begin{equation}
 WV_N\le M_{j,k,\ell}W^2+\eta_{j,k,\ell}|W|
 \label{obs:cone-endpoint-criterion}
\end{equation}
suffices when $M_{j,k,\ell}\ge0$ is locally bounded for each fixed
pair of indices and $\eta_{j,k,\ell}\ge0$ tends to zero uniformly over the
approaching starts for each fixed $j,k$, with the same endpoint
cutoffs, interval coverage, and independent prescription as that theorem.
\end{corollary}
\begin{proof}
Every sufficiently close matching predecessor lies in the fixed cone.
For $\alpha=e^{-\delta_k}$, \eqref{obs:cone-endpoint-criterion} gives
$W^2-\alpha WV_N\ge-\alpha M_{j,k,\ell}E
-\alpha\eta_{j,k,\ell}\sqrt E$.
The earlier proof fixes the length, then an accurate cutoff, and then
approaches the event on those same predecessors. Nothing else changes.
\end{proof}

Scalar geometry alone cannot prove the hypothesis. The non-theta heat
polynomials $H=y^2+2s$ and $H=y^3+6sy$ have their respective
stationary centers at $y=0$, but both satisfy
\[
 \frac{W(T-s,\delta)}{W(T-s,0)}
 =1+\frac{\delta^2}{2s}\longrightarrow+\infty
 \qquad(\delta>0\text{ fixed}).
\]
They retain the proved scalar geometry and violate the bounded
positive endpoint ratio.

% Prepared for Edward Baker, 15 October 2026, with substantial LLM assistance.
% Coordinating configuration recorded by the parent: gpt-6-astra, effort ultra.
\section[Multiplicative theta identities and nonlocal boundary obstructions]{Multiplicative theta identities\\and nonlocal boundary obstructions}
\label{dil:section}

The finite-range obstruction of Section~\ref{adj:joint-source} does not
exclude arithmetic identities with nonlocal index dependence. We now give
one such family explicitly: multiplication of a square index is translation
of the source coordinate. Its prime sieve and M\"obius inversion retain
the complete heat transform, including the lower intervals created by
translation. Three other nonlocal constructions provide useful comparisons.
The identities below have uniform absolute remainder bounds, but none
currently estimates the positive part of the invariant trace at the radial
order required in Section~\ref{obs:scalar-cones}. Detailed derivations and
control comparisons are in \cite{ThetaDilation,ThetaNonlocal}; the underlying
divisor identities are classical \cite[\S27.5]{DLMF}.

\subsection{Multiplicative dilation and a uniform source remainder}
\label{dil:source-arithmetic}

Retain $p_n=e^{u-\pi n^2e^{4u}}$ and
$\phi_n=p_n(2\pi^2n^4e^{8u}-3\pi n^2e^{4u})$, so that
$\psi=\sum p_n$, $\Phi=\sum\phi_n$, and $8\phi_n=p_n''-p_n$.
For integers $d\ge1$, define
\begin{equation}
 a_d=\tfrac12\log d,\qquad
 (\mathsf D_df)(u)=d^{-1/2}f(u+a_d).
 \label{dil:operator}
\end{equation}
\begin{theorem}[Complete multiplicative source identities]
\label{dil:sieve}
At every real source coordinate,
\begin{equation}
 \mathsf D_dp_n=p_{dn},\qquad
 \mathsf D_d\phi_n=\phi_{dn},\qquad
 \mathsf D_d\mathsf D_e=\mathsf D_{de},\qquad
 \mathsf D_d\Phi=\sum_{d\mid n}\phi_n.
 \label{dil:mode-identities}
\end{equation}
For a finite prime set $P$, put $Q_P=\prod_{p\in P}p$. Then
\begin{equation}
 \begin{aligned}
 \Phi_P:=\prod_{p\in P}(I-\mathsf D_p)\Phi
 &=\sum_{d\mid Q_P}\mu(d)\mathsf D_d\Phi
 =\sum_{(n,Q_P)=1}\phi_n,\\
 \sum_{d\ge1}\mu(d)\mathsf D_d\Phi&=\phi_1.
 \end{aligned}
 \label{dil:sieve-mobius}
\end{equation}
The infinite identity converges absolutely at each fixed real $u$.
On $u\ge0$, $\Phi_P>0$. If $m$ is the smallest prime outside $P$, then
\begin{equation}
 0\le\Phi_P-\phi_1\le\sum_{n\ge m}\phi_n.
 \label{dil:sieve-tail}
\end{equation}
\end{theorem}
\begin{proof}
Substitution of $u+a_d$ multiplies $e^{4u}$ by $d^2$ and $e^u$ by
$d^{1/2}$; the prefactor in \eqref{dil:operator} cancels the latter.
This proves \eqref{dil:mode-identities}. Expanding the finite product
gives inclusion--exclusion. For the infinite identity, absolute values
of the double sum are bounded by $\sum_n\tau(n)|\phi_n(u)|$, where
$\tau(n)$ is the divisor count. Gaussian decay in $n$ makes this finite
at each fixed $u$. Collecting divisors uses
$\sum_{d\mid n}\mu(d)=\mathbf1_{n=1}$.
Every mode is positive on $u\ge0$, since $2\pi n^2e^{4u}>3$.
The least integer greater than one coprime to $Q_P$ is $m$: any smaller
integer has a prime factor in $P$, while $m$ itself is coprime.
These facts prove \eqref{dil:sieve-tail}.
\end{proof}

\begin{lemma}[Uniform moment tail]
\label{dil:moment-tail}
Fix $T,Y\ge0$ and an integer $j\ge0$. For $m\ge2$, assume
$\pi m^2\ge T/8$ and $b_m=4\pi m^2-9-Y>0$, and put
$r_m=(1+1/m)^4e^{-\pi(2m+1)}<1$. For every real $t\le T$,
\begin{equation}
 \sum_{n\ge m}\int_0^\infty u^je^{tu^2+Yu}\phi_n(u)\,du
 \le\varepsilon_j(m;T,Y)
 :=\frac{2\pi^2m^4j!e^{-\pi m^2}}
 {b_m^{j+1}(1-r_m)}.
 \label{dil:epsilon}
\end{equation}
\end{lemma}
\begin{proof}
Drop the negative summand in $\phi_n$ and use
$e^{4u}\ge1+4u+8u^2$. For $n\ge m$ the resulting integrand is bounded by
$2\pi^2n^4u^je^{-\pi n^2}e^{-b_mu}$: its omitted quadratic
exponent is $(T-8\pi n^2)u^2\le0$.
The ratio of successive $n^4e^{-\pi n^2}$ terms is at most $r_m$.
Sum the geometric majorant and integrate
$u^je^{-b_mu}$ to obtain \eqref{dil:epsilon}.
\end{proof}

The estimate is uniform in the real spectral frequency on
$|\operatorname{Im}z|\le Y$. Its displayed value depends on $T$ through
the admissibility condition. Every $m\ge2$ is admissible when
$0\le t\le1/2$ and $Y=0$.
For M\"obius truncation at $d=K\ge1$, the error is supported on $n>K$
with coefficient magnitude at most $\tau(n)\le n$.
The same proof, with $m=K+1$, replaces $m^4$ by $m^5$ and $r_m$ by
$(1+1/m)^5e^{-\pi(2m+1)}$. After discarding finitely many initial
indices if necessary, this proves locally uniform convergence of the
paired heat transforms and every fixed spectral derivative.

The sieve approximates the \emph{sifted} source by $\phi_1$; it does not
approximate the complete theta source by its first mode with an error
vanishing as $P$ grows. A complete-theta endpoint approximation can instead
retain the indices all of whose prime factors belong to $P$. Its omitted
indices also begin at $m$, so \eqref{dil:epsilon} applies, but for nonempty $P$ this retained
head is infinite. No advantage over ordinary index cutoff is asserted.

\subsection{Complete heat packets and primitive lower edges}
\label{dil:packets-edges}

Write
\[
 \mathfrak F(t,z)=\int_0^\infty e^{tu^2+izu}\Phi(u)\,du,
 \qquad \mathfrak F_n(t,z)=\int_0^\infty e^{tu^2+izu}\phi_n(u)\,du.
\]
These are entire in $z$; for real $t,x$, $H(t,x)=\operatorname{Re}
\mathfrak F(t,x)$. The complete transform of a divisible-index packet is
\begin{align}
 L_d(t,z)&=\int_0^{a_d}e^{tv^2+i(z+2ita_d)v}\Phi(v)\,dv,\notag\\
 \mathcal P_d(t,z)
 &=d^{-1/2}e^{ta_d^2-iza_d}
       [\mathfrak F(t,z+2ita_d)-L_d(t,z)]
 =\sum_{d\mid n}\mathfrak F_n(t,z).
 \label{dil:complete-packet}
\end{align}
Indeed $v=u+a_d$ in the transform of $\mathsf D_d\Phi$ gives the
displayed prefactor, shifted frequency and lower interval together.
Consequently
\begin{equation}
 \begin{aligned}
 \sum_{d\mid Q_P}\mu(d)\mathcal P_d=\mathfrak F_1+R_P,
 &\qquad |\partial_z^jR_P|\le\varepsilon_j(m;T,Y),\\
 \sum_{d\ge1}\mu(d)\mathcal P_d=\mathfrak F_1.&
 \end{aligned}
 \label{dil:packet-identities}
\end{equation}
The last sum converges absolutely with its stated derivatives.
Differentiation in $z$ acts on the prefactor and on $L_d$ as well as on
the shifted transform; it retains $(v-a_d)^j$, not $v^j$.

The unrestricted pieces
$\mathcal T_d^tf(z)=d^{-1/2}e^{ta_d^2-iza_d}f(z+2ita_d)$
obey $\mathcal T_d^t\mathcal T_e^t=\mathcal T_{de}^t$.
Their lower-edge correction is equally exact. For additive shifts $a\ge0$,
set $\mathsf D_af(u)=e^{-a}f(u+a)$,
$\mathcal T_a^tf(z)=e^{-a+ta^2-iza}f(z+2ita)$, and
\[
 C_a[f](t,z)=e^{-a+ta^2-iza}
       \int_0^a e^{tv^2+i(z+2ita)v}f(v)\,dv.
\]
Splitting the interval $[0,a+b]$ at $b$ gives
\begin{equation}
 C_{a+b}[f]=\mathcal T_a^t C_b[f]+C_a[\mathsf D_bf].
 \label{dil:lower-cocycle}
\end{equation}
Thus the transform semigroup alone is not the half-line source action.

\begin{proposition}[Complete primitive lower edge]
\label{dil:primitive-boundary}
For a polynomial-growth mask $c_n$, put
$\psi_c=\sum c_np_n$ and $b_c=\sum c_n\beta_n$, with
$\beta_n=(4\pi n^2-1)e^{-\pi n^2}$.
If $w\in C^2([0,\infty))$, the integrals below converge, and
$\psi_c'w-\psi_cw'\to0$ at infinity, then
\begin{equation}
 \int_0^\infty\Big(\sum_n c_n\phi_n\Big)w\,du
 =\frac{b_cw(0)+\psi_c(0)w'(0)}8
   +\frac18\int_0^\infty\psi_c(w''-w)\,du.
 \label{dil:masked-reconstruction}
\end{equation}
For an endpoint mask $c_n$ in the reconstruction of
Section~\ref{adj:joint-source}, the complete starting observations give
\begin{equation}
 B_{\rm src}=\frac{\mathbf1_{j=0}b_c-A/2}{8}.
 \label{dil:masked-source-edge}
\end{equation}
\end{proposition}
\begin{proof}
Integrate $\psi_c''w-\psi_cw$ twice by parts, using
$\psi_c'(0)=-b_c$. The negative lower bracket is
$b_cw(0)+\psi_c(0)w'(0)$, which proves
\eqref{dil:masked-reconstruction}. The physical weights are even, hence
$w'(0)=0$; the endpoint contributes $\mathbf1_{j=0}b_c$ and the complete
starting $H$ contributes $-A\sum\beta_n=-A/2$.
\end{proof}

For a divisible-index packet,
$b_{[d]}=\sum_k\beta_{dk}=-d^{-1/2}\psi'(a_d)$.
For the sieve, $b_P=\sum_{(n,Q_P)=1}\beta_n$, and absolute M\"obius
inversion gives $\sum_d\mu(d)b_{[d]}=\beta_1$.
Neither a sieve nor a prime-smooth endpoint has lower coefficient $1/2$
without its omitted modes. The prescribed choice
$A=2\mathbf1_{j=0}b_c$ cancels \eqref{dil:masked-source-edge}, but does
not cancel the remaining trace. For an endpoint approximation with
error \eqref{dil:epsilon}, the complete stationary start obeys
\begin{equation}
 |W\{h_j(t,x+\delta)-h_{j,c}(t,x+\delta)\}|
 \le\varepsilon_j|W|\le\varepsilon_j\sqrt E.
 \label{dil:complete-start-error}
\end{equation}
This is an admissible mixed error. Replacing the starting $W$ by a sifted
or truncated field would change this conclusion.

\subsection{Divergence of separated positive-time inversion}
\label{dil:divergence}

\begin{proposition}[The lower intervals cannot be separated]
\label{dil:separated-divergence}
For fixed real $t>0$ and complex $z$, the two M\"obius series obtained
by separating the complete-transform and lower-interval terms of
\eqref{dil:complete-packet} both fail the term test. Their paired
difference is the convergent identity \eqref{dil:packet-identities}.
\end{proposition}
\begin{proof}
With $a\to\infty$, write $\mathfrak F(t,z+2ita)$ as the Laplace
transform of $e^{tv^2}\Phi(v)$ at $2ta-iz$. This smooth function has
value $\Phi(0)>0$ at zero and decays faster than every exponential.
Taylor expansion at zero and a split at any fixed positive endpoint give
\begin{equation}
 \mathfrak F(t,z+2ita)=\frac{\Phi(0)}{2ta-iz}+O(a^{-2}).
 \label{dil:laplace-endpoint}
\end{equation}
The same leading term holds for $L_d$, since their difference is a
complete packet controlled by Lemma~\ref{dil:moment-tail}. Therefore
\begin{equation}
 \left|d^{-1/2}e^{ta_d^2-iza_d}\mathfrak F(t,z+2ita_d)\right|
 \sim\frac{\Phi(0)}{t\log d}
 e^{(t/4)(\log d)^2}d^{-1/2+\operatorname{Im}z/2}
 \longrightarrow\infty.
 \label{dil:separated-growth}
\end{equation}
Infinitely many indices are primes, with $\mu(d)=-1$. Thus neither
separated series has terms tending to zero.
\end{proof}

At $t=0$, separation would instead introduce
$\sum_d\mu(d)d^{-s}$ with $s=1/2+iz/2$.
For real $z$ this is the critical line; the ordinary absolutely convergent
Euler-product argument requires $\operatorname{Re}s>1$.
The paired identity invokes no critical-line convergence assertion.
Proposition~\ref{dil:separated-divergence} concerns the specified separation,
not all nonpolynomial multipliers or all coupled estimates.

Even reflection does not remove this cost. For nonempty finite $P$,
reflection changes the forward source shifts to backward shifts.
The largest divisor $Q_P$ dominates as $u\to-\infty$, giving
\begin{equation}
 \Phi_P(u)\sim\mu(Q_P)2\pi^2Q_P^{-5}e^{-9u}
                    \exp[-\pi e^{-4u}/Q_P^2].
 \label{dil:sieve-negative-tail}
\end{equation}
This follows by inserting the positive-tail asymptotic of the even
complete $\Phi$ in each $d^{-1/2}\Phi(u+a_d)$.
It also differs from the positive-tail asymptotic of $\Phi_P$, whose
least retained index is one. An odd number of sieving primes makes the
far negative tail negative, although $\Phi_P>0$ on $u\ge0$.

Pointwise whole-line M\"obius inversion remains valid, but its heat
integral cannot be interchanged with the series. At frequency zero,
\begin{equation}
 d^{-1/2}\int_{\R}e^{tu^2}\Phi(u+a_d)\,du
 =2d^{-1/2}e^{ta_d^2}H(t,2ita_d)
 \ge2H(t,0)d^{-1/2}e^{ta_d^2}.
 \label{dil:whole-line-growth}
\end{equation}
Substitution, evenness, and $\cosh(2ta_du)\ge1$ prove the equality and
inequality. For $t>0$ the terms again fail to vanish along primes.
Consistently, $\phi_1(u)\sim-3\pi e^{5u}$ as $u\to-\infty$;
its ordinary whole-line positive-time heat integral does not exist.

\subsection{An infinite-range recurrence current}
\label{nonlocal:massive}

The following construction leaves the finite-range class while
retaining a complete lower trace. Use $r_n,K_n,g_n,G_n$ and
$\mathcal M(g)$ from \eqref{adj:joint-source-reduction} and
Section~\ref{adj:joint-source}. For $\kappa>0$, put
$r=e^{-\kappa}$, $\lambda_\kappa=r+r^{-1}-2$, and
\begin{equation}
 \mathcal G_\kappa(n,m)
 =\frac{r^{|n-m|}-r^{n+m}}{r^{-1}-r},\qquad n,m\ge1.
 \label{nonlocal:green}
\end{equation}
\begin{proposition}[Massive current with its full trace]
\label{nonlocal:massive-current}
Set $J=-\mathcal G_\kappa g$, $J_0=0$, and
$R=\lambda_\kappa\mathcal G_\kappa g$. Then
$r_n=\Delta J_n+\partial_uK_n+R_n$. On every finite rectangle,
\begin{equation}
 \begin{aligned}
 \sum_{n=1}^M\int_0^U r_n\,du
 ={}&\int_0^U(-J_1+J_{M+1}-J_M)\,du\\
 &+\sum_{n=1}^M[K_n(U)-K_n(0)]
  +\sum_{n=1}^M\int_0^UR_n\,du.
 \end{aligned}
 \label{nonlocal:rectangle}
\end{equation}
Under the local source bounds already stated, the upper limits yield
\begin{align}
 V_N&=AH+BH_x+CS+B_{\rm src}+B_\kappa+\mathcal R_\kappa,
   \notag\\
 B_\kappa&=\sum_{m\ge1}r^mG_m,
 \qquad \mathcal R_\kappa=\sum_{m\ge1}(1-r^m)G_m,
 \qquad |\mathcal R_\kappa|\le\kappa\mathcal M(g).
 \label{nonlocal:massive-reconstruction}
\end{align}
Here $S=h_{j+1}$ and $B_{\rm src}$ is unchanged from
\eqref{adj:joint-source-boundary}.
\end{proposition}
\begin{proof}
The positive kernel \eqref{nonlocal:green} is the Dirichlet inverse of
$\lambda_\kappa-\Delta$. Direct substitution gives the identity for
$r_n$, and telescoping gives \eqref{nonlocal:rectangle}.
Its two sum identities are
$\mathcal G_\kappa(1,m)=r^m$ and
$\lambda_\kappa\sum_n\mathcal G_\kappa(n,m)=1-r^m$.
They prove \eqref{nonlocal:massive-reconstruction}, including
$1-e^{-\kappa m}\le\kappa m$. The earlier Gaussian index and source
majorants justify the limits for each fixed $\kappa$.
\end{proof}

For a common weight $w$, $B_\kappa=\frac18\int\psi_\kappa(w''-w)$,
where $\psi_\kappa=e^u\sum_{n\ge1}r^ne^{-\pi n^2e^{4u}}$.
Writing $\Theta_y(v)=\sum_{n\in\mathbb Z}e^{-\pi n^2y}e^{inv}$,
its exact circle Poisson representation is
\begin{equation}
 \psi_\kappa(u)=\frac{e^u}{2}
 \left[\frac1{2\pi}\int_{-\pi}^{\pi}
 \frac{1-r^2}{1-2r\cos v+r^2}\Theta_{e^{4u}}(v)\,dv-1\right].
 \label{nonlocal:poisson-kernel}
\end{equation}
This positive kernel still acts on the signed physical test $w''-w$.
For polynomial-growth coefficients $a_n$ with $a_0=1$, summation by parts
instead leaves $\sum_nJ_n\Delta a_n$ in
$\sum_na_ng_n=-J_1+\sum_nJ_n\Delta a_n+\sum_na_nR_n$.
Theta removes this coefficient defect. It does not sign $B_\kappa$:
as $\kappa\downarrow0$, that trace tends to
$V_N-AH-BH_x-CS-B_{\rm src}$. The small remainder becomes
$\kappa\mathcal M(g)|W|$ after multiplication by the complete start,
but gives no bound for this limiting main term.

\subsection{Euler summation and Poisson duality: the required cancellations}
\label{nonlocal:continuous-index}

Two further arithmetic representations expose the danger of estimating
their pieces absolutely. Put $\alpha=\pi e^{4u}$ and
$f_u(v)=e^u(2\alpha^2v^4-3\alpha v^2)e^{-\alpha v^2}$.
Then $f_u(n)=\phi_n(u)$, $f_u(0)=0$, and
$f_u=-\alpha e^u\partial_v(v^3e^{-\alpha v^2})$, so
$\int_0^\infty f_u(v)\,dv=0$.
Euler summation with $B_1(\{v\})=\{v\}-1/2$ gives
\begin{equation}
 \Phi(u)=\int_0^\infty B_1(\{v\})\partial_vf_u(v)\,dv
        =-\int_1^\infty\lfloor v\rfloor\partial_vf_u(v)\,dv.
 \label{nonlocal:floor-repair}
\end{equation}
For the second equality subtract the null integral
$\int_0^\infty(v-1/2)\partial_vf_u(v)\,dv=0$.
The repair matters quantitatively. If
$F_*(z)=(2z^4-3z^2)e^{-z^2}$, scaling $z=\sqrt\alpha v$ gives, as $u\to+\infty$,
\begin{equation}
 e^{-u}\int_0^\infty|B_1(\{v\})\partial_vf_u(v)|\,dv
 \longrightarrow\frac12\int_0^\infty|F_*'(z)|\,dz>0.
 \label{nonlocal:bernoulli-loss}
\end{equation}
Dominated convergence follows from $|B_1|\le1/2$.
Thus the absolute Bernoulli majorant is not heat-integrable even at
$t=0$. In contrast, for $u\ge0$, $v\ge1$,
\[
 \partial_vf_u(v)
 =-2\alpha v(2\alpha v^2-1)(\alpha v^2-3)e^{u-\alpha v^2}<0,
\]
so the repaired floor integral is a convergent positive mixture.
Its primitive counterpart has lower trace
$-\int_1^\infty\lfloor v\rfloor\beta'(v)\,dv
=\sum_n\beta_n=1/2$, where
$\beta(v)=(4\pi v^2-1)e^{-\pi v^2}$.
An endpoint cutoff uses $\min(\lfloor v\rfloor,N)$ only at that endpoint.
Nonnegative square weights of polynomial growth replace $\lfloor v\rfloor$ by
$\sum_{n\le v}a_n$. Consequently a coarse positive envelope admits the
weighted-square controls: any new signed gain must use the exact counting
function rather than positivity alone.

For the second representation, extend the index of $p_n,\phi_n$ to
$a\in\R$ and use the Fourier convention
$\widehat f(\xi)=\int_\R f(a)e^{-2\pi ia\xi}\,da$.
Gaussian transformation and $8\phi_a=(\partial_u^2-1)p_a$ give
\begin{equation}
 \widehat{p_\cdot(u)}(\xi)=p_\xi(-u),\qquad
 \widehat{\phi_\cdot(u)}(\xi)=\phi_\xi(-u),\qquad
 \Phi(u)=\sum_{k\ge1}\phi_k(-u).
 \label{nonlocal:index-poisson}
\end{equation}
The last equality is Poisson summation with $\phi_0=0$; it is the
complete theta reflection identity
\cite[\S1.8(iv), \S20.7(viii)]{DLMF}.
For finite $U$, the index function
$\int_0^U e^{tu^2}u^j\cos(yu+j\pi/2)\phi_a(u)\,du$ is Schwartz,
so the finite-source transformed identity is valid. Termwise removal
of $U$ is not. For fixed $k\ne0$,
\[
 \phi_k(-u)=-3\pi k^2e^{-5u}+5\pi^2k^4e^{-9u}+O_k(e^{-13u}),
\]
and endpoint integration by parts, for fixed $t>0$ and real $y$, yields
\begin{equation}
 \int_0^Uu^je^{tu^2+iyu}\phi_k(-u)\,du
 =-\frac{3\pi k^2U^je^{tU^2-5U+iyU}}{2tU-5+iy}
       [1+O(U^{-2})+O_k(e^{-4U})].
 \label{nonlocal:dual-divergence}
\end{equation}
Already every $j=0$ real dual mode has no improper limit; oscillation
does not control its growing amplitude. No such assertion is needed for
the identically zero test $j=1,y=0$.
The collective cancellation is large:
\begin{equation}
 \begin{aligned}
 \phi_k(-u)<0&\quad\Longleftrightarrow\quad
 0<k<\sqrt{3/(2\pi)}e^{2u},\\
 \sum_{k\ge1}|\phi_k(-u)|&=e^u(C_{\rm abs}+o(1)),
 \end{aligned}
 \label{nonlocal:dual-cancellation}
\end{equation}
where $C_{\rm abs}=\int_0^\infty
|(2\pi^2v^4-3\pi v^2)e^{-\pi v^2}|\,dv>0$.
The last asymptotic is a Riemann sum with mesh $e^{-2u}$.
The signed sum, however, is the rapidly decaying complete $\Phi(u)$.
A dual cutoff $o(e^{2u})$ eventually retains only negative modes, and
absolute bounds erase the cancellation required for source integrability.
For finite source and index cutoffs the omitted dual sum must remain
inside its source integral; it is not a negligible tail at positive time.

\subsection{Arithmetic discrimination and the unresolved scalar sign}
\label{dil:control-discriminator}

For a polynomial-growth square source $\Phi_a=\sum a_n\phi_n$,
prime excision has the exact defect
\begin{equation}
 \Phi_a-\mathsf D_p\Phi_a
 =\sum_{p\nmid n}a_n\phi_n
  +\sum_{p\mid n}(a_n-a_{n/p})\phi_n.
 \label{dil:weighted-defect}
\end{equation}
By uniqueness of Gaussian square expansions, the second source vanishes
exactly when $a_{pn}=a_n$ for every $n$. Imposing this for every prime
and $a_1=1$ forces $a_n=1$. Finitely many primes leave the coefficients
on their prime-free parts undetermined; positivity of the excised source
is weaker still. The finite weighted-square collisions fail the complete
prime identities, while the reflected quartic control has no such
square-source expansion. Generic source positivity and primitive
endpoints do not supply these arithmetic premises.

These distinctions identify legitimate arithmetic information in the
new formulas. They do not establish its signed physical consequence.
An absolute estimate of a surviving main trace gives $C|W|$, which does
not improve to $CW^2$ through a hypothetical collision. The massive
current, repaired floor integral and complete prime packets therefore
remain representations with controlled errors. A useful next estimate
must retain their exact coupling, all lower edges, full stationarity,
and the complete starting factor, and prove a bound of the form
\eqref{obs:scalar-product-target}. None of the preceding cancellations
supplies that bound by itself.

\section{The remaining arithmetic-to-physical estimate}\label{next:section}

The most focused open target combines exact theta coefficient coupling
with the signed physical boundary observation. For both
$j=0,1$, put $W=h_j$ and retain full stationarity $h_{j+1}(t,x)=0$.
For source-prescribed endpoint approximations $V_N(t,x+\delta)$, seek
\begin{equation}
 W(t,x)^2-e^{-\delta}W(t,x)V_N(t,x+\delta)
 \ge-K_{\delta,N}E(t,x)
       -\varepsilon_{\delta,N}\sqrt{E(t,x)}.
 \label{next:chord-target}
\end{equation}
For each chosen length and cutoff, the coefficient must stay locally
bounded through a hypothetical event. The endpoint approximation error
and $\varepsilon_{\delta,N}$ must tend to zero along a prescribed
cutoff sequence. The independent-cutoff criterion fixes a sufficiently
small length, then an accurate cutoff, and finally approaches the event;
it imposes no growth restriction on $K_{\delta,N}$ across cutoffs.
By Corollary~\ref{obs:cone-mixed-exclusion}, the starting domains may
be restricted to any fixed positive cones \eqref{obs:prescribed-cone}
without losing the required predecessor coverage. The theta field has
not been shown to satisfy this bound.

The original chord bound $\mathcal C_j\ge-KE^\gamma$,
$\gamma>1/2$, and the stationary and angular criteria remain
alternative sufficient hypotheses with their own coverage requirements.
For a family available only at shrinking lengths, a sufficient uniform condition is that the error decay
faster than every fixed power of those lengths, covering arbitrary finite
multiplicity. Polynomial cutoff error permits lengths such as
$1/\log(N+2)$; it does not justify an arbitrary faster schedule.

\subsection{One complete trace and one physical scalar}
The joint source identity has now been constructed in
Section~\ref{adj:joint-source}. Its two lower fluxes combine into the
invariant trace \eqref{obs:invariant-total-trace}; eliminating the
source flux leaves the same unknown observation at the index boundary.
The finite-range polynomial obstruction includes polynomial source
divergences, so further enlargement within that declared class is not
a route to an exact residual-free reconstruction beyond its boundary
range. Identities with a controlled residual and independent signed
estimates remain available questions.

Section~\ref{dil:section} gives another explicit arithmetic formulation:
finite prime identities preserve the complex shifted transforms together
with their lower intervals. The corresponding endpoint errors have the
required mixed order, but their main signed pairing remains unestimated.
Separating the M\"obius terms, integrating Poisson dual modes individually,
or replacing the exact floor multiplier by a positive envelope loses
essential cancellation. The massive current regularizes the same trace.
These conclusions restrict those specified arguments; they do not rule
out a coupled one-sided estimate retaining the full arithmetic data.

The geometry permits the next estimate to be stated with one scalar.
Fix $\kappa>0$ before considering any event, and on both prescribed
domains \eqref{obs:prescribed-cone} seek
\begin{equation}
 \boxed{WB_\partial\le C_0W^2+\varepsilon_0|W|,}
 \qquad W=h_j,\quad h_{j+1}=0,
 \label{next:signed-trace}
\end{equation}
for the complete combined trace from an independently prescribed
reconstruction. Its coefficient contribution can be paid by
$q_{j,\kappa}$ in \eqref{obs:cone-reconstruction-transfer}; bounding
individual $A,B$ is sufficient but not necessary. Alternatively, an
estimate for the complete endpoint itself of the form
\eqref{obs:cone-endpoint-criterion} avoids reconstruction coefficients.
The source formula \eqref{obs:scalar-source-certificate} states one
possible division-free certificate for the trace.

The cone restriction retains every finite-event predecessor. Complete
starting factors, the full stationary equations, complete endpoint
intervals and uniform mixed errors remain required. A decomposition
or gauge chosen by dividing by an unknown physical observation gives
no estimate. The outstanding information is a source-derived favorable
sign or sufficient positive-side vanishing order for the invariant
trace, together with its remainder bounds. No such theta theorem has
been established.

\subsection{Earlier unsuccessful constructions}
The coupled-product Gram ansatz, finite elliptic projection, and
primitive normalization correction have the specific limitations
recorded above and in their notes. These conclusions do not exclude all
arithmetic currents, integrated cancellations or one-sided estimates.
The period-two contour identity in \cite{ThetaContour} retains its
shifted ray and bottom edge, whose leading terms cancel. Separate
absolute sector bounds lose that cancellation and have exponential
cost in spectral height. The detailed asymptotics remain in the note.
Other geometric and outer-bridge alternatives are assessed in
\cite{ThetaOverview,ThetaStrategy,ThetaResearchPaths}. None supplies
the missing signed trace or extends a theorem's physical domain.

\subsection{Validation and coverage}
A proposed coefficient rule should be tested against the asymmetric
heat-polynomial and weighted-square controls, using the complete
stationary condition. The positive recurrence hierarchy supplies an
additional quantitative test: an argument claiming Gaussian convergence
from its unanchored finite constraints conflicts with the sharp ramp
controls. Conversely, even Gaussian absolute proximity to theta does
not establish the sign of theta's observation.

A theta pilot must evaluate the complete signed pairing, full stationary
root, interval conditions, endpoints and tails with coefficients fixed
before the calculation. A finite success motivates an analytic extension; it does not prove
the required cutoff family or coverage of every positive-time
finite-height event. The declared local uniformity, event coverage and
imported global hypotheses in Remark~\ref{phy:imported-bridge} remain
indispensable for an RH conclusion. A single constant uniform as
$t\downarrow0$ or $|x|\to\infty$ is not asserted or required by the
fixed-length mixed-error criterion. The joint-current, scalar and multiplicative continuations are
\cite{ThetaJointContinuation,ThetaScalarGeometry,ThetaDilation,ThetaNonlocal}.

\section*{Acknowledgments and research status}
This research manuscript was prepared for Edward Baker with
substantial assistance from OpenAI language models.  The coordinating
session used \texttt{gpt-6-astra} with reasoning effort \texttt{ultra};
internal collaborators inherited that configuration.  LLM assistance
was used for drafting, derivations, literature checks, exact-arithmetic
programs and internal reviews.  The accompanying reviews distinguish analytic arguments, validated
interval computations, and finite algebra checks.  These internal checks
are not external mathematical refereeing or formal verification.
Edward B.~Baker~III's 2014 quaternion-form notes are acknowledged as
exploratory motivation.  The author and future reviewers retain
responsibility for assessing the claims and their applicability.

\appendix
\section{Verification scope and provenance}\label{app:verification}
The verification package distinguishes analytic arguments from finite
calculations. Seven exact checkers cover the geometry, Chern/ASD
identities, complete signed bridge and heat profiles, with 1,230
overlapping assertions. Two further exact checkers add 5,667 assertions
for the chord controls and reflection algebra. The weighted-square
certificate and contour checker add 123 validated or mixed assertions
using Arb/Acb with explicit complete tails. These totals are not counts
of distinct theorems. The reflection-rigidity and arbitrary-multiplicity
arguments are analytic proofs, not consequences of finite samples.

The chord replay uses 24 asymmetric or sign-reversed heat germs,
144 full stationary roots, and 1,152 zero-free intervals at three
scales. Early intervals need not show the eventual negative sign;
the asymptotic theorem is not replaced by universal finite-step
negativity. The weighted-square calculation certifies a non-theta
collision. The contour samples check complete identities and are not
stationary theta sign certificates.

The later coefficient-core and endpoint package records 2,136 exact
and 26 mixed or validated checks; the coupled-product package records
3,278 exact and 38 mixed or validated checks. The recurrence checker
adds 6,892 exact rational assertions, independently replayed. Its rational
surrogate weights test finite current, defect and hierarchy algebra;
they are not actual-theta numerical sign data. The physical beta bound,
sharp hierarchy rate, sequence-space range and quotient-norm statements
are analytic proofs. These packages are indexed in the
\researchnote{numerics/COEFFICIENT_CORE_AND_TWO_ENDPOINT_CHECKS_20261015.json}
{coefficient-core manifest},
\researchnote{numerics/COUPLED_PRODUCT_CONTINUATION_CHECKS_20261015.json}
{coupled-product manifest}, and
\researchnote{numerics/ADJOINT_RECURRENCE_CONTINUATION_CHECKS_20261015.json}
{recurrence research manifest}.

The joint-source and signed-quotient checker adds 3,803 exact rational
assertions; the stationary scalar and gauge checker adds 1,291.
Their receipts are
\researchnote{numerics/JOINT_SOURCE_CURRENT_CHECKS_20261015.json}
{joint-current checks} and
\researchnote{numerics/STATIONARY_SCALAR_GEOMETRY_CHECKS_20261015.json}
{stationary scalar checks}, with the independent scalar replay indexed in
\researchnote{numerics/SCALAR_STATIONARY_GEOMETRY_REVIEW_CHECKS_20261015.json}
{its review manifest}. These checks cover finite identities, leading heat
coefficients, boundary signs, and quotient or cone algebra. The complete
local stationary classification, infinite-current limits and scoped
polynomial no-inverse statements rely on the analytic proofs; no finite
test proves a theta signed estimate. The
\researchnote{reviews/10_JOINT_SOURCE_CURRENT_AND_SIGNED_QUOTIENT_REVIEW_20261015.md}
{joint-current review} and
\researchnote{reviews/11_SCALAR_STATIONARY_GEOMETRY_REVIEW_20261015.md}
{scalar geometry review} distinguish those claims and their scopes.

The multiplicative-dilation checker adds 47 exact assertions, independently
replayed, for finite divisor and sieve masks, formal source scaling,
complex-shift prefactors, lower-interval substitution, masked primitive
edges, complete-start factorization and rational remainder constants.
Its \researchnote{numerics/ARITHMETIC_DILATION_CHECKS_20261015.json}
{arithmetic receipt} separates these checks from the analytic claims.
Universal divisor inversion, the infinite convergence and divergence
arguments, the moment inequalities and the nonlocal-current limits are
proved in Section~\ref{dil:section} or its cited detailed notes; finite
enumeration does not establish them. No new theta quadrature, favorable
signed domain, or improvement over ordinary cutoff is asserted.

The reproducible scripts, output receipts and source hashes are indexed
in the \researchnote{numerics/CONNECTION_BRIDGE_AND_INNER_TURNING_CHECKS_20261015.json}
{connection and profile manifest} and
\researchnote{numerics/ARITHMETIC_CONTOUR_CHORD_CHECKS_20261015.json}
{arithmetic, contour and chord manifest}. The
\researchnote{reviews/09_RECURRENCE_MANUSCRIPT_INTEGRATION_REVIEW_20261015.md}
{earlier recurrence integration review} records that checkpoint's proof
and layout checks. The \researchnote{reviews/12_SCALAR_MANUSCRIPT_AND_TAIL_CONE_REVIEW_20261015.md}{scalar integration review}
records that earlier checkpoint. The
\researchnote{reviews/15_ARITHMETIC_MANUSCRIPT_INTEGRATION_REVIEW_20261015.md}{arithmetic integration review}
records the present proof and layout checks, with the
\researchnote{numerics/ARITHMETIC_MANUSCRIPT_INTEGRATION_CHECKS_20261015.json}{integration manifest}
and a separate research-chat integrity record. The
\researchnote{manuscript/BUILD_RECORD.json}{build record}
identifies this source and its exported PDF. Older manifests describe
their own research checkpoints, including the then-unchanged initial
manuscript. They are not silently rewritten as current build records.

The foundational material comes from Investigation 21
\cite{BakerTheta} and its Notes 56--63. Principal geometric arguments
are restated, while the global theta-to-RH bridge retains its imported
status. The present integration includes exponential-average transfer,
quantitative recurrence stability, joint currents with complete boundaries,
the signed quotient, complete local stationary geometry, fixed-cone
scalar estimates and complete multiplicative heat identities with their
nonlocal boundary audit. It retains the unresolved theta sign explicitly. Negative findings are
summarized with their precise scopes and links to detailed notes; less
developed proposals remain in \cite{ThetaOverview,ThetaResearchPaths}. Local research links are relative
to the repository reading copy in \texttt{output/pdf}; the source
remains the stable editable manuscript, with milestones indexed in
\texttt{DRAFT\_HISTORY.md} and no new snapshot folders.

\begin{thebibliography}{99}
\bibitem{BakerTheta}
E. Baker (research project),
\emph{Theta correlations, Laguerre signs, and hypercomplex jets of the theta heat flow},
Investigation 21, stable research manuscript, 15 October 2026;
drafted for Edward Baker with substantial LLM assistance.
The accompanying research continuation includes Notes 56--63.

\bibitem{BakerQuaternion}
E. B. Baker III,
\emph{Quaternion-form working notes with applications to Yang--Mills theory},
unpublished exploratory drafts, 21 January and 14 February 2014.
The January version specifies $SU(2)$ in its title.
Both versions are retained in Investigation 21's source documents.

\bibitem{BakerControl}
E. Baker (research project),
\emph{Positive primitive reconstruction can coexist with a positive time collision},
Investigation 21, Note 56, 15 October 2026;
prepared for Edward Baker with substantial LLM assistance.

\bibitem{Joyce2001}
D. Joyce,
\emph{Lectures on Calabi--Yau and special Lagrangian geometry},
\href{https://arxiv.org/abs/math/0108088}{arXiv:math/0108088} (2001),
revised 2002. An expanded version appears in
M. Gross, D. Huybrechts and D. Joyce, \emph{Calabi--Yau Manifolds and
Related Geometries}, Springer, 2003.

\bibitem{Lotay2018}
J. D. Lotay,
\emph{Calibrated submanifolds},
\href{https://arxiv.org/abs/1810.08709}{arXiv:1810.08709} (2018).

% Entries checked against the original papers and their authors' institutional records.
\bibitem{AndersenBlaavand2011}
J. E. Andersen and J. L. Blaavand,
\emph{Asymptotics of Toeplitz operators and applications in TQFT},
Travaux Math\'ematiques \textbf{19} (2011), 167--201;
\href{https://arxiv.org/abs/1106.1564}{arXiv:1106.1564}.
% Author's institutional publication record:
% https://pure.au.dk/portal/en/publications/asymptotics-of-toeplitz-operators-and-applications-in-tqft/
% Original proceedings: https://math.uni.lu/geoquant/vol19.pdf

\bibitem{GrossWilson2000}
M. Gross and P. M. H. Wilson,
\emph{Large complex structure limits of K3 surfaces},
Journal of Differential Geometry \textbf{55} (2000), no.~3, 475--546;
\href{https://doi.org/10.4310/jdg/1090341262}{doi:10.4310/jdg/1090341262};
\href{https://arxiv.org/abs/math/0008018}{arXiv:math/0008018}.
% Original published PDF: https://mate.unipv.it/stoppa/didattica/GrossWilson.pdf

\bibitem{Joyce2004}
D. Joyce,
\emph{The exceptional holonomy groups and calibrated geometry},
in S. Akbulut, T. \"Onder, and R. J. Stern (eds.),
\emph{Proceedings of the G\"okova Geometry-Topology Conference 2005},
International Press, Somerville, MA, 2006, 110--139;
\href{https://arxiv.org/abs/math/0406011}{arXiv:math/0406011} (2004).
% Journal/volume details transcribed from the author's arXiv record.
% The citation key reflects the preprint year; it is distinct from Joyce2001.

\bibitem{ThetaTurning}
E. Baker (research project),
\researchnote{notes/02_signed_turning_control/01_INNER_TURNING_PROFILE_20261015.md}{\emph{Inner turning profile and physical energy on both stationary branches}},
Investigation 22, Signed turning Note 01, 15 October 2026;
prepared for Edward Baker with substantial LLM assistance.

\bibitem{ThetaChord}
E. Baker (research project),
\researchnote{notes/02_signed_turning_control/02_CENTRAL_ENERGY_CHORD_AND_CUTOFF_TRANSFER_20261015.md}{\emph{Central-energy chord and cutoff transfer}},
Investigation 22, Signed turning Note 02, 15 October 2026;
prepared for Edward Baker with substantial LLM assistance.

\bibitem{ThetaMetric}
E. Baker (research project),
\researchnote{notes/03_theta_bundles_and_connections/02_ASD_METRIC_FREEDOM_AND_PHYSICAL_RESTRICTION_20261015.md}{\emph{ASD metric freedom and physical restriction}},
Investigation 22, Theta connections Note 02, 15 October 2026;
prepared for Edward Baker with substantial LLM assistance.

\bibitem{ThetaBridge}
E. Baker (research project),
\researchnote{notes/03_theta_bundles_and_connections/03_COMPLETE_SIGNED_CORRELATION_BRIDGE_20261015.md}{\emph{Complete signed correlation bridge}},
Investigation 22, Theta connections Note 03, 15 October 2026;
prepared for Edward Baker with substantial LLM assistance.

\bibitem{ThetaContour}
E. Baker (research project),
\researchnote{notes/03_theta_bundles_and_connections/04_PERIOD_TWO_CONTOUR_PACKET_AND_SECTOR_OBSTRUCTION_20261015.md}{\emph{Period-two contour packet and sector obstruction}},
Investigation 22, Theta connections Note 04, 15 October 2026;
prepared for Edward Baker with substantial LLM assistance.

\bibitem{ThetaSquareControl}
E. Baker (research project),
\researchnote{notes/05_control_families/01_SQUARE_SUPPORT_ENDPOINT_COLLISION_OBSTRUCTION_20261015.md}{\emph{Square-support endpoint collision obstruction}},
Investigation 22, Control families Note 01, 15 October 2026;
prepared for Edward Baker with substantial LLM assistance.

\bibitem{ThetaReflection}
E. Baker (research project),
\researchnote{notes/05_control_families/02_WEIGHTED_COMB_REFLECTION_RIGIDITY_20261015.md}{\emph{Weighted-comb reflection rigidity}},
Investigation 22, Control families Note 02, 15 October 2026;
prepared for Edward Baker with substantial LLM assistance.

\bibitem{ThetaOverview}
E. Baker (research project),
\researchnote{research-overview/02_CHORD_CRITERION_AND_ARITHMETIC_BRIDGE_20261015.md}{\emph{Chord criterion, arithmetic bridge and next research steps}},
Investigation 22, Research overview 02, 15 October 2026;
prepared for Edward Baker with substantial LLM assistance.

\bibitem{ThetaStrategy}
E. Baker (research project),
\researchnote{notes/00_research_coordination/03_PRIORITIZED_SIGNED_BOUND_RESEARCH_PATHS_20261015.md}{\emph{Prioritized signed-bound research paths}},
Investigation 22, Research coordination Note 03, 15 October 2026;
prepared for Edward Baker with substantial LLM assistance.

\bibitem{ThetaFactorChord}
E. Baker (research project),
\researchnote{notes/02_signed_turning_control/03_FACTOR_PRESERVING_AND_AVERAGED_CHORD_TARGETS_20261015.md}{\emph{Factor-preserving and averaged chord targets}},
Investigation 22, Signed turning Note 03, 15 October 2026;
prepared for Edward Baker with substantial LLM assistance.

\bibitem{ThetaExponential}
E. Baker (research project),
\researchnote{notes/02_signed_turning_control/04_EXPONENTIAL_AVERAGE_AND_SHRINKING_CHORD_TARGETS_20261015.md}{\emph{Exponential average and shrinking chord targets}},
Investigation 22, Signed turning Note 04, 15 October 2026;
prepared for Edward Baker with substantial LLM assistance.

\bibitem{ThetaEllipticCore}
E. Baker (research project),
\researchnote{notes/03_theta_bundles_and_connections/05_ELLIPTIC_SHIFT_CORE_AND_PROJECTION_OBSTRUCTION_20261015.md}{\emph{Elliptic shift core and projection obstruction}},
Investigation 22, Theta connections Note 05, 15 October 2026;
prepared for Edward Baker with substantial LLM assistance.

\bibitem{ThetaCoupled}
E. Baker (research project),
\researchnote{notes/03_theta_bundles_and_connections/06_COUPLED_THETA_PRODUCT_AND_PARITY_SECTORS_20261015.md}{\emph{Coupled theta product and parity sectors}},
Investigation 22, Theta connections Note 06, 15 October 2026;
prepared for Edward Baker with substantial LLM assistance.

\bibitem{ThetaParityStrip}
E. Baker (research project),
\researchnote{notes/03_theta_bundles_and_connections/07_PARITY_STRIP_CUTOFF_AND_FULL_STARTING_FACTOR_20261015.md}{\emph{Parity strip cutoff and full starting factor}},
Investigation 22, Theta connections Note 07, 15 October 2026;
prepared for Edward Baker with substantial LLM assistance.

\bibitem{ThetaTwoEndpoint}
E. Baker (research project),
\researchnote{notes/05_control_families/03_TWO_ENDPOINT_SQUARE_COLLISION_20261015.md}{\emph{Two-endpoint square collision}},
Investigation 22, Control families Note 03, 15 October 2026;
prepared for Edward Baker with substantial LLM assistance.

\bibitem{ThetaReflectedControl}
E. Baker (research project),
\researchnote{notes/05_control_families/04_REFLECTION_POSITIVE_SOURCE_COLLISION_AND_MIXED_SECTORS_20261015.md}{\emph{Reflection-positive source collision and mixed sectors}},
Investigation 22, Control families Note 04, 15 October 2026;
prepared for Edward Baker with substantial LLM assistance.

\bibitem{ThetaAdjoint}
E. Baker (research project),
\researchnote{notes/03_theta_bundles_and_connections/08_ADJOINT_RECURRENCE_ENDPOINT_EXPERIMENT_20261015.md}{\emph{Adjoint recurrence endpoint experiment and quantitative stability}},
Investigation 22, Theta connections Note 08, 15 October 2026;
prepared for Edward Baker with substantial LLM assistance.

\bibitem{ThetaBoundaryGeometry}
E. Baker (research project),
\researchnote{notes/01_physical_jet_geometry/01_ARITHMETIC_BOUNDARY_TRACE_AND_SIGNED_OBSERVATION_20261015.md}{\emph{Arithmetic boundary traces and signed physical observation}},
Investigation 22, Physical jet geometry Note 01, 15 October 2026;
prepared for Edward Baker with substantial LLM assistance.


\bibitem{ThetaJointCurrent}
E. Baker (research project),
\researchnote{notes/03_theta_bundles_and_connections/09_JOINT_SOURCE_CURRENT_AND_POLYNOMIAL_OBSTRUCTION_20261015.md}{\emph{Joint coefficient and source currents with complete boundaries}},
Investigation 22, Theta connections Note 09, 15 October 2026;
prepared for Edward Baker with substantial LLM assistance.

\bibitem{ThetaSignedQuotient}
E. Baker (research project),
\researchnote{notes/01_physical_jet_geometry/02_SIGNED_OBSERVATION_QUOTIENT_AND_COEFFICIENT_SINGULARITIES_20261015.md}{\emph{The exact signed observation quotient and allowable coefficient singularities}},
Investigation 22, Physical jet geometry Note 02, 15 October 2026;
prepared for Edward Baker with substantial LLM assistance.

\bibitem{ThetaScalarGeometry}
E. Baker (research project),
\researchnote{notes/01_physical_jet_geometry/03_SCALAR_STATIONARY_GEOMETRY_AND_INVARIANT_TRACE_20261015.md}{\emph{Scalar boundary geometry on the complete stationary locus}},
Investigation 22, Physical jet geometry Note 03, 15 October 2026;
prepared for Edward Baker with substantial LLM assistance.

\bibitem{ThetaJointContinuation}
E. Baker (research project),
\researchnote{notes/00_research_coordination/10_JOINT_SOURCE_CURRENT_AND_SIGNED_QUOTIENT_CONTINUATION_20261015.md}{\emph{Joint source current and signed quotient continuation}},
Investigation 22, Research coordination Note 10, 15 October 2026;
prepared for Edward Baker with substantial LLM assistance.

\bibitem{ThetaResearchPaths}
E. Baker (research project),
\researchnote{notes/00_research_coordination/07_RESEARCH_PATHS_AFTER_COUPLED_PRODUCT_20261015.md}{\emph{Research paths after the coupled-product experiment}},
Investigation 22, Research coordination Note 07, 15 October 2026;
prepared for Edward Baker with substantial LLM assistance.

\bibitem{ThetaContinuation}
E. Baker (research project),
\researchnote{notes/00_research_coordination/09_MANUSCRIPT_RECURRENCE_INTEGRATION_AND_NEXT_CHAT_20261015.md}{\emph{Recurrence manuscript integration and next-chat continuation}},
Investigation 22, Research coordination Note 09, 15 October 2026;
prepared for Edward Baker with substantial LLM assistance.


\bibitem{ThetaDilation}
E. Baker (research project),
\researchnote{notes/03_theta_bundles_and_connections/11_MULTIPLICATIVE_DILATION_AND_COMPLETE_HEAT_PACKETS_20261015.md}{\emph{Multiplicative dilation and complete heat packets}},
Investigation 22, Theta connections Note 11, 15 October 2026;
prepared for Edward Baker with substantial LLM assistance.

\bibitem{ThetaNonlocal}
E. Baker (research project),
\researchnote{notes/03_theta_bundles_and_connections/12_NONLOCAL_ARITHMETIC_MULTIPLIERS_AND_BOUNDARY_OBSTRUCTIONS_20261015.md}{\emph{Nonlocal arithmetic multipliers and boundary obstructions}},
Investigation 22, Theta connections Note 12, 15 October 2026;
prepared for Edward Baker with substantial LLM assistance.

\bibitem{DLMF}
NIST Digital Library of Mathematical Functions,
\href{https://dlmf.nist.gov/27.5}{\S27.5, Inversion formulas},
\href{https://dlmf.nist.gov/1.8.iv}{\S1.8(iv), Poisson's summation formula},
and \href{https://dlmf.nist.gov/20.7.viii}{\S20.7(viii), Theta transformations};
consulted 15 October 2026.

\end{thebibliography}
\end{document}
